% Created 2026-04-24 sex 11:22
% Intended LaTeX compiler: lualatex
\documentclass[12pt,a4paper]{article}
\usepackage{fgv-dissertacao}
\autor{Gustavo M. Mendes de Tarso}
\titulo{Inteligência Artificial no Governo Público Federal}
\subtitulo{Governança, Ética, Compliance e Gestão de Riscos na Administração Pública Federal Brasileira (2019–2026)}
\cidade{Brasília}
\instituicao{FUNDAÇÃO GETÚLIO VARGAS}
\programa{Mestrado Profissional em Políticas Públicas e Governo}
\curso{Mestrado Profissional em Políticas Públicas e Governo}
\areadeconcentracao{Políticas Públicas e Governo}
\naturezatrabalho{Dissertação apresentada à Fundação Getúlio Vargas, como requisito para obtenção do título de Mestre em Políticas Públicas e Governo.}
\dataaprovacao{A definir}
\linhapesquisa{Governança, Estado e Políticas Públicas}
\author{Gustavo M. Mendes de Tarso}
\date{\today}
\title{Inteligência Artificial no Governo Público Federal: Governança, Ética, Compliance e Gestão de Riscos na Administração Pública Federal Brasileira (2019–2026)}
\usepackage[backend=biber,style=apa]{biblatex}
\addbibresource{dissertacao_ia_governo_publico_federal.bib}
\begin{document}


\capa
\folhaderosto
\fichacatalograficaaviso
\folhadeaprovacao
\section*{DEDICATÓRIA}
\label{sec:org9468969}
À minha família, pelo apoio silencioso e constante, e a todos que acreditam no aprimoramento das instituições públicas por meio do conhecimento crítico e responsável.
\section*{AGRADECIMENTOS}
\label{sec:orgf1d8c1b}
Agradeço à Fundação Getulio Vargas pelo ambiente acadêmico estimulante, aos professores do Mestrado Profissional em Políticas Públicas e Governo pelas contribuições formativas, e aos colegas de trajetória pelo diálogo qualificado. Registro, ainda, minha gratidão à minha família, pelo incentivo, compreensão e confiança ao longo deste percurso.
\section*{EPÍGRAFE}
\label{sec:org35314ff}
"A ciência é, portanto, uma perversão de si mesma, a menos que tenha como fim último melhorar a humanidade."
Nikola Tesla
\section*{RESUMO}
\label{sec:org6b70e9a}
Esta dissertação examina, por meio de revisão de literatura, o uso de sistemas e aplicações de inteligência artificial na administração pública federal brasileira como instrumento de governança, ética, compliance e gestão de riscos, no período de 2019 a 2026. O estudo parte do reconhecimento de que a incorporação de soluções baseadas em dados e algoritmos no setor público vem ampliando capacidades estatais em automação de processos, apoio à decisão, análise documental, fiscalização e prestação de serviços digitais, ao mesmo tempo em que introduz desafios institucionais relacionados à opacidade decisória, responsabilização, vieses discriminatórios, segurança da informação, proteção de dados, conformidade normativa e legitimidade administrativa. Metodologicamente, a pesquisa caracteriza-se como revisão de literatura de natureza qualitativa, com foco em produção acadêmica e técnico-institucional pertinente ao tema, em português e inglês, buscando identificar categorias analíticas recorrentes, convergências, tensões e lacunas interpretativas. Os resultados indicam que a literatura converge para a necessidade de arranjos de governança capazes de combinar estratégia institucional, supervisão humana significativa, transparência proporcional, rastreabilidade, auditoria, avaliação de impacto, controles internos e mecanismos de responsabilização ao longo do ciclo de vida dos sistemas de IA. Também se observa que a ética, quando tratada de forma isolada e principiológica, mostra-se insuficiente sem sua tradução em instrumentos operacionais de compliance e gestão de riscos. No contexto do governo público federal brasileiro, a revisão sugere que a implementação responsável da IA depende menos da mera adoção tecnológica e mais da construção de capacidades organizacionais, normativas e analíticas que assegurem aderência ao interesse público, integridade decisória e proteção de direitos. Conclui-se que a IA pode contribuir para a modernização do Estado, desde que inserida em estruturas robustas de governança pública digital, com critérios de auditabilidade, proporcionalidade, controle institucional e revisão contínua dos riscos associados ao seu uso.
Palavras-chave: inteligência artificial; administração pública federal; governança da IA; compliance público; gestão de riscos; ética algorítmica.
\section*{ABSTRACT}
\label{sec:org38d5d36}
This dissertation examines, through a literature review, the use of artificial intelligence systems and applications in the Brazilian federal public administration as an instrument of governance, ethics, compliance, and risk management from 2019 to 2026. The study is grounded in the recognition that the incorporation of data- and algorithm-based solutions in the public sector has expanded state capacities in process automation, decision support, document analysis, oversight, and digital service delivery, while simultaneously introducing institutional challenges related to decision opacity, accountability, discriminatory bias, information security, data protection, regulatory compliance, and administrative legitimacy. Methodologically, the research is a qualitative literature review focused on academic and technical-institutional production in Portuguese and English, aiming to identify recurring analytical categories, convergences, tensions, and interpretive gaps. The findings indicate that the literature converges on the need for governance arrangements capable of combining institutional strategy, meaningful human oversight, proportional transparency, traceability, auditing, impact assessment, internal controls, and accountability mechanisms throughout the AI system life cycle. The review also shows that ethics, when addressed in an isolated and merely principled manner, is insufficient unless translated into operational compliance and risk management instruments. In the context of the Brazilian federal government, the review suggests that the responsible implementation of AI depends less on technological adoption alone and more on the development of organizational, regulatory, and analytical capacities that ensure alignment with the public interest, decisional integrity, and rights protection. It is concluded that AI may contribute to state modernization, provided that it is embedded in robust digital public governance structures, with criteria for auditability, proportionality, institutional control, and continuous review of the risks associated with its use.
Keywords: artificial intelligence; federal public administration; AI governance; public compliance; risk management; algorithmic ethics.
\section*{LISTA DE ILUSTRAÇÕES}
\label{sec:org10eb41a}
\begin{itemize}
\item Figura 1 — Dimensões analíticas da governança da inteligência artificial no setor público
\item Figura 2 — Ciclo de vida da IA e pontos de controle institucional
\item Figura 3 — Relação entre ética, compliance e gestão de riscos na administração pública federal
\end{itemize}
\section*{LISTA DE TABELAS}
\label{sec:orgde79516}
\begin{itemize}
\item Tabela 1 — Síntese do corpus bibliográfico selecionado
\item Tabela 2 — Aplicações de inteligência artificial identificadas no setor público
\item Tabela 3 — Principais riscos, salvaguardas e mecanismos de mitigação mapeados na literatura
\end{itemize}
\section*{LISTA DE ABREVIATURAS E SIGLAS}
\label{sec:org3683e9e}
\begin{itemize}
\item IA — Inteligência artificial
\item APF — Administração Pública Federal
\item FGV — Fundação Getulio Vargas
\item LGPD — Lei Geral de Proteção de Dados Pessoais
\item TCU — Tribunal de Contas da União
\item CGU — Controladoria-Geral da União
\item OCDE — Organização para a Cooperação e Desenvolvimento Econômico
\item SEI — Sistema Eletrônico de Informações
\end{itemize}
\section*{LISTA DE SÍMBOLOS}
\label{sec:org1e638af}
\section{INTRODUÇÃO}
\label{sec:org612e52a}
\subsection{Contextualização}
\label{sec:org64084ee}
A incorporação da inteligência artificial (IA) no setor público deixou de ser um horizonte meramente prospectivo para se tornar um vetor efetivo de transformação administrativa, com impactos sobre rotinas decisórias, desenho de serviços, fiscalização, processamento documental e coordenação organizacional. A literatura recente converge ao demonstrar que a adoção de IA na administração pública não se resume à digitalização incremental, mas introduz uma camada sociotécnica capaz de reconfigurar a relação entre burocracia, dados, discricionariedade e accountability \autocite{madan2022ai} \autocite{criado2024artificial}. Nesse movimento, a promessa de maior eficiência, celeridade e capacidade analítica convive com preocupações substantivas relativas à opacidade algorítmica, à qualidade dos dados, à assimetria de capacidades institucionais e aos riscos de reprodução de vieses em contextos decisórios sensíveis \autocite{parycek2023artificial} \autocite{batool2025ai}.
No campo da administração pública, esse debate adquire densidade particular porque a IA opera em ambientes normativamente mais exigentes do que aqueles do setor privado. Enquanto empresas podem priorizar desempenho, redução de custos e vantagem competitiva, órgãos públicos estão vinculados a legalidade, impessoalidade, motivação, transparência, devido processo e proteção do interesse público. Por isso, a literatura tem insistido que a difusão de IA em governos deve ser lida menos como simples adoção tecnológica e mais como problema de governança pública, envolvendo atores, regras, capacidades organizacionais e escolhas políticas em níveis micro, meso e macro \autocite{criado2024artificial}. Essa perspectiva é especialmente relevante para o contexto brasileiro, em que a modernização estatal convive com heterogeneidade institucional, pressões por transformação digital e necessidade de ampliar controles sobre sistemas automatizados utilizados em atividades-fim e atividades-meio \autocite{serra2024a}.
As revisões de literatura sobre adoção e difusão da IA mostram que os usos governamentais se expandem em múltiplas frentes, como triagem e classificação de documentos, análise preditiva, atendimento digital, apoio à formulação e implementação de políticas, monitoramento regulatório e detecção de irregularidades \autocite{madan2022ai}. Contudo, a mesma literatura também assinala que a expansão das aplicações não foi acompanhada, com igual intensidade, por arranjos institucionais maduros para supervisão, auditoria, explicabilidade e responsabilização. Em outras palavras, o avanço funcional da IA frequentemente supera o avanço normativo e gerencial de seus mecanismos de controle \autocite{batool2025ai}. Essa assimetria é central para o setor público federal brasileiro, no qual a adoção de soluções algorítmicas pode afetar direitos, acesso a benefícios, priorização de ações fiscalizatórias e distribuição de capacidades administrativas.
Ao analisar a inserção da IA na administração pública, Serra e Machado destacam que a literatura recente tende a organizar o debate em torno de aplicações, resultados esperados e exigência de governança adequada, enfatizando que a implementação responsável demanda sistematização, coordenação institucional e atenção explícita às dimensões éticas \autocite{serra2024a}. Esse ponto é decisivo porque desloca a discussão de uma visão instrumental da tecnologia para uma compreensão institucional de seu uso. Em vez de perguntar apenas onde a IA pode ser aplicada, torna-se necessário perguntar sob quais condições ela deve ser empregada, com quais salvaguardas, por quais estruturas de supervisão e com quais mecanismos de correção de falhas.
A literatura internacional reforça essa inflexão. Parycek, Schmid e Novak mostram que, em procedimentos administrativos, o potencial da IA depende diretamente de condições de contorno como qualidade e disponibilidade dos dados, adequação jurídica, transparência dos modelos e preservação da possibilidade de revisão humana \autocite{parycek2023artificial}. Já Batool, Zowghi e Bano, ao sistematizarem o campo da governança de IA, evidenciam que estruturas robustas de governança precisam articular princípios, processos, papéis organizacionais, monitoramento contínuo e tratamento de incidentes, superando abordagens puramente declaratórias \autocite{batool2025ai}. Assim, governança, ética, compliance e gestão de riscos não aparecem como preocupações periféricas, mas como dimensões constitutivas da própria viabilidade administrativa e democrática da IA.
Outro aspecto relevante decorre do fato de que a adoção de IA no setor público é marcada por desigualdades de capacidade institucional e por diferentes percepções dos gestores sobre benefícios e restrições. Estudos empíricos em governos locais indicam que dirigentes públicos reconhecem ganhos potenciais em eficiência e inovação, mas identificam obstáculos importantes relacionados a recursos, competências, integração organizacional e incertezas regulatórias \autocite{yiitcanlar2022artificial}. Em chave semelhante, análises sobre gestão de pessoas no setor público mostram que a IA pode apoiar recrutamento, planejamento de força de trabalho e gestão de desempenho, mas também introduz novos riscos éticos, de exclusão e de dependência tecnológica, exigindo competências especializadas e governança adaptativa \autocite{chilunjika2022artificial}. Embora esses estudos incidam sobre níveis e contextos distintos, eles oferecem insumos valiosos para o caso federal brasileiro ao demonstrar que a maturidade do uso da IA depende menos do fascínio tecnológico e mais da capacidade estatal de regular, supervisionar e aprender institucionalmente.
Nesse cenário, a presente dissertação se insere em um debate contemporâneo e estratégico: como compreender o uso da IA no governo público federal brasileiro não apenas como inovação administrativa, mas como instrumento que precisa ser governado à luz de critérios de ética, compliance e gestão de riscos. A relevância do tema decorre do crescimento de aplicações algorítmicas em funções de apoio à decisão, automação de fluxos, análise documental, fiscalização e prestação de serviços, ao mesmo tempo em que persistem lacunas quanto aos modelos de governança, às salvaguardas institucionais e aos parâmetros para implementação responsável. Em síntese, a contextualização aqui delineada indica que o problema não reside simplesmente em adotar ou não adotar IA, mas em definir como o Estado federal brasileiro pode fazê-lo de modo legítimo, auditável, controlável e orientado ao interesse público \autocite{madan2022ai} \autocite{criado2024artificial} \autocite{serra2024a}.
\subsection{Delimitação do estudo}
\label{sec:orga79722b}
Este estudo delimita seu objeto à análise do uso de sistemas e aplicações de inteligência artificial no âmbito da administração pública federal brasileira, compreendendo órgãos e entidades da administração direta e indireta, com foco nas condições institucionais que tornam essa adoção compatível com governança pública, integridade administrativa e proteção do interesse coletivo. Não se trata, portanto, de uma investigação geral sobre transformação digital do Estado, governo eletrônico ou inovação pública em sentido amplo. O recorte concentra-se especificamente em iniciativas, debates e formulações analíticas relativas à IA aplicada à automação de processos, apoio à decisão, análise documental, fiscalização e prestação de serviços digitais, desde que esses usos suscitem questões de governança, ética, compliance e gestão de riscos \autocite{madan2022ai} \autocite{criado2024artificial} \autocite{lyrio2025emerging}.
A delimitação também é temporal. O período de 2019 a 2026 foi escolhido por corresponder à fase em que a difusão da IA no setor público se intensifica, tanto pelo amadurecimento de aplicações baseadas em aprendizado de máquina e automação inteligente quanto pela crescente preocupação com marcos de governança, auditabilidade e uso responsável. Nesse intervalo, a literatura passa a deslocar o debate do entusiasmo tecnológico para a discussão sobre condições de implementação, capacidades organizacionais e mecanismos de supervisão. Tal deslocamento é particularmente visível em revisões que enfatizam o setor público como ambiente de adoção regulado por exigências de legitimidade, accountability e controle, e não apenas por métricas de eficiência \autocite{sigfrids2022how} \autocite{batool2025ai} \autocite{lyrio2025emerging}.
No plano temático, a dissertação não pretende avaliar desempenho técnico de algoritmos específicos, comparar arquiteturas computacionais ou propor modelos matemáticos de otimização. Também não abrange, de modo exaustivo, toda a administração pública brasileira em seus níveis federal, estadual e municipal. Embora estudos sobre governos locais, experiências internacionais e casos setoriais apareçam na revisão como referências comparativas, eles serão mobilizados apenas na medida em que contribuam para compreender o caso da administração pública federal brasileira e para iluminar padrões de adoção, barreiras institucionais e salvaguardas relevantes \autocite{yiitcanlar2022artificial} \autocite{david2023understanding} \autocite{criado2024inteligencia}. O foco permanece no governo federal, por ser a esfera em que se concentram capacidades normativas, infraestruturas digitais estratégicas, funções regulatórias e arranjos de coordenação com efeitos mais amplos sobre o aparato estatal.
Além disso, o estudo se delimita como revisão de literatura de caráter analítico e integrador. Seu propósito não é mensurar empiricamente a incidência de IA em cada órgão federal, nem realizar auditoria direta de sistemas em operação. A ênfase recai sobre como a produção acadêmica e técnico-institucional descreve usos, identifica riscos e propõe instrumentos de governança. Esse recorte é metodologicamente coerente com a constatação de que a literatura recente já permite mapear aplicações recorrentes, assimetrias de maturidade e padrões de capacidade institucional, inclusive no contexto público e latino-americano \autocite{mariano2026inteligencia} \autocite{goulart2025processos} \autocite{criado2024inteligencia}.
A delimitação substantiva do estudo privilegia quatro eixos interdependentes. O primeiro é a governança da IA, entendida como conjunto de princípios, estruturas, processos decisórios e responsabilidades organizacionais necessários para orientar o uso da tecnologia no setor público \autocite{sigfrids2022how}. O segundo é a ética aplicada, especialmente no que diz respeito a transparência, justiça, não discriminação, supervisão humana e explicabilidade em contextos administrativos \autocite{brkan2020legal} \autocite{sigfrids2022how}. O terceiro é o compliance, compreendido não apenas como aderência formal a normas, mas como incorporação de controles, rastreabilidade e responsabilização ao longo do ciclo de vida dos sistemas. O quarto é a gestão de riscos, incluindo riscos operacionais, legais, reputacionais, organizacionais e sociais decorrentes do uso de IA em atividades públicas \autocite{batool2025ai} \autocite{pi2021machine}.
Nesse ponto, a literatura sobre auditoria de IA é central para delimitar o enfoque da dissertação. Mökander mostra que auditorias de IA vêm sendo estruturadas a partir de abordagens legais, éticas e técnicas, e que sua efetividade depende de clareza quanto ao objeto auditado, aos critérios de conformidade e aos efeitos institucionais esperados \autocite{mkander2023auditing}. Já Minkkinen, Laine e Mäntymäki argumentam que sistemas adaptativos e dinâmicos desafiam auditorias periódicas tradicionais, exigindo formas de monitoramento contínuo e instrumentos capazes de acompanhar mudanças de comportamento, dados e desempenho ao longo do tempo \autocite{minkkinen2022continuous}. Para o recorte desta dissertação, essas contribuições são decisivas porque reforçam que governança e controle da IA no setor público não podem ser pensados como validações pontuais anteriores à implantação, mas como práticas permanentes de supervisão, revisão e correção.
A delimitação inclui ainda a dimensão das capacidades estatais requeridas para adoção responsável. A literatura recente sobre competências em IA na gestão pública indica que a implementação de sistemas inteligentes depende não apenas de infraestrutura tecnológica, mas de repertórios organizacionais e profissionais relacionados à interpretação de dados, avaliação crítica de modelos, gestão interdisciplinar, compreensão regulatória e tomada de decisão ética \autocite{melo2026competencias}. Esse ponto restringe o estudo a uma abordagem institucional da tecnologia: mais do que perguntar quais ferramentas estão disponíveis, interessa compreender se o governo federal dispõe, ou precisa desenvolver, capacidades para governá-las de modo consistente com legalidade, integridade e accountability. Assim, o objeto da dissertação não é a IA em abstrato, mas a IA enquanto problema de governo, de coordenação administrativa e de controle público \autocite{melo2026competencias} \autocite{sigfrids2022how} \autocite{mkander2023auditing}.
\subsection{Problema de pesquisa}
\label{sec:org1287e85}
A incorporação de sistemas de inteligência artificial na administração pública tem sido frequentemente apresentada como vetor de eficiência, automação, incremento analítico e ampliação da capacidade estatal de processar grandes volumes de dados. Contudo, quando tais sistemas passam a influenciar procedimentos administrativos, triagens, priorizações, fiscalização, análise documental e apoio à decisão, o problema deixa de ser apenas tecnológico e passa a ser eminentemente institucional. A literatura recente indica que a adoção de IA no setor público não pode ser compreendida somente pela ótica dos ganhos operacionais, pois sua implementação envolve disputas sobre responsabilidade decisória, qualidade dos dados, supervisão humana, transparência, legitimidade administrativa e capacidade de controle \autocite{madan2022ai} \autocite{parycek2023artificial} \autocite{criado2024artificial}. No caso brasileiro, e especialmente no âmbito federal, esse desafio é agravado pela coexistência de diferentes graus de maturidade digital entre órgãos, pela heterogeneidade de capacidades técnico-burocráticas e pela necessidade de compatibilizar inovação com juridicidade, accountability e proteção de direitos.
É nesse ponto que emerge o problema central desta dissertação. Embora a difusão da IA no setor público federal brasileiro avance em múltiplas frentes, a literatura ainda se mostra dispersa quanto às condições concretas que permitem transformar adoção tecnológica em uso governado, ético, auditável e aderente a exigências de compliance e gestão de riscos. Em vez de um quadro consolidado, observa-se um campo marcado por assimetrias. Mariano, Ceolin e Correia-Neto identificam precisamente esse padrão ao destacar que a administração pública brasileira apresenta níveis desiguais de maturidade no emprego de IA e IA generativa, com coexistência de iniciativas avançadas e arranjos institucionais ainda incipientes para orientar seu uso responsável \autocite{mariano2026inteligencia}. De modo convergente, estudos sobre governança pública mediada por algoritmos apontam que o sucesso da adoção depende menos da mera disponibilidade da tecnologia e mais da existência de estruturas de coordenação, regras internas, competências organizacionais e mecanismos de monitoramento contínuo \autocite{goulart2025processos} \autocite{batool2025ai}.
A formulação do problema requer reconhecer, ainda, que os riscos não se distribuem uniformemente entre tipos de aplicação. Em áreas como gestão de pessoas, por exemplo, o uso de analytics e IA pode apoiar recrutamento, alocação, previsão de desempenho e planejamento da força de trabalho, mas também amplia a possibilidade de reproduzir vieses, obscurecer critérios decisórios e reduzir a inteligibilidade das escolhas administrativas. Cho, Choi e Choi mostram que, na gestão pública de pessoas, o valor analítico desses sistemas convive com cautelas importantes relacionadas à qualidade dos dados, à interpretação dos resultados e à necessidade de preservar julgamento humano e sensibilidade contextual \autocite{cho2023human}. Na mesma direção, Barros enfatiza que a adoção de IA na gestão de pessoas do setor público não pode ser separada de questões éticas, organizacionais e laborais, pois a promessa de inovação tende a tensionar valores como equidade, dignidade e responsabilidade institucional \autocite{barros2025inteligencia}. Esses achados são relevantes porque demonstram que o problema não reside apenas em “usar ou não usar” IA, mas em definir sob quais critérios, controles e limites seu uso se torna compatível com o interesse público.
Além disso, a literatura internacional sobre implementação adverte que barreiras à adoção responsável são multidimensionais. Sharma et al. demonstram que desafios de implementação de IA em organizações públicas incluem restrições de infraestrutura, escassez de competências, dificuldades de integração com processos legados, resistência organizacional, incerteza estratégica e fragilidades de governança \autocite{sharma2021implementing}. Embora o estudo não se restrinja ao caso brasileiro, sua contribuição é particularmente útil para esta dissertação porque ajuda a deslocar o debate do plano normativo abstrato para o plano das condições organizacionais efetivas de implementação. Em outras palavras, não basta prescrever princípios éticos ou diretrizes gerais; é necessário compreender como tais princípios são traduzidos em rotinas, responsabilidades, controles, trilhas de auditoria e arranjos decisórios concretos.
No contexto latino-americano, essa preocupação ganha contornos adicionais. Criado argumenta que a expansão da IA no setor público da região ocorre em ambiente de capacidades estatais desiguais, pressões por modernização e marcos normativos ainda em consolidação, o que torna especialmente relevante a discussão sobre governança, direitos, transparência e alinhamento às orientações da Carta Ibero-Americana de Inteligência Artificial na Administração Pública \autocite{criado2024inteligencia}. Tal perspectiva comparada reforça a pertinência do caso federal brasileiro: trata-se de um espaço institucional em que se combinam ambição de transformação digital, centralidade regulatória e riscos de adoção fragmentada. Assim, o problema de pesquisa não consiste apenas em mapear aplicações de IA, mas em investigar como a literatura acadêmica e técnico-institucional tem tratado a relação entre usos concretos da tecnologia e os mecanismos necessários para governá-la de forma responsável.
Diante desse quadro, a questão que organiza esta subseção pode ser formulada nos seguintes termos: como a literatura acadêmica e técnico-institucional descreve, explica e avalia o uso de inteligência artificial na administração pública federal brasileira, especialmente no que se refere aos instrumentos de governança, às salvaguardas éticas, às exigências de compliance e às práticas de gestão de riscos necessárias para uma implementação responsável, auditável e compatível com o interesse público? A relevância dessa pergunta reside precisamente em enfrentar uma lacuna analítica: a produção existente já mostra aplicações, benefícios e obstáculos, mas ainda carece de maior integração entre os debates sobre adoção tecnológica, governança institucional, integridade administrativa e supervisão de riscos \autocite{pi2021machine} \autocite{sigfrids2022how} \autocite{mkander2023auditing}. É dessa insuficiência de síntese crítica e comparativa, aplicada ao recorte federal brasileiro, que emerge o problema de pesquisa da presente dissertação.
\subsection{Objetivo geral}
\label{sec:org1944a0f}
O objetivo geral desta dissertação é investigar, à luz da literatura acadêmica e técnico-institucional, como a inteligência artificial vem sendo concebida, adotada e problematizada no governo público federal brasileiro como instrumento de governança, ética, compliance e gestão de riscos, de modo a identificar não apenas suas principais aplicações, mas sobretudo as condições institucionais que tornam seu uso compatível com a legalidade, a transparência, a auditabilidade e a proteção do interesse público. Trata-se, portanto, de um objetivo que ultrapassa o mapeamento descritivo de iniciativas tecnológicas para alcançar uma compreensão analítica dos arranjos normativos, organizacionais e procedimentais que sustentam — ou fragilizam — a implementação responsável de sistemas de IA no setor público federal \autocite{pi2021machine} \autocite{madan2022ai} \autocite{criado2024artificial}.
Em termos substantivos, a investigação pretende examinar a adoção de IA em diferentes funções governamentais, como automação de processos, apoio à decisão, análise documental, fiscalização e oferta de serviços digitais, observando como essas aplicações reconfiguram práticas administrativas e deslocam o eixo do debate da eficiência para a governança. Esse deslocamento é central porque a literatura demonstra que os ganhos potenciais associados à IA em serviços governamentais — maior produtividade, melhor capacidade analítica, ampliação da personalização e rapidez na resposta estatal — coexistem com desafios relevantes relacionados à opacidade, à qualidade dos dados, à dependência tecnológica e à responsabilização por decisões automatizadas \autocite{alhosani2024opportunities} \autocite{pi2021machine}. Assim, o objetivo geral também envolve comparar benefícios prometidos e riscos efetivamente identificados, evitando tanto o entusiasmo tecnocrático quanto o ceticismo abstrato.
No recorte brasileiro, o estudo busca compreender como experiências setoriais já documentadas ajudam a iluminar o problema mais amplo da governança da IA na administração pública federal. O caso da saúde pública é especialmente relevante nesse sentido, pois evidencia que a utilização de IA pela Administração Pública brasileira já ocorre em contextos sensíveis, nos quais eficiência gerencial e impacto sobre direitos caminham conjuntamente, exigindo cautelas reforçadas de regulação, supervisão e controle \autocite{lemes2020o}. De forma análoga, a literatura sobre algoritmos preditivos em serviços públicos de emprego demonstra que sistemas utilizados para classificação, direcionamento e matching de usuários podem aperfeiçoar a alocação de recursos e a focalização de políticas, mas também intensificam problemas de discriminação indireta, perfilamento e redução da inteligibilidade das escolhas administrativas quando não submetidos a critérios claros de governança e revisão humana \autocite{krtner2021predictive}. Esses exemplos são relevantes para o objetivo geral porque oferecem parâmetros comparativos para interpretar o caso federal brasileiro a partir de aplicações concretas e não apenas de formulações normativas genéricas.
Além disso, esta dissertação pretende identificar como a literatura trata a exigência de explicabilidade, tema decisivo quando sistemas algorítmicos passam a influenciar decisões administrativas. A discussão desenvolvida por Brkan e Bonnet mostra que a explicação de decisões algorítmicas envolve tensões entre viabilidade técnica, inteligibilidade jurídica e efetividade do controle, indicando que a mera invocação de transparência não resolve, por si só, o problema da accountability em ambientes de alta complexidade computacional \autocite{brkan2020legal}. Desse modo, o objetivo geral inclui examinar em que medida os referenciais de governança, ética e compliance disponíveis são capazes de converter princípios abstratos em exigências operacionais, como documentação de modelos, definição de responsabilidades, trilhas de auditoria, monitoramento contínuo, avaliação de riscos e mecanismos de contestação.
Por fim, ao se constituir como revisão de literatura com foco no período de 2019 a 2026, o objetivo geral desta pesquisa é sistematizar criticamente o estado do conhecimento sobre IA no governo público federal brasileiro, articulando contribuições nacionais e internacionais para propor uma compreensão integrada do fenômeno. Busca-se, assim, oferecer uma base analítica que permita distinguir adoção tecnológica de adoção institucionalmente governada, evidenciando que a implementação de IA no setor público somente se justifica, em perspectiva republicana e administrativa, quando orientada por salvaguardas éticas, conformidade normativa, gestão de riscos e capacidade efetiva de escrutínio público \autocite{sigfrids2022how} \autocite{batool2025ai} \autocite{goulart2025processos}.
\subsection{Objetivos específicos}
\label{sec:org8ce9952}
Em coerência com o objetivo geral e com o recorte desta dissertação, os objetivos específicos foram definidos para decompor analiticamente um problema que não pode ser tratado apenas como questão de adoção tecnológica. A literatura recente sugere que, no setor público, a incorporação de inteligência artificial tende a produzir resultados muito distintos conforme a existência — ou ausência — de capacidades institucionais, arranjos decisórios, critérios de responsabilização e mecanismos de supervisão. Nesse sentido, a formulação dos objetivos específicos procura evitar um tratamento meramente descritivo das aplicações de IA, orientando a investigação para a identificação das condições sob as quais tais aplicações podem ser consideradas institucionalmente legítimas, administrativamente controláveis e compatíveis com o interesse público \autocite{madan2022ai} \autocite{criado2024artificial} \autocite{david2023understanding}.
O primeiro objetivo específico consiste em mapear, na literatura acadêmica e técnico-institucional, os principais usos de sistemas e aplicações de IA no governo público federal brasileiro, com atenção às funções de automação de processos, apoio à decisão, análise documental, fiscalização e prestação de serviços digitais. Esse mapeamento não se limita a enumerar casos de uso, mas busca compreender como diferentes aplicações redistribuem tarefas, redefinem fluxos administrativos e alteram a relação entre decisão humana, processamento algorítmico e entrega de políticas públicas. A relevância desse objetivo decorre do fato de que a literatura sobre estratégias de adoção de tecnologias digitais no setor público demonstra que a implementação não ocorre de maneira homogênea, sendo condicionada por capacidades organizacionais, prioridades institucionais e diferentes níveis de maturidade digital \autocite{david2023understanding} \autocite{alhosani2024opportunities}.
O segundo objetivo específico é identificar e sistematizar os modelos, instrumentos e processos de governança da IA descritos no corpus selecionado, especialmente aqueles relacionados à definição de responsabilidades, coordenação intersetorial, padronização procedimental, supervisão humana e controle organizacional. Aqui, a ênfase recai sobre a distinção entre adoção de IA como inovação operacional e adoção de IA como prática sujeita a governança. A revisão integrativa de Goulart et al. evidencia, no contexto brasileiro, que os processos de governança em órgãos públicos ainda se encontram em consolidação e que sua efetividade depende da articulação entre diretrizes formais, instâncias de decisão e rotinas de acompanhamento capazes de reduzir improvisação institucional e assimetrias de implementação \autocite{goulart2025processos}. Em diálogo com revisões mais amplas sobre governança de IA, esse objetivo permitirá examinar em que medida o debate brasileiro converge com referenciais internacionais ou apresenta lacunas próprias do ambiente federativo e administrativo nacional \autocite{batool2025ai} \autocite{sigfrids2022how}.
O terceiro objetivo específico consiste em analisar criticamente os dilemas éticos e os requisitos de compliance associados ao uso de IA no setor público federal, com especial atenção aos problemas de transparência, explicabilidade, vieses, proteção de dados, rastreabilidade e prestação de contas. A pertinência desse objetivo torna-se mais nítida quando se observam aplicações em áreas sensíveis, nas quais a inteligência artificial não apenas automatiza rotinas, mas influencia a priorização de recursos, a classificação de demandas e a formulação de diagnósticos administrativos. A revisão de Souza, Chagas e Bulgareli sobre o uso de IA no processo decisório de alocação de recursos na saúde pública brasileira é particularmente elucidativa ao mostrar que ganhos esperados de eficiência e precisão convivem com desafios de equidade, legitimidade decisória e necessidade de supervisão robusta \autocite{souza2025uso}. Assim, a investigação buscará compreender como a literatura traduz princípios éticos em exigências concretas de conformidade institucional e quais lacunas permanecem entre normatividade abstrata e operacionalização administrativa \autocite{brkan2020legal} \autocite{mkander2023auditing}.
O quarto objetivo específico é examinar como a literatura trata a gestão de riscos vinculada ao ciclo de vida dos sistemas de IA utilizados pelo poder público, incluindo riscos técnicos, jurídicos, organizacionais e reputacionais. Esse exame abrange desde a qualidade e governança dos dados até o monitoramento de desempenho, a auditoria contínua, os mecanismos de revisão e os protocolos para correção de falhas. A literatura internacional tem mostrado que a gestão de riscos em IA não pode ser dissociada de estruturas permanentes de monitoramento, sobretudo quando os sistemas são empregados em decisões administrativas de impacto relevante \autocite{minkkinen2022continuous} \autocite{mkander2023auditing} \autocite{parycek2023artificial}. No caso brasileiro, esse objetivo também permitirá verificar se a produção identificada avança da retórica da inovação para a institucionalização de salvaguardas verificáveis.
Por fim, um quinto objetivo específico consiste em integrar criticamente os achados da revisão para propor um quadro analítico de implementação responsável da IA no governo público federal brasileiro, articulando aplicações, governança, ética, compliance e gestão de riscos. Mais do que encerrar o estudo com uma síntese temática, pretende-se produzir uma interpretação capaz de explicitar relações entre capacidades institucionais, escolhas de desenho organizacional e exigências de controle público. Esse esforço de integração é importante porque a literatura revisada indica que a qualidade da adoção tecnológica depende menos da presença isolada de ferramentas avançadas e mais da coerência entre estratégia, supervisão, competências e accountability \autocite{david2023understanding} \autocite{goulart2025processos} \autocite{souza2025uso}.
\subsection{Justificativa do estudo}
\label{sec:orge07dcc3}
A justificativa desta dissertação decorre, em primeiro lugar, da crescente centralidade da inteligência artificial nas rotinas, capacidades e formas de intervenção do Estado, sem que esse movimento tenha sido acompanhado, no mesmo ritmo, por um amadurecimento equivalente dos referenciais analíticos e normativos aplicáveis ao contexto da administração pública federal brasileira. A literatura recente mostra que a adoção de IA no setor público deixou de ser hipótese prospectiva para se tornar objeto concreto de reorganização administrativa, com efeitos sobre automação de procedimentos, apoio à decisão, prestação de serviços, fiscalização e tratamento de grandes volumes de dados \autocite{madan2022ai} \autocite{serra2024a} \autocite{pi2021machine}. Contudo, o debate permanece frequentemente fragmentado entre promessas de eficiência, estudos setoriais e formulações principiológicas amplas, o que torna necessário um esforço de revisão que integre aplicações, governança, ética, compliance e gestão de riscos em uma mesma moldura interpretativa.
Essa lacuna é particularmente relevante porque, no setor público, a adoção de IA não pode ser avaliada apenas por métricas de desempenho operacional. Diferentemente de contextos privados, sistemas algorítmicos empregados por órgãos governamentais incidem sobre direitos, prioridades distributivas, critérios de elegibilidade, fluxos procedimentais e padrões de responsabilização administrativa. Nesse sentido, a literatura tem enfatizado que a questão central não é somente se a IA melhora a capacidade estatal, mas sob quais condições institucionais ela pode fazê-lo sem comprometer legalidade, transparência, equidade e controle democrático \autocite{criado2024artificial} \autocite{parycek2023artificial}. Revisões sistemáticas e integrativas apontam, ademais, que os processos de difusão da IA na administração pública são desiguais, condicionados por capacidades organizacionais, disponibilidade de dados, desenho regulatório, arranjos de coordenação e competências dos servidores, produzindo assimetrias de maturidade e de governabilidade entre instituições e setores \autocite{madan2022ai} \autocite{mariano2026inteligencia} \autocite{melo2026competencias}.
A importância acadêmica do estudo reside, portanto, em organizar criticamente um campo ainda em consolidação. A produção disponível já é suficiente para demonstrar a amplitude do fenômeno, mas ainda carece de maior sistematização quanto aos seus mecanismos explicativos e às suas implicações administrativas. Há revisões que destacam o papel da governança e das considerações éticas como condições para implementação responsável \autocite{serra2024a}; outras enfatizam estruturas de governança de IA em termos mais gerais, sublinhando princípios como accountability, fairness, transparência e supervisão humana \autocite{batool2025ai} \autocite{sigfrids2022how}. Também há contribuições que deslocam o foco para auditoria e monitoramento contínuo, argumentando que a confiança institucional em sistemas de IA depende menos de declarações abstratas de conformidade e mais de mecanismos verificáveis de inspeção, rastreabilidade e revisão \autocite{mkander2023auditing} \autocite{minkkinen2022continuous}. A relevância desta dissertação está em articular esses eixos em uma revisão orientada especificamente ao governo público federal brasileiro, onde a combinação entre complexidade administrativa, sensibilidade decisória e exigências de interesse público torna a discussão especialmente crítica.
A justificativa prática e institucional do estudo também é robusta. O avanço do uso de IA em administrações públicas tem sido acompanhado por desafios recorrentes relacionados à opacidade de modelos, qualidade dos dados, vieses discriminatórios, dificuldade de explicação das decisões automatizadas e insuficiência de capacidades internas para supervisão técnica e jurídica \autocite{brkan2020legal} \autocite{parycek2023artificial}. Em áreas como saúde, gestão de pessoas, serviços de emprego, análise documental e procedimentos administrativos, a literatura sugere que ganhos de celeridade e racionalização coexistem com riscos de reforço de desigualdades, erro sistêmico, dependência tecnológica e erosão da responsabilização humana \autocite{lemes2020o} \autocite{cho2023human} \autocite{krtner2021predictive} \autocite{chilunjika2022artificial} \autocite{barros2025inteligencia}. Assim, investigar o tema é relevante não apenas para descrever inovações, mas para oferecer base analítica capaz de subsidiar decisões públicas mais prudentes, auditáveis e institucionalmente legítimas.
Além disso, o recorte temporal de 2019 a 2026 é justificável porque corresponde ao período em que a IA, incluindo aplicações mais recentes de IA generativa, passou a ocupar posição mais estruturante nas agendas de transformação digital e modernização estatal, ao mesmo tempo em que se intensificaram os debates sobre governança, regulação e uso responsável \autocite{lyrio2025emerging} \autocite{mariano2026inteligencia}. No contexto latino-americano e ibero-americano, estudos comparados indicam que há difusão do tema, mas também heterogeneidade importante quanto à institucionalização de diretrizes, salvaguardas e capacidades de implementação, o que reforça a necessidade de examinar o caso brasileiro com maior precisão analítica \autocite{criado2024inteligencia}.
Por fim, esta dissertação se justifica por responder a uma necessidade simultaneamente científica e pública: produzir uma síntese qualificada sobre como a literatura acadêmica e técnico-institucional tem tratado a adoção de IA no governo federal brasileiro, identificando não só usos e promessas, mas também limites, pré-condições e salvaguardas. Em um cenário em que a incorporação de sistemas inteligentes tende a redefinir procedimentos, competências e formas de controle, a ausência de reflexão sistemática pode ampliar riscos justamente onde se exige maior prudência institucional. Ao reunir e comparar evidências sobre aplicações, modelos de governança, dilemas éticos, mecanismos de compliance e práticas de gestão de riscos, o estudo pretende contribuir para que a inovação tecnológica no setor público seja compreendida menos como imperativo de modernização acrítica e mais como objeto de governo, regulação e responsabilidade administrativa \autocite{goulart2025processos} \autocite{alhosani2024opportunities} \autocite{yiitcanlar2022artificial} \autocite{sharma2021implementing}.
\section{REFERENCIAL TEÓRICO}
\label{sec:org31b721f}
\subsection{Inteligência artificial e administração pública}
\label{sec:orgd4e6e66}
A relação entre inteligência artificial (IA) e administração pública deve ser compreendida para além da simples incorporação de uma nova tecnologia ao aparato estatal. A literatura recente indica que a IA altera rotinas administrativas, modos de produção de informação, formas de coordenação interorganizacional e critérios de decisão, produzindo implicações institucionais que alcançam tanto a eficiência operacional quanto a legitimidade da ação governamental \autocite{madan2022ai} \autocite{criado2024artificial}. Nesse sentido, a IA não aparece apenas como ferramenta instrumental de automação, mas como vetor de reconfiguração da capacidade administrativa, uma vez que interfere nos meios pelos quais o Estado coleta dados, processa demandas, prioriza casos, distribui recursos e interage com cidadãos e servidores \autocite{pi2021machine} \autocite{serra2024a}. Para uma dissertação voltada ao governo público federal brasileiro, esse enquadramento é decisivo, pois desloca o debate do plano tecnicista para uma análise propriamente político-administrativa, centrada nas condições de uso, nos efeitos organizacionais e nos limites institucionais da adoção algorítmica.
Uma primeira convergência importante na literatura diz respeito ao reconhecimento de que a difusão da IA no setor público tem sido impulsionada por agendas de transformação digital, pressão por produtividade, escassez relativa de recursos e crescente complexidade dos problemas governamentais \autocite{madan2022ai} \autocite{alhosani2024opportunities}. Em contraste com fases anteriores do governo eletrônico, a IA amplia a capacidade de tratar volumes massivos de dados, identificar padrões, realizar classificações, apoiar prognósticos e automatizar etapas procedimentais em escala \autocite{pi2021machine} \autocite{lyrio2025emerging}. Esse potencial explica seu avanço em áreas como triagem documental, atendimento ao usuário, fiscalização, saúde, gestão de pessoas e serviços públicos digitais. Todavia, os autores também convergem ao sustentar que o setor público constitui um ambiente de adoção singular: diferentemente das organizações privadas, as administrações públicas estão submetidas a exigências reforçadas de legalidade, motivação, impessoalidade, publicidade e controle, de modo que o sucesso da IA não pode ser medido apenas por indicadores de velocidade ou redução de custos \autocite{parycek2023artificial} \autocite{criado2024artificial}.
Essa distinção ajuda a entender por que parte da literatura critica leituras excessivamente solucionistas. Em revisões amplas do campo, Madan e Ashok mostram que a adoção e a difusão da IA na administração pública foram frequentemente discutidas sob uma lógica de promessa tecnológica, com ênfase em benefícios potenciais, mas com menor densidade explicativa sobre os fatores organizacionais e institucionais que condicionam a implementação efetiva \autocite{madan2022ai}. Serra, em revisão integrativa voltada especificamente à administração pública, reforça esse diagnóstico ao observar que o debate tende a destacar inovação, eficiência e melhoria de serviços, mas ainda necessita de maior sistematização sobre governança, ética e adequação institucional \autocite{serra2024a}. Há, portanto, uma lacuna entre a retórica de modernização e a análise concreta dos arranjos pelos quais sistemas de IA passam a participar de procedimentos administrativos, influenciando fluxos de trabalho e cadeias de responsabilização.
Essa lacuna se torna mais visível quando se observa a heterogeneidade das aplicações. Parycek, Schmid e Novak distinguem com precisão os diferentes graus de automação em procedimentos administrativos e advertem que nem toda utilização de IA implica substituição decisória plena \autocite{parycek2023artificial}. Em muitos casos, a tecnologia atua como mecanismo de apoio, recomendação, classificação ou pré-processamento, permanecendo a decisão formal com o agente público. Em outros, contudo, a automação pode adquirir papel mais substantivo ao estruturar a própria base informacional sobre a qual a decisão humana é tomada. Essa distinção é teoricamente relevante porque o impacto institucional da IA não depende apenas de sua presença, mas da posição que ocupa no ciclo decisório. Um sistema de triagem documental, por exemplo, pode parecer neutro do ponto de vista jurídico-administrativo; entretanto, se ele prioriza determinados tipos de casos, produz efeitos distributivos indiretos sobre tempos de resposta, atenção estatal e acesso a direitos. Assim, a análise da IA na administração pública exige considerar não só a decisão final, mas toda a arquitetura sociotécnica do procedimento.
A literatura também sugere que a IA opera em múltiplos níveis analíticos. Criado, Sandoval-Almazán e Gil-Garcia propõem compreender a relação entre IA e administração pública a partir de perspectivas micro, meso e macro, articulando atores individuais, organizações públicas e ambientes institucionais mais amplos \autocite{criado2024artificial}. No nível micro, ganham relevo as capacidades dos servidores, suas percepções sobre utilidade, confiança, risco e adequação da tecnologia. No nível meso, sobressaem desenho organizacional, governança interna, infraestrutura de dados e rotinas de implementação. No nível macro, entram em cena regulação, estratégias nacionais, ecossistemas de inovação e expectativas sociais sobre o papel do Estado. Essa formulação é particularmente útil porque evita tanto o reducionismo tecnológico quanto o institucionalismo abstrato: a IA é vista como fenômeno relacional, dependente da interação entre capacidades técnicas, estruturas administrativas e orientações normativas.
Em termos empíricos, as aplicações identificadas na literatura são amplas, mas desigualmente desenvolvidas. Em estudos de revisão, aparecem com frequência usos em atendimento automatizado, classificação de requerimentos, detecção de fraudes, análise preditiva, alocação de recursos, monitoramento regulatório e suporte à formulação de políticas \autocite{madan2022ai} \autocite{serra2024a} \autocite{alhosani2024opportunities}. Pi destaca que o aprendizado de máquina em governos tende a concentrar-se em tarefas intensivas em dados, nas quais há grande volume de registros administrativos e necessidade de reconhecer padrões não triviais \autocite{pi2021machine}. Essa constatação ajuda a explicar por que áreas como arrecadação, controle, saúde e assistência social se tornam terreno fértil para experimentação algorítmica. No entanto, a literatura diverge quanto ao grau de maturidade dessas experiências. Enquanto alguns trabalhos enfatizam a rápida expansão da IA como parte da transformação digital do setor público, outros ressaltam que muitos governos ainda se encontram em estágios iniciais, marcados por projetos-piloto, baixa integração institucional e forte dependência de fornecedores externos \autocite{yiitcanlar2022artificial} \autocite{david2023understanding} \autocite{mariano2026inteligencia}.
A noção de maturidade assimétrica é central para o recorte desta dissertação. Mariano, ao mapear IA e IA generativa na administração pública, enfatiza a existência de assimetrias de maturidade entre órgãos, setores e tipos de aplicação, indicando que a difusão tecnológica não ocorre de forma homogênea nem linear \autocite{mariano2026inteligencia}. Essa leitura dialoga com a revisão de Melo sobre competências em IA na gestão pública, segundo a qual a capacidade de adoção depende de habilidades técnicas, analíticas e gerenciais que ainda são escassas em muitos contextos administrativos \autocite{melo2026competencias}. Assim, a presença de IA em determinados nichos do Estado não autoriza inferir transformação generalizada da administração pública. Ao contrário, a literatura indica um quadro de adoção fragmentada, condicionado por disponibilidade de dados, infraestrutura, liderança, desenho institucional e capacidade de coordenação intersetorial \autocite{madan2022ai} \autocite{sharma2021implementing}.
Outro aspecto relevante diz respeito às promessas de eficiência e qualidade do serviço público. Yigitcanlar, Agdas e Degirmenci, ao examinar percepções de gestores municipais, mostram que a IA é frequentemente associada à melhoria de desempenho, agilidade, inovação e capacidade de resposta a ambientes complexos \autocite{yiitcanlar2022artificial}. Em chave semelhante, Alhosani e coautores ressaltam oportunidades relacionadas à personalização de serviços, otimização de processos e melhor uso de dados governamentais \autocite{alhosani2024opportunities}. Contudo, essa agenda de benefícios convive com obstáculos persistentes. Os mesmos estudos registram barreiras ligadas a custo, qualidade de dados, resistência organizacional, fragilidade regulatória e incertezas quanto a accountability e confiança pública \autocite{yiitcanlar2022artificial} \autocite{alhosani2024opportunities}. A contribuição desses trabalhos está em demonstrar que a adoção de IA não é meramente questão de disponibilidade tecnológica, mas de alinhamento entre capacidades organizacionais e valores públicos.
Em setores específicos, a literatura revela de forma ainda mais nítida a ambivalência entre potencial e risco. No campo da saúde, Lemes mostra que o uso de IA pela administração pública brasileira pode ampliar capacidades diagnósticas, otimizar fluxos e apoiar políticas sanitárias, mas também intensifica preocupações com sensibilidade dos dados, supervisão profissional e equidade no acesso aos benefícios tecnológicos \autocite{lemes2020o}. Na gestão de pessoas, Chilunjika identifica oportunidades para recrutamento, planejamento de força de trabalho e análise de desempenho no setor público, ao mesmo tempo em que chama atenção para riscos de desumanização, discriminação e inadequação contextual \autocite{chilunjika2022artificial}. Cho, ao discutir analytics para gestão de pessoal, reforça que sistemas orientados por dados podem apoiar decisões mais informadas, mas não eliminam os vieses presentes nos próprios registros administrativos e nas categorias organizacionais herdadas \autocite{cho2023human}. Já Barros, em discussão mais recente sobre IA e gestão de pessoas na administração pública, recoloca o problema em termos éticos e institucionais, argumentando que inovação sem salvaguardas pode comprometer confiança interna, justiça organizacional e legitimidade administrativa \autocite{barros2025inteligencia}.
Os serviços públicos de emprego oferecem um caso paradigmático dessa tensão. Ao analisar algoritmos preditivos na oferta de serviços de emprego, Körtner demonstra que ferramentas voltadas à segmentação e priorização de usuários podem aumentar focalização e eficiência administrativa, mas também correm o risco de cristalizar padrões excludentes se operarem com bases históricas enviesadas ou critérios opacos \autocite{krtner2021predictive}. Embora o estudo não se restrinja ao contexto brasileiro, seu argumento é altamente transferível para a administração pública federal, sobretudo em áreas nas quais decisões automatizadas ou semiautomatizadas influenciam acesso a políticas, benefícios ou oportunidades. O ponto central é que a IA, ao organizar a atenção administrativa por meio de probabilidades e classificações, tende a redefinir a própria gramática da intervenção estatal: casos deixam de ser tratados apenas como demandas singulares e passam a ser processados como perfis, categorias e riscos estimados.
Esse deslocamento afeta igualmente a compreensão do trabalho burocrático. Uma parte importante da literatura sustenta que a IA não elimina a burocracia, mas a reconfigura. Em vez de uma substituição simples de servidores por sistemas, o que emerge é uma redistribuição de tarefas entre humanos e máquinas, com novas exigências de interpretação, validação, curadoria de dados e supervisão \autocite{parycek2023artificial} \autocite{criado2024artificial}. Melo enfatiza justamente que a incorporação de IA na gestão pública demanda competências que combinam letramento de dados, compreensão de modelos, capacidade crítica e tradução entre linguagem técnica e jurídica \autocite{melo2026competencias}. A consequência é dupla. De um lado, servidores passam a depender de novas habilidades para usar a tecnologia de forma substantiva e não meramente ritualística. De outro, a qualidade do controle administrativo passa a depender menos da familiaridade abstrata com inovação e mais da capacidade concreta de questionar outputs, limites e pressupostos dos sistemas empregados.
No plano comparado, estudos sobre América Latina mostram que a adoção de IA no setor público ocorre em contexto de capacidades desiguais e marcos normativos em consolidação. Criado, em análise da região a partir da Carta Ibero-Americana de Inteligência Artificial na Administração Pública, sugere que o debate latino-americano combina forte expectativa modernizadora com desafios persistentes de institucionalização, coordenação e adaptação às realidades administrativas nacionais \autocite{criado2024inteligencia}. Essa observação é importante para o caso brasileiro: embora o país possua aparato administrativo relativamente sofisticado e agenda robusta de transformação digital, a adoção de IA no governo federal não pode ser presumida como uniforme, nem automaticamente alinhada a padrões consolidados de governança. O quadro mais plausível, em consonância com a literatura, é o de avanços setoriais coexistindo com lacunas de coordenação, capacidades díspares e tensão entre inovação acelerada e controles institucionais ainda em amadurecimento \autocite{goulart2025processos} \autocite{mariano2026inteligencia}.
Também merece destaque a emergência da IA generativa como novo fator de complexificação do campo. Lyrio e Mariano observam que a disseminação recente dessas aplicações altera o repertório de usos no setor público, expandindo possibilidades de apoio à redação, síntese documental, interação conversacional e produção de conteúdo \autocite{lyrio2025emerging} \autocite{mariano2026inteligencia}. Ao mesmo tempo, essa expansão amplia incertezas sobre confiabilidade factual, proteção de dados, rastreabilidade e uso adequado em atividades que exigem precisão técnico-jurídica. Embora a literatura ainda esteja em fase de consolidação, ela sugere que a IA generativa intensifica um problema já presente em outros sistemas: a distância entre facilidade de adoção e capacidade institucional de governança. Em outras palavras, quanto mais acessível e versátil a tecnologia, maior o risco de usos difusos sem critérios claros, especialmente em organizações públicas pressionadas por ganhos rápidos de produtividade.
Em síntese, o referencial teórico sobre inteligência artificial e administração pública aponta para uma conclusão fundamental: a IA deve ser tratada como fenômeno organizacional, decisório e institucional, e não apenas computacional \autocite{madan2022ai} \autocite{criado2024artificial} \autocite{serra2024a}. Seus efeitos não derivam exclusivamente da sofisticação técnica dos modelos, mas do modo como são inseridos em procedimentos, distribuídos entre órgãos, apropriados por servidores e legitimados perante a sociedade. A literatura converge quanto ao potencial da IA para ampliar capacidades estatais, mas diverge ou nuanceia fortemente as condições sob as quais esse potencial se converte em valor público. As principais lacunas residem justamente na compreensão de como arranjos de governança, capacidades burocráticas, controles institucionais e critérios éticos moldam a adoção efetiva da tecnologia no setor público. Por isso, esta subseção estabelece uma premissa que orienta as seguintes: discutir IA na administração pública implica examinar, de modo articulado, seus usos concretos e as estruturas que devem discipliná-los. É nessa passagem, do uso à regulação organizacional do uso, que se situam os debates sobre governança, ética, compliance e gestão de riscos desenvolvidos nas subseções posteriores.
\subsection{Governança da IA no setor público}
\label{sec:org3ec3239}
A governança da inteligência artificial no setor público pode ser definida, em sentido amplo, como o conjunto de arranjos institucionais, normas, processos decisórios, competências organizacionais e mecanismos de controle que orientam o desenvolvimento, a aquisição, a implementação, o monitoramento e a revisão de sistemas algorítmicos usados por órgãos estatais. Essa definição é mais exigente do que a mera gestão de projetos tecnológicos, porque incorpora uma preocupação central com legitimidade administrativa, conformidade jurídica, capacidade institucional e preservação do interesse público. A literatura recente converge ao sustentar que, no setor público, a governança da IA não pode ser tratada como camada acessória posterior à inovação; ela constitui condição de possibilidade para que a adoção tecnológica seja compatível com legalidade, impessoalidade, transparência, motivação decisória e controle democrático \autocite{sigfrids2022how} \autocite{batool2025ai} \autocite{serra2024a}.
Essa compreensão representa um deslocamento importante em relação a abordagens iniciais que descreviam a IA sobretudo pelos seus ganhos funcionais. Revisões de escopo amplo mostram que a literatura sobre adoção de IA na administração pública cresceu em torno de promessas de eficiência, automação e melhoria de serviços, mas amadureceu progressivamente para reconhecer que tais benefícios dependem de estruturas de governança capazes de disciplinar usos, delimitar responsabilidades e mitigar efeitos adversos \autocite{madan2022ai} \autocite{alhosani2024opportunities}. Em outras palavras, o problema deixou de ser apenas “como introduzir IA” e passou a incluir “sob quais condições institucionais essa introdução é aceitável e governável”. Esse ponto é decisivo para o recorte desta dissertação, porque o governo público federal brasileiro opera em ambiente normativo e procedimental no qual decisões automatizadas ou semi-automatizadas podem produzir efeitos distributivos, probatórios e reputacionais relevantes, exigindo mais do que coordenação técnica entre áreas de tecnologia.
A revisão sistemática de Batool et al. propõe uma visão abrangente da governança da IA ao mostrar que o campo se consolidou em torno de eixos como accountability, transparência, fairness, gestão de riscos, supervisão humana, regulação e auditoria \autocite{batool2025ai}. Embora esses elementos apareçam em diferentes combinações conforme o contexto, a contribuição do estudo está em evidenciar que governança da IA é um conceito necessariamente multidimensional. Não se trata apenas de compliance regulatório nem apenas de desenho ético; trata-se da articulação entre instrumentos organizacionais e normativos para orientar o ciclo de vida dos sistemas. Essa leitura converge com Sigfrids et al., que, ao examinar propostas para o desenvolvimento ético e o uso responsável de IA por administrações públicas, enfatizam que princípios abstratos somente se tornam operacionais quando traduzidos em estruturas concretas de governança, com atribuições definidas, processos de avaliação, participação de múltiplos atores e capacidade de revisão institucional \autocite{sigfrids2022how}.
Nesse ponto, uma distinção analítica útil é entre governança “da” IA e governança “por meio da” IA. A primeira diz respeito ao controle da tecnologia pelo Estado e pelas organizações públicas; a segunda, ao uso da própria IA como instrumento de governar processos, pessoas e políticas. Criado, Sandoval-Almazán e Gil-Garcia ajudam a esclarecer essa dupla dimensão ao propor uma leitura em níveis micro, meso e macro, na qual atores individuais, organizações e ecossistemas institucionais interagem na conformação dos usos públicos da IA \autocite{criado2024artificial}. A governança da IA no setor público, portanto, não se esgota em documentos estratégicos ou em comitês centrais; ela depende de como gestores, analistas, juristas, equipes de dados, fornecedores e instâncias de controle distribuem autoridade e interpretam riscos em contextos concretos de implementação. Essa perspectiva evita tratar a governança como sinônimo de centralização normativa absoluta e permite compreender que o problema central é a coordenação legítima entre múltiplos centros decisórios.
A literatura também mostra que a governança da IA precisa ser sensível ao tipo de função administrativa afetada. Parycek, Schmid e Novak, ao analisar IA e automação em procedimentos administrativos, sustentam que o desenho institucional adequado varia conforme o grau de automação, a natureza do procedimento e o peso jurídico da decisão \autocite{parycek2023artificial}. Sistemas usados para apoio interno, como classificação documental ou priorização de filas, já requerem critérios de qualidade de dados, rastreabilidade e supervisão. Porém, quando a IA intervém em procedimentos com impactos diretos sobre direitos, benefícios, sanções ou elegibilidade, as exigências de governança tornam-se mais intensas. A contribuição desse argumento é mostrar que não existe modelo único de governança aplicável indistintamente a todos os usos. Quanto mais próximo o sistema estiver do núcleo decisório estatal, maior deve ser a densidade de garantias institucionais, inclusive quanto à documentação do modelo, à explicabilidade suficiente para fins administrativos e à possibilidade de contestação.
Serra, em revisão integrativa da literatura sobre IA na administração pública, reforça esse ponto ao indicar que a expansão das aplicações governamentais vem acompanhada de preocupações crescentes com governança adequada e considerações éticas \autocite{serra2024a}. O mérito do estudo está em não reduzir a governança a um problema técnico de implantação, mas situá-la como requisito para sistematizar o avanço da IA no setor público. Essa sistematização é especialmente importante em administrações complexas e fragmentadas, como a federal, nas quais diferentes órgãos podem adotar soluções algorítmicas com níveis distintos de maturidade, documentação, controles e competências internas. Sem coordenação mínima, cria-se um cenário de assimetria institucional: algumas unidades desenvolvem práticas robustas de avaliação e supervisão, enquanto outras operam com baixa capacidade de governança, ampliando riscos de arbitrariedade, dependência de fornecedores ou uso opaco da automação.
Esse diagnóstico se aproxima das evidências de Goulart sobre processos de governança de IA em órgãos públicos brasileiros. Ainda que o campo nacional esteja em formação, a revisão integrativa indica que a governança tem sido discutida em torno da necessidade de políticas institucionais, definição de papéis, integração entre áreas finalísticas e tecnológicas, e criação de fluxos formais para avaliação de sistemas \autocite{goulart2025processos}. O aspecto mais relevante, para esta dissertação, é que o debate brasileiro parece menos centrado em modelos acabados e mais em processos de construção institucional. Isso sugere que, no caso da administração pública federal, a governança da IA deve ser compreendida como trajetória de maturação organizacional, e não como simples importação de frameworks prontos oriundos do setor privado ou de organismos internacionais.
A literatura internacional, contudo, oferece referenciais analíticos importantes para essa trajetória. Madan e Ashok identificam que a adoção e difusão de IA na administração pública dependem de fatores como disponibilidade de dados, infraestrutura digital, liderança, apoio político, capacidades organizacionais e compatibilidade com normas e rotinas existentes \autocite{madan2022ai}. Embora o foco do estudo não seja exclusivamente governança, suas conclusões demonstram que a qualidade da adoção é inseparável do ambiente institucional em que a tecnologia é inserida. Em termos comparativos, isso significa que governos não falham apenas por insuficiência técnica, mas por ausência de arranjos de coordenação e aprendizado capazes de integrar inovação com accountability. Já David, ao revisar estratégias de adoção de tecnologias digitais em governos locais, mostra que a adoção bem-sucedida depende de combinação entre estratégia, capacidades internas e adaptação institucional \autocite{david2023understanding}. Transposto para a IA, esse argumento reforça que governança eficaz pressupõe alinhamento entre objetivo público, arquitetura organizacional e mecanismos de controle.
Outra convergência robusta na literatura refere-se à centralidade das capacidades estatais. Governar IA requer competências que não são apenas computacionais. Melo destaca, em revisão sistemática, que a gestão pública precisa desenvolver competências em IA que incluam compreensão técnica básica, leitura crítica de modelos, interpretação de resultados, avaliação de riscos e capacidade de coordenação entre áreas \autocite{melo2026competencias}. Sem esse repertório, a organização pública tende a terceirizar também o juízo sobre adequação, tornando-se dependente de fornecedores ou de pequenos núcleos especializados. Essa dependência fragiliza a governança porque desloca o poder efetivo de decisão para atores que nem sempre respondem diretamente aos princípios da administração pública. Em sentido semelhante, estudos sobre recursos humanos e analytics no setor público advertem que a introdução de sistemas analíticos sem qualificação institucional suficiente pode reproduzir vieses e obscurecer critérios de gestão \autocite{cho2023human} \autocite{chilunjika2022artificial} \autocite{barros2025inteligencia}. A lição mais ampla é que governança da IA é também governança de capacidades humanas.
A dimensão organizacional da governança inclui, ademais, a definição de estruturas responsáveis por decisões ao longo do ciclo de vida dos sistemas. A revisão de Batool et al. mostra que o campo tem valorizado mecanismos como comitês de ética ou governança, políticas corporativas de IA, avaliações de impacto, padrões de documentação e procedimentos de monitoramento contínuo \autocite{batool2025ai}. Todavia, a literatura ainda diverge sobre o grau ideal de formalização. Parte dos autores entende que estruturas centrais mais fortes favorecem padronização e coerência; outra parte alerta que modelos excessivamente centralizados podem reduzir adaptabilidade e desconsiderar especificidades setoriais \autocite{sigfrids2022how} \autocite{criado2024artificial}. No setor público, essa divergência é particularmente relevante porque a heterogeneidade de missões institucionais é grande: a governança adequada para uma autarquia reguladora não coincide integralmente com a de um órgão de prestação massiva de serviços digitais ou de uma entidade voltada à fiscalização. A saída mais promissora, sugerida indiretamente por vários estudos, parece residir em modelos híbridos: princípios comuns e instâncias centrais de orientação combinados com responsabilização distribuída e controles contextualizados nas unidades usuárias.
A governança da IA também deve lidar com a materialidade dos dados e dos modelos. Parycek, Schmid e Novak assinalam que a validade das previsões e classificações depende criticamente da disponibilidade, qualidade e adequação dos dados utilizados \autocite{parycek2023artificial}. Em consequência, governar IA implica governar bases de dados, critérios de rotulação, processos de atualização, interoperabilidade e integridade informacional. Essa observação desloca o debate para um ponto às vezes negligenciado: muitos problemas atribuídos genericamente à “IA” decorrem, na verdade, de governança deficiente dos dados, de bases históricas enviesadas ou de usos indevidos fora do contexto original de treinamento. Para administrações públicas, isso é especialmente sensível porque registros estatais são produzidos em arranjos burocráticos diversos, com assimetrias de cobertura, temporalidade e qualidade. Logo, a governança da IA precisa incluir governança de dados como dimensão constitutiva, e não periférica.
Além disso, a literatura vem enfatizando a auditoria como mecanismo estruturante da governança. Mökander sustenta que auditorias de IA podem cumprir função semelhante à de auditorias financeiras, de segurança ou de conformidade em outros domínios, desde que adequadas às especificidades técnicas e sociotécnicas dos sistemas \autocite{mkander2023auditing}. A relevância dessa proposição está em oferecer uma ponte entre princípios abstratos e verificações concretas. Auditoria, nesse sentido, não é apenas inspeção posterior, mas parte de um regime institucional de asseguração, capaz de testar documentação, desempenho, aderência a critérios jurídicos e riscos de discriminação ou opacidade. Minkkinen, por sua vez, argumenta em favor da auditoria contínua de IA, destacando que modelos e ambientes de dados mudam ao longo do tempo, o que torna insuficiente uma avaliação única no momento da implantação \autocite{minkkinen2022continuous}. Para o setor público, essa conclusão é particularmente importante: se políticas públicas, bases administrativas e perfis de usuários se alteram, a governança precisa ser dinâmica, com monitoramento reiterado, thresholds de alerta e gatilhos de revisão.
Essa ênfase na auditabilidade revela uma convergência entre governança e prestação de contas, mas também expõe lacunas. Nem toda aplicação de IA no setor público é hoje suficientemente documentada para ser auditável em sentido forte. A revisão de Mariano sobre IA e IA generativa na administração pública aponta assimetrias de maturidade entre organizações e usos, o que sugere diferenças acentuadas na capacidade de registrar decisões de projeto, justificar escolhas de modelos e acompanhar efeitos em produção \autocite{mariano2026inteligencia}. Em outras palavras, a governança da IA não enfrenta apenas o desafio de definir princípios; enfrenta também o de institucionalizar rotinas documentais, métricas e trilhas de decisão que permitam controle posterior. Essa lacuna é crítica no contexto federal brasileiro, em que a multiplicidade de órgãos, contratações e soluções pode dificultar visibilidade sistêmica sobre quem usa o quê, para qual finalidade e sob quais salvaguardas.
No plano comparado, Criado et al. sobre o setor público latino-americano ressaltam que a governança da IA na região é atravessada por assimetrias institucionais, capacidades administrativas desiguais e diferentes estágios de formulação normativa \autocite{criado2024inteligencia}. A importância desse estudo para o presente trabalho está em relativizar a transposição automática de referenciais desenvolvidos em contextos com maior infraestrutura digital e tradição regulatória consolidada. Em países latino-americanos, inclusive o Brasil, a governança da IA precisa ser pensada em articulação com desafios mais amplos de coordenação estatal, profissionalização burocrática, transformação digital e desigualdades de acesso. Isso não reduz a relevância dos princípios internacionais; ao contrário, exige sua tradução institucional. O risco, aqui, é adotar retóricas sofisticadas de IA responsável sem construir os meios administrativos para torná-las exequíveis.
As evidências de governos locais e contextos subnacionais também enriquecem essa discussão. Yigitcanlar, ao analisar percepções de gestores municipais, mostra que as escolhas sobre IA são condicionadas por recursos, competências, percepção de risco e objetivos estratégicos \autocite{yiitcanlar2022artificial}. Embora o nível federativo seja distinto, a conclusão é transferível: governança não decorre apenas de normas superiores, mas da capacidade concreta de organizações públicas de interpretar, adaptar e operacionalizar diretrizes. Sharma, ao estudar desafios de implementação em setor público de economia emergente, reforça que barreiras como resistência organizacional, infraestrutura insuficiente, déficits de qualificação e ambiguidades regulatórias afetam diretamente a sustentação da inovação \autocite{sharma2021implementing}. Assim, qualquer modelo de governança da IA que desconsidere condicionantes materiais e organizacionais corre o risco de permanecer nominal.
Em síntese, a literatura permite sustentar que a governança da IA no setor público constitui um campo de interseção entre estratégia estatal, desenho institucional, capacidade burocrática, governança de dados, mecanismos de auditoria e princípios normativos de legitimidade administrativa. Há convergência quanto à necessidade de supervisão, transparência procedimental, responsabilização e avaliação contínua, mas permanecem divergências sobre o grau de centralização, o nível de formalização e a melhor forma de adaptar modelos de governança a diferentes funções estatais \autocite{batool2025ai} \autocite{sigfrids2022how} \autocite{parycek2023artificial}. Para o caso do governo público federal brasileiro, essa literatura sugere que a questão central não é apenas expandir o uso de IA, mas construir arranjos capazes de orientar sua adoção de modo auditável, proporcional ao risco, compatível com o direito administrativo e sustentado por competências institucionais distribuídas. Essa conclusão prepara o passo seguinte da análise, pois evidencia que governança da IA, no setor público, só se completa quando articulada às exigências de ética, accountability e explicabilidade que condicionam sua legitimidade.
\subsection{Ética, accountability e explicabilidade}
\label{sec:org2fab9ee}
A discussão sobre ética, accountability e explicabilidade da inteligência artificial no setor público ocupa posição central no debate contemporâneo porque, diferentemente do uso privado de sistemas algorítmicos, a administração pública exerce competências que afetam direitos, distribuem benefícios, impõem restrições, ordenam prioridades e produzem efeitos jurídicos e sociais em larga escala. Por isso, a questão ética não pode ser reduzida a um apêndice reputacional da inovação tecnológica. No contexto governamental, ela se relaciona diretamente com a legitimidade do agir administrativo, com a possibilidade de controle democrático e com a preservação dos princípios que estruturam a atuação estatal. A literatura recente converge nesse ponto ao sustentar que a adoção de IA no setor público somente é defensável quando acompanhada de salvaguardas institucionais capazes de tornar seus usos justificáveis, auditáveis e contestáveis \autocite{sigfrids2022how} \autocite{serra2024a} \autocite{batool2025ai}.
Esse deslocamento é relevante porque os estudos iniciais sobre IA na administração pública enfatizavam, com frequência, ganhos de produtividade, automação e melhoria de serviços, mas dedicavam menos atenção às assimetrias de poder que emergem quando algoritmos passam a influenciar procedimentos estatais. As revisões sistemáticas mostram que o campo amadureceu ao reconhecer que eficiência administrativa, isoladamente, não constitui critério suficiente para legitimar sistemas algorítmicos em governos \autocite{madan2022ai} \autocite{alhosani2024opportunities}. Em termos analíticos, isso significa admitir que o desempenho técnico de um modelo não esgota a avaliação pública de sua aceitabilidade. Um sistema pode ser acurado e ainda assim eticamente problemático se reproduzir vieses, obscurecer critérios decisórios, enfraquecer a responsabilização de agentes públicos ou dificultar a revisão de decisões que afetam cidadãos.
A ética da IA no setor público, portanto, precisa ser compreendida em chave institucional. Batool et al. identificam, na literatura de governança da IA, a recorrência de princípios como fairness, accountability, transparência, explicabilidade, privacidade, segurança e supervisão humana, observando que esses elementos aparecem como componentes de um mesmo esforço de responsabilização tecnológica \autocite{batool2025ai}. Contudo, um dos achados mais importantes dessa literatura é que tais princípios não operam de forma automática nem harmoniosa. Há tensões persistentes entre eficiência e devido processo, entre precisão estatística e inteligibilidade, entre proteção de dados e abertura informacional, e entre autonomia técnica de especialistas e controle por instâncias jurídicas, administrativas e sociais. A principal dificuldade, assim, não está em enunciar valores desejáveis, mas em traduzi-los em arranjos de decisão, documentação, monitoramento e revisão que resistam ao uso cotidiano da tecnologia.
Nesse sentido, Sigfrids et al. argumentam que as administrações públicas devem fomentar o desenvolvimento ético e o uso responsável da IA por meio de mecanismos organizacionais concretos, e não apenas por cartas de princípios ou códigos genéricos \autocite{sigfrids2022how}. A contribuição dessa abordagem é mostrar que a ética pública da IA depende de institucionalização. Isso envolve definir responsáveis pelo sistema ao longo de todo o seu ciclo de vida, estabelecer critérios de admissibilidade para usos sensíveis, prever avaliação prévia de impacto, registrar bases de dados, hipóteses de modelagem, objetivos do sistema e limitações conhecidas, além de assegurar vias de contestação para os indivíduos afetados. Em outras palavras, a ética deixa de ser uma camada discursiva e passa a constituir uma gramática operacional da decisão pública mediada por algoritmos.
A noção de accountability é particularmente decisiva nesse quadro. Em administração pública, accountability não se confunde com mera prestação de contas formal; ela abrange a obrigação de explicar, justificar e assumir responsabilidade por decisões, processos e resultados perante diferentes fóruns de controle, internos e externos. Quando sistemas de IA são inseridos em fluxos decisórios, surge o risco de difusão da responsabilidade entre gestores, desenvolvedores, fornecedores, analistas de dados e operadores finais. Criado, Sandoval-Almazán e Gil-Garcia ajudam a compreender esse problema ao enfatizar que a IA em administração pública deve ser analisada nos níveis micro, meso e macro, isto é, considerando atores individuais, organizações e ambientes institucionais mais amplos \autocite{criado2024artificial}. A accountability algorítmica, nesse contexto, não pode ser atribuída apenas ao agente que “clicou” na decisão final nem apenas ao fornecedor que desenvolveu o modelo. Ela exige mapear como competências, discricionariedades e dependências técnicas são distribuídas ao longo da cadeia decisória.
Esse ponto é especialmente relevante para o governo público federal brasileiro, no qual a fragmentação administrativa, a contratação de soluções tecnológicas de múltiplos fornecedores e a coexistência entre equipes jurídicas, de negócio e de tecnologia podem tornar opaca a identificação de quem responde por falhas algorítmicas. A literatura comparada sugere que, sem definição explícita de papéis, a IA tende a produzir um paradoxo institucional: sistemas são adotados para apoiar decisões de alto impacto, mas, quando surgem erros, discriminações ou resultados injustificáveis, a responsabilidade é deslocada para a “complexidade técnica” do modelo \autocite{mkander2023auditing} \autocite{minkkinen2022continuous}. Em vez de reduzir a accountability, a incorporação de IA deveria ampliar os deveres de documentação e justificativa, precisamente porque a mediação algorítmica dificulta a percepção imediata dos critérios empregados.
A explicabilidade entra nesse debate como dimensão normativa e prática da accountability. No plano mais básico, explicar um sistema de IA significa oferecer razões inteligíveis sobre como ele produz classificações, recomendações, priorizações ou previsões. No entanto, a literatura distingue diferentes níveis de explicação. Há explicações técnicas, úteis para especialistas entenderem estrutura do modelo, variáveis e parâmetros; explicações procedimentais, voltadas a esclarecer como o sistema foi treinado, validado e integrado ao processo administrativo; e explicações decisórias, necessárias para que a pessoa afetada compreenda por que determinada conclusão foi produzida em seu caso \autocite{brkan2020legal} \autocite{parycek2023artificial}. Essa distinção é crucial, porque grande parte do debate público sobre “direito à explicação” fracassa ao presumir que uma única forma de transparência satisfaz todas as exigências jurídicas e éticas.
Brkan, ao discutir a viabilidade jurídica e técnica da busca por explicações de decisões algorítmicas no contexto do GDPR, é cética quanto à ideia de que sistemas complexos possam sempre fornecer explicações plenas, individualizadas e tecnicamente satisfatórias \autocite{brkan2020legal}. Sua análise é importante para esta dissertação porque impede uma solução simplista. Exigir explicabilidade não equivale a supor que toda caixa-preta possa ser integralmente “aberta” de maneira compreensível para qualquer público. Em muitos casos, sobretudo em modelos de aprendizagem de máquina de alta complexidade, o objetivo mais realista não é uma transparência total sobre a arquitetura interna, mas a combinação entre rastreabilidade processual, documentação robusta, justificativas contextualizadas, testes de desempenho e mecanismos efetivos de contestação. Para a administração pública, isso implica reconhecer que a explicabilidade relevante é a que sustenta controle e revisão, não necessariamente a que satisfaz uma idealização de transparência absoluta.
Parycek, Schmid e Novak reforçam esse argumento ao mostrarem que a adequação do uso de IA em procedimentos administrativos depende das condições de enquadramento jurídico e organizacional de cada caso \autocite{parycek2023artificial}. Quanto mais intenso o impacto do sistema sobre direitos subjetivos, elegibilidade a benefícios, sanções ou acesso a serviços essenciais, maior deve ser a exigência de explicabilidade e de intervenção humana qualificada. A contribuição dos autores está em rejeitar a ideia de que automação administrativa seja uniforme. Há usos relativamente benignos, como triagem documental e apoio à organização interna, nos quais o ônus explicativo pode ser menor, embora não inexistente. Mas há também usos sensíveis, nos quais a combinação entre opacidade técnica e assimetria de poder estatal pode comprometer o devido processo administrativo. Nesses casos, o déficit de explicabilidade deixa de ser apenas problema técnico e se converte em problema de legitimidade.
A literatura sobre adoção de IA em governos também evidencia que os dilemas éticos não decorrem apenas do algoritmo, mas do contexto organizacional em que ele é aplicado. Madan e Ashok observam que a difusão de IA na administração pública é condicionada por fatores institucionais, capacidade organizacional, disponibilidade de dados, pressões por desempenho e estruturas de governança \autocite{madan2022ai}. Essa observação é importante porque sistemas algorítmicos não operam em abstração. Eles são treinados com bases históricas, integrados a processos existentes e utilizados por agentes cujas competências e incentivos influenciam os resultados. Assim, preconceitos históricos, desigualdades administrativas e assimetrias no acesso a dados podem ser convertidos em padrões aparentemente neutros de decisão automatizada. A ética da IA pública exige, portanto, analisar não só o comportamento do modelo, mas as condições institucionais que produzem e legitimam seus outputs.
Essa preocupação aparece com nitidez em aplicações voltadas à gestão de pessoas e recursos humanos no setor público. Estudos sobre analytics e IA em gestão de pessoal alertam que sistemas de classificação, ranqueamento, recrutamento ou avaliação de desempenho podem reforçar discriminações preexistentes, sobretudo quando utilizam dados históricos enviesados ou proxies socialmente sensíveis \autocite{cho2023human} \autocite{chilunjika2022artificial} \autocite{barros2025inteligencia}. Embora parte dessa literatura esteja situada fora do governo federal brasileiro, ela oferece advertências substantivas para a administração pública em geral. Em ambientes marcados por forte formalização e por efeitos de carreira, a introdução de modelos preditivos sem transparência suficiente pode corroer confiança organizacional, dificultar contestação por servidores e naturalizar decisões discriminatórias sob a aparência de neutralidade estatística. A lacuna, nesse campo, não está apenas em identificar riscos, mas em desenvolver critérios públicos para definir quais usos de IA são aceitáveis em gestão de pessoas e quais devem ser restringidos ou vedados.
Problema semelhante emerge nos serviços públicos orientados por modelos preditivos. Kärtner mostra que algoritmos aplicados à intermediação e entrega de serviços de emprego podem alterar prioridades de atendimento, classificar perfis de risco e influenciar a distribuição de apoio estatal \autocite{krtner2021predictive}. Ainda que apresentados como instrumentos de eficiência e focalização, tais sistemas levantam questões sobre opacidade dos critérios de classificação, possibilidade de profiling discriminatório e capacidade dos cidadãos de compreender ou contestar a lógica empregada. Para o setor público federal, isso é particularmente sensível em políticas distributivas, fiscalização, triagem de demandas e priorização de casos. O ponto central é que a accountability algorítmica não deve ser acionada apenas após incidentes; ela precisa estar embutida no desenho das aplicações, sobretudo quando a tecnologia interfere na alocação de atenção administrativa e recursos públicos.
A auditoria de IA aparece, nesse debate, como ponte entre ética abstrata e accountability verificável. Mökander argumenta que a auditoria de IA constitui campo em rápida consolidação e que suas abordagens podem combinar perspectivas legais, éticas e técnicas \autocite{mkander2023auditing}. Seu argumento de que a auditoria deve aprender com tradições mais maduras, como auditoria financeira, engenharia de segurança e métodos das ciências sociais, é particularmente fecundo para o setor público. Isso porque a confiabilidade de um sistema não pode ser inferida apenas da autodeclaração do fornecedor ou da equipe interna responsável. É preciso construir rotinas independentes ou, ao menos, funcionalmente separadas de verificação sobre qualidade de dados, desempenho, robustez, vieses, aderência normativa e coerência entre finalidade declarada e uso efetivo do sistema. Em administração pública, a auditoria também precisa considerar se o sistema está compatível com deveres de motivação, publicidade adequada, revisão administrativa e controle externo.
Minkkinen, por sua vez, aprofunda a noção de auditoria contínua de IA ao sugerir que o controle não deve ocorrer apenas na fase ex ante ou em avaliações pontuais \autocite{minkkinen2022continuous}. Como modelos podem degradar com o tempo, sofrer drift, responder a mudanças no ambiente ou produzir efeitos inesperados em novas populações, a accountability exige monitoramento permanente. Essa contribuição é central para distinguir governança formal de governança efetiva. Um órgão pode aprovar diretrizes éticas sofisticadas e ainda assim falhar se não dispuser de métricas, indicadores, registros e capacidade analítica para acompanhar o comportamento do sistema em operação. Em termos práticos, a explicabilidade relevante para a auditoria contínua envolve não só explicações sobre decisões passadas, mas visibilidade sistemática sobre padrões de erro, grupos afetados, frequência de revisão humana e desvios entre objetivos declarados e resultados observados.
A literatura sobre administrações locais e contextos comparados reforça que a percepção dos gestores sobre IA costuma oscilar entre entusiasmo instrumental e prudência ética. Yigitcanlar, Agdas e Degirmenci mostram que gestores municipais reconhecem potencial de melhoria de serviços e ganho de eficiência, mas também identificam obstáculos ligados a confiança, transparência, qualidade dos dados e aceitação social \autocite{yiitcanlar2022artificial}. David, em revisão sobre estratégias de adoção tecnológica por governos locais, igualmente sugere que escolhas tecnológicas são inseparáveis de capacidade institucional, cultura organizacional e alinhamento estratégico \autocite{david2023understanding}. Esses estudos, embora não centrados no nível federal brasileiro, ajudam a evidenciar que ética e accountability não são exigências externas impostas à inovação; elas condicionam a sustentabilidade política e administrativa da própria adoção. Sistemas percebidos como arbitrários, opacos ou injustos tendem a encontrar resistência interna, litígios e perda de legitimidade.
No contexto latino-americano, Criado et al. observam assimetrias relevantes na incorporação da IA no setor público, tanto em termos de capacidade estatal quanto de desenvolvimento de marcos orientadores \autocite{criado2024inteligencia}. Essa constatação é relevante para o caso brasileiro porque indica que a agenda de ética e explicabilidade não pode ser transplantada mecanicamente de contextos com alta maturidade regulatória e tecnológica. Em países e administrações com capacidades heterogêneas, o desafio é duplo: evitar tanto a adoção acrítica de soluções opacas quanto a importação de padrões normativos cuja implementação prática seja inviável sem investimentos em competências, infraestrutura de dados e coordenação institucional. Nessa linha, Melo destaca que competências em IA na gestão pública são parte constitutiva de uma implementação responsável, não mero complemento técnico \autocite{melo2026competencias}. Sem capacidades para compreender limitações do modelo, interpretar relatórios, formular requisitos e exercer supervisão, princípios éticos permanecem retóricos.
No Brasil, revisões sobre governança e maturidade de IA em órgãos públicos apontam avanços discursivos e normativos, mas também heterogeneidade institucional e assimetrias relevantes entre organizações \autocite{goulart2025processos} \autocite{mariano2026inteligencia}. Isso sugere que a ética da IA no governo federal não deve ser pensada apenas como conformidade com princípios gerais, e sim como problema de implementação diferenciada. Órgãos com maior capacidade técnica podem estruturar trilhas de documentação, testes e monitoramento mais sofisticadas; outros podem depender fortemente de fornecedores ou de soluções prontas, ampliando riscos de opacidade e lock-in informacional. O desafio teórico e prático, portanto, consiste em formular padrões mínimos comuns de accountability e explicabilidade que sejam exigíveis em toda a administração, ao mesmo tempo em que se reconhece a necessidade de gradação conforme o risco, a finalidade e o impacto do uso algorítmico.
Em síntese, a literatura revisada indica uma convergência importante: ética, accountability e explicabilidade devem ser tratadas como dimensões integradas da governança da IA no setor público, e não como agendas isoladas. A ética fornece os critérios substantivos de justiça, não discriminação, dignidade e orientação ao interesse público; a accountability organiza os deveres de justificação, responsabilização e controle; e a explicabilidade torna praticável, ao menos em parte, a compreensão e a contestação de sistemas que influenciam decisões administrativas \autocite{sigfrids2022how} \autocite{batool2025ai} \autocite{mkander2023auditing}. Persistem, contudo, lacunas relevantes. A literatura ainda carece de maior detalhamento sobre como transformar princípios em rotinas administrativas estáveis, como calibrar níveis adequados de explicação conforme o risco da aplicação e como distribuir responsabilidades em ecossistemas sociotécnicos complexos. Para o recorte desta dissertação, essa lacuna é decisiva: no governo público federal brasileiro, a questão não é apenas se a IA pode ser usada, mas sob quais condições ela pode ser usada de modo eticamente aceitável, institucionalmente responsável e efetivamente explicável perante cidadãos, gestores e órgãos de controle.
\subsection{Compliance, integridade e controles institucionais}
\label{sec:org973c0bf}
A passagem do debate sobre ética e accountability para o campo do compliance e da integridade institucional representa um deslocamento analítico importante: não basta reconhecer que sistemas de inteligência artificial podem afetar direitos, introduzir vieses ou obscurecer decisões; é necessário compreender como a administração pública constrói mecanismos internos capazes de prevenir desvios, detectar falhas, responsabilizar agentes e sustentar conformidade contínua entre inovação tecnológica e deveres jurídicos e organizacionais. No setor público, compliance não pode ser entendido apenas em sentido empresarial, como aderência normativa voltada à redução de perdas reputacionais ou sancionatórias. Em contexto estatal, trata-se de um arranjo mais amplo de conformidade com princípios administrativos, deveres procedimentais, controles de legalidade, proteção do interesse público e salvaguardas de integridade na formulação, contratação, operação e supervisão de sistemas algorítmicos. A literatura revisada converge ao sugerir que a adoção de IA em governos exige institucionalização de controles desde o desenho até o uso cotidiano, especialmente porque a automação, quando incorporada aos procedimentos administrativos, pode tornar erros mais escaláveis, mais opacos e mais difíceis de corrigir \autocite{serra2024a} \autocite{parycek2023artificial} \autocite{batool2025ai}.
Esse ponto é particularmente relevante para a administração pública federal brasileira, onde a incorporação de IA ocorre em ambiente marcado por forte juridicidade, pluralidade de órgãos de controle e crescente pressão por transformação digital. Nessa moldura, a conformidade não se reduz à verificação final da legalidade de um ato, mas passa a depender de um sistema de controles institucionais distribuídos ao longo do ciclo de vida da tecnologia. Madan e Ashok mostram, em sua revisão sistemática, que a literatura sobre adoção e difusão de IA na administração pública identifica barreiras relacionadas não apenas à infraestrutura e à capacidade técnica, mas também à ausência de estruturas de governança, clareza normativa e confiança institucional \autocite{madan2022ai}. Esse diagnóstico é decisivo para a presente subseção porque evidencia que problemas de compliance não surgem apenas quando um sistema já está em operação e produz dano; eles frequentemente se originam em fases anteriores, como a definição inadequada do problema público, a seleção acrítica de fornecedores, a baixa qualidade dos dados, a falta de documentação e a inexistência de critérios formais para decidir em quais atividades a automação é admissível.
A noção de integridade, por sua vez, amplia o debate ao incluir padrões de conduta, coerência institucional e proteção contra usos indevidos da tecnologia. Em vez de restringir-se ao combate à corrupção em sentido estrito, integridade institucional envolve a capacidade da organização pública de manter consistência entre finalidades públicas, meios tecnológicos empregados e controles aplicados sobre esses meios. Serra e Machado sustentam que a literatura mais recente sobre IA no setor público já reconhece que governança adequada e considerações éticas precisam ser combinadas para orientar a implementação \autocite{serra2024a}. Contudo, quando se observa esse argumento pela lente do compliance, percebe-se uma lacuna frequente: muitos estudos defendem princípios gerais, mas dedicam menor atenção à arquitetura operacional que converte tais princípios em rotinas verificáveis. Em outras palavras, a integridade da IA pública depende menos da proclamação de valores abstratos e mais da existência de registros, trilhas de auditoria, segregação de funções, revisão humana qualificada, critérios de escalonamento de incidentes e procedimentos de resposta institucional.
A literatura sobre governança da IA reforça essa leitura ao indicar que responsible AI, ethical AI e AI governance são campos fortemente interligados, mas não equivalentes. Batool et al. observam que a governança da IA evoluiu justamente em resposta a incidentes, falhas e externalidades, de modo que estruturas de controle passaram a ser tratadas como condição de legitimidade tecnológica, e não apenas como mecanismo corretivo posterior \autocite{batool2025ai}. Essa formulação é central para o compliance público, pois desloca o foco de uma lógica reativa para uma lógica preventiva e monitorada. Em termos institucionais, isso significa reconhecer que a conformidade de sistemas algorítmicos deve ser construída ex ante, testada durante a implementação e acompanhada ex post por meio de auditoria, revisão e mecanismos de correção. A ausência de qualquer uma dessas etapas enfraquece a integridade do arranjo e aumenta o risco de que decisões automatizadas ou assistidas por IA passem a operar de modo informal, sem lastro documental suficiente para controle administrativo, judicial ou social.
No âmbito dos procedimentos administrativos, Parycek, Schmid e Novak oferecem contribuição particularmente relevante ao argumentar que a inserção de IA e automação em processos administrativos depende de avaliação rigorosa de capacidades, limitações e condições estruturais de uso \autocite{parycek2023artificial}. Os autores insistem na centralidade dos dados, da transparência e das condições jurídicas do procedimento, o que permite interpretar o compliance algorítmico como um problema de compatibilização entre desenho técnico e garantias procedimentais. Para a administração pública, isso implica que controles institucionais não podem ser concebidos apenas como anexos externos, aplicados depois que o sistema já foi comprado ou desenvolvido. Ao contrário, devem integrar o próprio procedimento administrativo de inovação: justificativa para adoção da ferramenta, análise de adequação ao caso de uso, avaliação sobre impacto em direitos, identificação de bases legais para tratamento de dados, definição dos responsáveis funcionais e critérios para revisão de outputs. Sem isso, a automação pode produzir uma espécie de “regularidade aparente”, na qual o processo parece mais eficiente, mas se torna menos controlável e menos inteligível para gestores, órgãos de controle e cidadãos.
Essa preocupação se intensifica quando se consideram as múltiplas camadas organizacionais envolvidas na IA pública. Criado, Sandoval-Almazán e Gil-Garcia propõem compreender o fenômeno em níveis micro, meso e macro, articulando atores, organizações e ambientes institucionais mais amplos \autocite{criado2024artificial}. Aplicada ao tema do compliance, essa abordagem permite perceber que controles institucionais eficazes dependem de coordenação entre dimensões distintas. No nível micro, importam competências, deveres funcionais e capacidade dos servidores de compreender limites e riscos dos sistemas. No nível meso, tornam-se centrais a governança interna, a distribuição de responsabilidades, a documentação e as rotinas de supervisão. No nível macro, ganham relevo os marcos regulatórios, as expectativas normativas e a atuação de instâncias externas de controle. A principal consequência analítica é que falhas de integridade raramente decorrem de um único erro isolado; elas emergem de desalinhamentos entre essas camadas, como quando diretrizes estratégicas são vagas, contratos são insuficientes e equipes operacionais não dispõem de formação para questionar os resultados do sistema.
A literatura sobre adoção tecnológica em governos subnacionais e locais reforça, por via comparativa, a relevância desse ponto. Yigitcanlar, Agdas e Degirmenci mostram que gestores públicos percebem a IA simultaneamente como oportunidade e como fonte de restrições associadas a capacidades organizacionais, dados, confiança e accountability \autocite{yiitcanlar2022artificial}. Embora o foco do estudo seja o governo local, sua contribuição é útil para o caso federal: a conformidade algorítmica não pode ser presumida a partir da sofisticação técnica da solução. Quanto mais complexa a ferramenta, maior a necessidade de controles claros sobre aquisição, parametrização, uso e avaliação. De modo semelhante, David et al., ao revisarem estratégias de adoção de tecnologias digitais por governos locais, indicam que a implementação bem-sucedida depende de alinhamento entre estratégia, estrutura organizacional e capacidades institucionais \autocite{david2023understanding}. Em chave de compliance, isso significa que a integridade de sistemas de IA é inseparável da qualidade da governança organizacional que os abriga.
Um dos temas mais sensíveis nessa agenda é a auditabilidade. A literatura recente não trata auditoria de IA como mero exame técnico de código ou performance, mas como prática híbrida, situada na interseção entre direito, ética, gestão e engenharia. Mökander argumenta que a auditoria de IA pode aprender com tradições anteriores de auditoria financeira, de segurança e de pesquisa social, ao mesmo tempo em que enfrenta limitações próprias de sistemas dinâmicos, opacos e dependentes de contexto \autocite{mkander2023auditing}. A relevância dessa contribuição para a administração pública é dupla. Primeiro, ela ajuda a afastar uma visão ingênua segundo a qual bastaria “abrir a caixa-preta” tecnicamente para resolver o problema da conformidade. Segundo, demonstra que auditoria útil exige definição prévia do que deve ser auditado: desempenho, fairness, aderência normativa, qualidade dos dados, robustez documental, governança contratual ou efeitos organizacionais. No setor público, essas dimensões tendem a coexistir. Assim, mecanismos de controle institucional precisam combinar auditorias técnicas com auditorias procedimentais e jurídicas, sob pena de se restringirem a verificar métricas sem avaliar se o uso concreto da IA permanece compatível com o devido processo administrativo.
Minkkinen aprofunda esse debate ao conceituar a auditoria contínua de IA e ao discutir ferramentas e frameworks destinados ao monitoramento constante de sistemas algorítmicos \autocite{minkkinen2022continuous}. Essa abordagem é especialmente pertinente para aplicações governamentais porque muitos sistemas não permanecem estáticos após sua implantação. Mudanças em bases de dados, contextos de uso e padrões comportamentais podem alterar gradualmente o desempenho e os efeitos distributivos do modelo. Para o compliance institucional, a consequência é clara: controles pontuais ou exclusivamente ex ante são insuficientes. É preciso instituir rotinas periódicas de verificação, indicadores de desvio, gatilhos para revalidação, revisão de logs e procedimentos de suspensão temporária quando forem detectadas inconsistências significativas. Em outras palavras, o controle da IA no setor público deve ser processual e contínuo, não episódico.
Esse raciocínio aproxima o tema de compliance da gestão de pessoas e das competências organizacionais. A literatura mostra que a fragilidade institucional no uso de IA decorre frequentemente de assimetrias de conhecimento entre desenvolvedores, gestores, usuários finais e órgãos de controle. Melo destaca, em revisão sistemática sobre competências em IA na gestão pública, a necessidade de capacidades técnicas, analíticas e ético-regulatórias para uso responsável dessas tecnologias \autocite{melo2026competencias}. De forma convergente, estudos sobre IA em gestão de pessoas no setor público apontam que a adoção tecnológica sem qualificação adequada amplia riscos de uso acrítico, reprodução de vieses e dependência excessiva de fornecedores \autocite{chilunjika2022artificial} \autocite{cho2023human} \autocite{barros2025inteligencia}. A implicação para o compliance é decisiva: não existe controle institucional robusto sem capacidade humana de interpretar, questionar e revisar os sistemas. Controles documentais e normativos são necessários, mas insuficientes se os agentes públicos não tiverem competências para operá-los substantivamente.
Essa dependência de capacidades também afeta a relação com o mercado e com fornecedores privados. Em muitos contextos de administração pública, soluções de IA são adquiridas de terceiros ou desenvolvidas em arranjos híbridos. Isso cria desafios de integridade ligados à opacidade contratual, à assimetria informacional e à dificuldade de exigir explicações tecnicamente inteligíveis e juridicamente relevantes. Brkan mostra, ao discutir a viabilidade legal e técnica da explicação de decisões algorítmicas, que a pretensão de total transparência encontra limites importantes, especialmente quando sistemas complexos e modelos de aprendizado de máquina estão em jogo \autocite{brkan2020legal}. Para o compliance público, isso não elimina a obrigação de explicabilidade; antes, exige refiná-la. Nem sempre será possível obter transparência total do modelo, mas é indispensável assegurar transparência suficiente sobre finalidade, dados utilizados, variáveis relevantes, limitações conhecidas, critérios de uso e possibilidades de contestação. Em sede institucional, isso demanda cláusulas contratuais adequadas, exigências de documentação e padrões mínimos de prestação de contas técnica por parte dos fornecedores.
A aplicação da IA em áreas sensíveis evidencia ainda mais essas exigências. Estudos sobre saúde pública, serviços de emprego e análise preditiva mostram que sistemas algorítmicos podem afetar acesso a políticas, triagem de demandas e priorização de beneficiários \autocite{lemes2020o} \autocite{krtner2021predictive}. Nesses casos, o déficit de compliance não é um problema meramente administrativo; ele pode converter-se em lesão material a direitos, discriminação indireta ou reforço de desigualdades preexistentes. Por isso, a literatura aponta que controles institucionais precisam ser calibrados conforme a sensibilidade do caso de uso. Quanto maior o potencial de impacto sobre direitos, maior deve ser a intensidade de documentação, validação, revisão humana e auditabilidade. Essa ideia também aparece em revisões mais amplas sobre oportunidades e desafios da IA em serviços governamentais, nas quais os ganhos de eficiência são sempre condicionados por preocupações com privacidade, justiça, confiança e coordenação institucional \autocite{alhosani2024opportunities} \autocite{pi2021machine}.
No plano comparado ibero-americano e brasileiro, a literatura aponta avanços discursivos, mas maturidade institucional desigual. Criado, ao examinar o setor público latino-americano à luz da Carta Ibero-americana de Inteligência Artificial na Administração Pública, sugere que o desafio não reside apenas em aderir a princípios, mas em convertê-los em capacidades administrativas, arranjos de coordenação e instrumentos de controle \autocite{criado2024inteligencia}. No Brasil, revisões sobre processos de governança em órgãos públicos e sobre assimetrias de maturidade em IA mostram cenário heterogêneo, com iniciativas relevantes convivendo com distintos níveis de formalização, capacidade e supervisão \autocite{goulart2025processos} \autocite{mariano2026inteligencia}. Essa heterogeneidade afeta diretamente o compliance: órgãos com maior capacidade digital podem avançar mais rapidamente, mas também correm o risco de expandir aplicações antes de consolidar salvaguardas; órgãos com menor maturidade, por sua vez, podem terceirizar excessivamente decisões técnicas e reduzir sua autonomia de controle.
Em síntese, a literatura permite sustentar que compliance, integridade e controles institucionais constituem o elo que transforma princípios de governança e ética em práticas administrativas verificáveis. A principal convergência entre os autores está em reconhecer que a IA no setor público requer mais do que inovação orientada a desempenho: requer estruturas de conformidade capazes de disciplinar escolhas tecnológicas, documentar decisões, monitorar resultados, viabilizar auditoria e preservar a responsabilidade humana e institucional sobre atos praticados com apoio algorítmico \autocite{serra2024a} \autocite{madan2022ai} \autocite{batool2025ai} \autocite{mkander2023auditing}. As divergências concentram-se menos na necessidade desses controles e mais na ênfase atribuída a diferentes instrumentos, como auditoria, transparência, capacidades organizacionais, desenho procedimental ou supervisão contínua. A principal lacuna, especialmente relevante para o recorte desta dissertação, é que ainda há limitada sistematização sobre como tais mecanismos podem ser adaptados ao contexto específico da administração pública federal brasileira, marcada por forte controle jurídico-formal, diversidade organizacional e crescente adoção de IA em funções administrativas e analíticas. Essa lacuna justifica tratar o compliance algorítmico não como dimensão acessória, mas como componente estruturante da governança pública da IA, em íntima relação com ética, accountability, gestão de riscos e supervisão humana.
\subsection{Gestão de riscos e supervisão humana}
\label{sec:org9c761f9}
Se, na subseção anterior, o problema central foi demonstrar que a conformidade institucional depende de controles distribuídos ao longo do ciclo de vida dos sistemas de IA, a discussão sobre gestão de riscos e supervisão humana aprofunda essa mesma lógica sob um prisma operacional e decisório. Em matéria de inteligência artificial aplicada ao setor público, risco não se confunde com mera possibilidade abstrata de erro tecnológico. Trata-se da probabilidade de que a automação produza efeitos incompatíveis com finalidades públicas, direitos dos administrados, racionalidade procedimental, integridade organizacional e legitimidade democrática. A literatura revisada converge em reconhecer que o setor público enfrenta uma configuração de risco distinta daquela observada em ambientes puramente mercadológicos, porque decisões estatais afetam acesso a políticas, distribuição de benefícios, exercício de poder de polícia, priorização de fiscalização e produção de evidências para decisão administrativa \autocite{madan2022ai} \autocite{criado2024artificial} \autocite{pi2021machine}. Nesses contextos, uma falha algorítmica não é apenas um defeito de desempenho: ela pode alterar padrões de tratamento estatal, consolidar assimetrias de informação e dificultar a revisão de atos que impactam diretamente a esfera jurídica dos cidadãos.
A literatura mais recente sobre governança da IA sustenta que a gestão de riscos deve ser compreendida como núcleo estruturante da governança, e não como apêndice técnico posterior à implementação. Batool et al. observam que a evolução do debate sobre AI governance ocorreu justamente porque incidentes, vieses, opacidade e danos distribuídos revelaram a insuficiência de abordagens centradas exclusivamente em eficiência e inovação \autocite{batool2025ai}. Essa constatação é particularmente relevante para a administração pública federal brasileira, onde sistemas de IA podem ser utilizados em triagem documental, classificação de demandas, apoio à fiscalização, análise de conformidade, atendimento digital e apoio a decisões complexas. Quanto maior a integração desses sistemas aos fluxos ordinários da burocracia, maior a necessidade de identificar antecipadamente riscos de base de dados, riscos de modelagem, riscos procedimentais, riscos organizacionais e riscos de uso indevido. Nesse sentido, a gestão de riscos não se limita a testar se o algoritmo “funciona”; ela exige examinar se o sistema é apropriado para aquela finalidade pública, se o contexto institucional comporta seu uso e se existem condições reais para contenção de danos e revisão de resultados \autocite{parycek2023artificial} \autocite{serra2024a}.
Parycek, Schmid e Novak fornecem contribuição importante ao mostrar que a inserção de IA em procedimentos administrativos depende de avaliação abrangente sobre capacidades, limitações e condições de enquadramento institucional \autocite{parycek2023artificial}. A ênfase dos autores na qualidade e disponibilidade dos dados, na validade das previsões e na necessidade de compatibilidade entre tecnologia e procedimento ajuda a distinguir dois planos frequentemente confundidos: o risco técnico e o risco jurídico-administrativo. Um sistema pode apresentar bom desempenho estatístico e, ainda assim, ser inadequado para um procedimento administrativo em que o dever de motivação, a possibilidade de contraditório, a rastreabilidade da decisão e a inteligibilidade para o administrado sejam elementos centrais. Em outras palavras, a boa performance preditiva não elimina o problema da adequação institucional. Essa distinção é essencial para o caso brasileiro, porque muitos usos potencialmente sedutores da IA na administração pública federal se concentram justamente em atividades nas quais a pressão por produtividade pode levar à naturalização de soluções automatizadas sem o devido exame de seus efeitos procedimentais.
Madan e Ashok, ao revisarem a literatura sobre adoção e difusão de IA na administração pública, mostram que os obstáculos mais significativos não decorrem apenas de limitações tecnológicas, mas também de lacunas de capacidades organizacionais, estruturas de governança frágeis, déficits regulatórios e resistência institucional \autocite{madan2022ai}. Lida à luz da gestão de riscos, essa formulação sugere que o risco algorítmico é inseparável do risco organizacional. Em administrações com baixa maturidade digital, pouca documentação de processos, dados heterogêneos e escassez de competências especializadas, a IA tende a ser incorporada em ambiente de controle imperfeito. O risco, portanto, não reside só no modelo, mas na combinação entre modelo, processo e organização. Essa é uma das principais convergências entre a literatura de administração pública e a literatura de governança da IA: ambas indicam que a segurança institucional da automação depende menos de soluções milagrosas e mais de capacidade estatal para delimitar usos, supervisionar resultados, registrar decisões, corrigir desvios e aprender com incidentes \autocite{criado2024artificial} \autocite{sigfrids2022how}.
A supervisão humana aparece, nesse debate, como salvaguarda decisiva, mas a literatura também alerta contra leituras simplistas do chamado \emph{human in the loop}. Não basta afirmar genericamente que haverá um humano supervisionando o sistema. É preciso perguntar quem supervisiona, em que momento, com quais informações, com qual grau de autonomia frente à recomendação algorítmica e com quais responsabilidades em caso de erro. Sigfrids et al., ao discutir propostas para fomentar o desenvolvimento ético e o uso responsável de IA nas administrações públicas, ressaltam que a governança adequada exige arranjos institucionais concretos, e não apenas declaração de princípios \autocite{sigfrids2022how}. Aplicado à supervisão humana, isso significa que a intervenção humana só funciona como mecanismo real de mitigação de riscos quando é qualificada, documentada e situada em pontos relevantes do processo decisório. Um agente público chamado a apenas “carimbar” a saída do sistema, sem tempo, sem formação e sem possibilidade material de contestá-la, não realiza supervisão substantiva; apenas transfere legitimidade formal à automação.
Esse problema se conecta ao fenômeno da automação complacente, isto é, à tendência de operadores humanos atribuírem confiança excessiva às recomendações de sistemas computacionais, especialmente quando estes são apresentados como tecnicamente sofisticados ou estatisticamente precisos. Embora a literatura do corpus nem sempre trate esse ponto com a mesma terminologia, ele emerge de forma clara nos estudos que problematizam a adoção acrítica de IA em decisões públicas e em gestão de pessoas \autocite{cho2023human} \autocite{chilunjika2022artificial} \autocite{barros2025inteligencia}. No campo de recursos humanos, por exemplo, o uso de analítica avançada e IA para recrutamento, avaliação ou alocação pode reproduzir vieses históricos e gerar discriminações difíceis de detectar se o agente público assumir que o sistema é neutro por definição. Essa percepção é reforçada por estudos sobre gestão pública e IA que assinalam o risco de se converter conhecimento administrativo tácito em confiança indevida na modelagem de dados, sem considerar que indicadores e categorias administrativas são eles próprios produtos de escolhas institucionais e históricas \autocite{chilunjika2022artificial} \autocite{cho2023human}.
A necessidade de supervisão humana também se intensifica em usos preditivos. Kärtner, ao examinar algoritmos preditivos na prestação de serviços públicos de emprego, evidencia tensões entre eficiência alocativa, categorização de usuários e riscos de tratamento desigual \autocite{krtner2021predictive}. Ainda que o contexto não seja o da administração federal brasileira em sentido estrito, a análise é altamente pertinente porque ilustra como ferramentas de predição podem deslocar a discricionariedade burocrática para camadas técnicas menos visíveis, reconfigurando critérios de priorização e intervenção. Em situações desse tipo, a supervisão humana não deve ser concebida apenas como última barreira antes da decisão final, mas como presença contínua na definição do problema, na escolha das variáveis, na interpretação dos resultados e na revisão das consequências distributivas do sistema. Sem esse acompanhamento, o aparato estatal corre o risco de institucionalizar desigualdades sob aparência de neutralidade analítica.
A literatura sobre explicabilidade e direito à explicação contribui para qualificar esse ponto. Brkan mostra que a busca por explicações de decisões algorítmicas enfrenta limites técnicos e jurídicos, especialmente em sistemas complexos e opacos \autocite{brkan2020legal}. Para a gestão de riscos, a implicação é clara: quando a explicabilidade plena é difícil ou custosa, a exigência de supervisão humana se torna ainda mais importante, mas também mais exigente. O supervisor humano não pode depender exclusivamente da inteligibilidade interna do modelo; ele precisa contar com documentação de finalidade, descrição dos dados utilizados, métricas de desempenho, limites de uso, critérios de intervenção e canais de contestação. A supervisão, portanto, deve ser apoiada por uma infraestrutura documental e organizacional que reduza a opacidade operacional do sistema, ainda que não elimine completamente a complexidade técnica \autocite{mkander2023auditing}.
É nesse ponto que o tema da auditoria se articula de modo mais direto com a gestão de riscos. Mökander argumenta que a auditoria de IA constitui campo em expansão e pode funcionar como importante mecanismo de governança, desde que se reconheçam seus limites e se evite tratá-la como solução automática para todos os problemas de legitimidade tecnológica \autocite{mkander2023auditing}. A principal contribuição dessa abordagem para a administração pública é mostrar que a gestão de riscos requer controles verificáveis, capazes de produzir evidências sobre conformidade, robustez, justiça procedimental e aderência ao desenho institucional. Já Minkkinen et al., ao conceituarem a auditoria contínua de IA, reforçam a ideia de que o controle não deve ocorrer apenas em momentos excepcionais, como após um incidente grave ou durante avaliação inicial do sistema \autocite{minkkinen2022continuous}. Em ambientes públicos dinâmicos, nos quais dados, contextos e padrões de uso mudam, o risco algorítmico também é mutável. Isso exige monitoramento contínuo, revisão periódica de desempenho, detecção de deriva do modelo e reavaliação dos impactos administrativos ao longo do tempo.
A noção de risco dinâmico é particularmente útil para evitar uma visão estática da supervisão humana. Não se trata somente de acompanhar a decisão individual produzida com apoio de IA, mas de supervisionar a trajetória institucional do sistema. Um modelo inicialmente calibrado e juridicamente aceitável pode se tornar problemático após mudança de base de dados, ampliação indevida de finalidade, transferência para outro contexto decisório ou simples transformação do comportamento dos usuários. Serra e Machado, ao sistematizarem o estado da pesquisa sobre IA no setor público, enfatizam a importância de governança adequada e considerações éticas para a implementação \autocite{serra2024a}. O passo adicional que a presente subseção propõe é interpretar essa governança adequada como capacidade institucional de operar sob incerteza, reconhecendo que o risco nunca é definitivamente eliminado. Em vez de buscar segurança total, a administração precisa construir mecanismos de detecção precoce, resposta proporcional, aprendizado organizacional e revisão do próprio desenho sociotécnico.
Essa perspectiva é compatível com as leituras de Criado, Sandoval-Almazán e Gil-Garcia, que propõem compreender a IA na administração pública a partir de níveis micro, meso e macro \autocite{criado2024artificial}. No plano micro, a supervisão humana envolve competências, julgamento e interação entre servidor e sistema. No plano meso, envolve estruturas organizacionais, processos, rotinas e governança interna. No plano macro, envolve ambiente regulatório, expectativas sociais, capacidades estatais e formas de legitimação democrática da automação. A gestão de riscos eficaz depende precisamente da articulação entre esses níveis. Concentrar-se apenas no operador humano individual é insuficiente, porque muitos riscos derivam de decisões organizacionais anteriores, como contratação inadequada, definição opaca de escopo, pressão por metas de produtividade ou ausência de padrões institucionais para revisão. Do mesmo modo, focar apenas em regulações gerais sem enfrentar rotinas concretas de uso deixa intacta a zona em que os danos efetivamente se materializam.
Os estudos sobre adoção de tecnologias digitais por governos locais e organizações públicas ajudam a reforçar esse argumento ao mostrar que escolhas tecnológicas são condicionadas por estratégias institucionais, recursos disponíveis e percepções gerenciais \autocite{yiitcanlar2022artificial} \autocite{david2023understanding}. Embora parte dessa literatura tenha foco subnacional, suas conclusões são úteis para o governo federal porque evidenciam que a disposição para adotar IA frequentemente antecede a maturidade para governá-la. Gestores podem perceber ganhos de eficiência, inovação e responsividade, mas subestimar custos de supervisão, manutenção de controles e necessidade de competências especializadas. Essa defasagem entre entusiasmo adotante e capacidade de supervisão constitui, em si mesma, um risco organizacional. Quando a instituição internaliza a tecnologia sem internalizar simultaneamente os encargos de controle que ela exige, a supervisão humana tende a se converter em ficção procedimental.
A dimensão das competências merece destaque próprio. Melo, ao revisar a literatura sobre competências em IA na gestão pública, evidencia que a capacidade estatal para lidar com essas tecnologias depende de repertórios técnicos, analíticos, éticos e gerenciais \autocite{melo2026competencias}. Isso afeta diretamente a qualidade da supervisão humana. Um supervisor sem conhecimentos mínimos sobre funcionamento do sistema, qualidade dos dados, limites estatísticos e potenciais vieses dificilmente conseguirá questionar saídas automatizadas de forma substantiva. Por outro lado, uma leitura exclusivamente técnica também é insuficiente: a supervisão em contexto público requer compreender finalidades administrativas, impactos sobre direitos, exigências procedimentais e responsabilidades institucionais. A lacuna identificada em parte da literatura está justamente em tratar competência digital como sinônimo de aptidão técnica, quando, para fins de gestão de riscos públicos, é necessário um perfil híbrido de competência sociotécnica e jurídico-organizacional \autocite{melo2026competencias} \autocite{criado2024artificial}.
No caso brasileiro, essa discussão adquire densidade adicional porque a administração pública federal opera em ambiente de alta heterogeneidade institucional. Órgãos com maior capacidade tecnológica podem desenvolver práticas mais sofisticadas de monitoramento, enquanto outros podem depender de soluções terceirizadas, infraestrutura limitada e equipes reduzidas. Mariano, ao mapear assimetrias de maturidade no uso de IA e IA generativa na administração pública, reforça a ideia de que a difusão dessas ferramentas não ocorre de forma homogênea \autocite{mariano2026inteligencia}. Sob a ótica da gestão de riscos, isso significa que não basta propor princípios universais de supervisão humana; é preciso reconhecer variações de capacidade, estágio de maturidade e complexidade das aplicações. Em contextos de baixa maturidade, usos de menor risco e maior reversibilidade podem ser mais compatíveis com a capacidade institucional disponível, ao passo que aplicações decisórias sensíveis demandam controles mais robustos, documentação intensiva e governança mais densa.
Em síntese, a literatura revisada permite sustentar que gestão de riscos e supervisão humana formam um binômio inseparável na adoção responsável de IA pela administração pública. A gestão de riscos sem supervisão humana degenera em tecnocracia procedimental incapaz de responder a contextos imprevistos, valores públicos concorrentes e efeitos distributivos não antecipados. A supervisão humana sem gestão de riscos, por sua vez, reduz-se a intervenção ad hoc, pouco estruturada e vulnerável à captura simbólica pela própria autoridade do sistema. A contribuição mais consistente do corpus analisado está em mostrar que a legitimidade da IA pública depende de converter a supervisão humana em prática institucionalizada, competente, documentada e auditável, articulada a avaliações ex ante, monitoramento contínuo, mecanismos de revisão e capacidade de correção \autocite{batool2025ai} \autocite{mkander2023auditing} \autocite{minkkinen2022continuous} \autocite{parycek2023artificial}. Para o recorte desta dissertação, isso implica reconhecer que a adoção de IA no governo público federal brasileiro só pode ser considerada compatível com o interesse público quando a busca por eficiência vier acompanhada de arranjos efetivos de contenção de risco, contestabilidade e comando humano significativo sobre sistemas que influenciam a ação estatal.
\section{METODOLOGIA}
\label{sec:org3ad4e19}
\subsection{Estratégia de pesquisa}
\label{sec:org10a6d11}
A presente dissertação adota uma estratégia de revisão de literatura de caráter sistemático, com finalidade analítica e interpretativa, orientada para mapear, comparar e integrar a produção acadêmica e técnico-institucional sobre o uso de inteligência artificial na administração pública, tendo como foco específico o governo público federal brasileiro no período de 2019 a 2026. A escolha desse desenho metodológico decorre da própria natureza do objeto. O uso de IA no setor público não constitui um campo estabilizado, com fronteiras conceituais totalmente consolidadas, mas um domínio em rápida expansão, marcado pela sobreposição entre debates de transformação digital, automação administrativa, governança algorítmica, ética aplicada, conformidade regulatória e gestão de riscos. Em razão disso, uma estratégia de revisão sistemática mostra-se mais adequada do que uma revisão narrativa ampla, pois permite explicitar critérios de busca, seleção e interpretação, reduzir arbitrariedades na composição do corpus e produzir uma síntese crítica capaz de distinguir tendências robustas de formulações ainda exploratórias \autocite{prisma2020statement} \autocite{madan2022ai}.
A revisão foi concebida não como simples inventário bibliográfico, mas como dispositivo de construção de um quadro analítico coerente com o problema de pesquisa. Em vez de tomar a literatura apenas como repositório de exemplos de uso de IA, a estratégia adotada procurou compreender como diferentes tradições de pesquisa definem o fenômeno, quais dimensões priorizam e quais lacunas permanecem abertas quando o debate é transposto para a administração pública federal brasileira. Essa orientação é importante porque a literatura internacional frequentemente enfatiza difusão tecnológica, capacidades organizacionais e inovação administrativa, enquanto os estudos mais diretamente relacionados ao contexto público latino-americano e brasileiro tendem a enfatizar enquadramentos normativos, assimetrias institucionais, maturidade digital desigual e desafios de responsabilização \autocite{criado2024artificial} \autocite{criado2024inteligencia} \autocite{serra2024a}. Assim, a estratégia de pesquisa buscou articular duas escalas simultâneas: uma escala mais ampla, voltada à compreensão do estado da arte sobre IA e administração pública, e uma escala aplicada, dirigida à extração de categorias relevantes para o recorte federal brasileiro.
Nesse sentido, o desenho da revisão combinou função mapeadora e função explanatória. A função mapeadora consistiu em identificar os principais eixos temáticos recorrentes na literatura: adoção e difusão da IA, aplicações em processos administrativos, governança institucional, ética, auditoria, compliance, capacidades organizacionais, competências estatais e gestão de riscos. A função explanatória, por sua vez, consistiu em examinar como esses eixos se relacionam entre si e quais pressupostos sustentam as recomendações de implementação responsável da IA no setor público. Madan e Ashok demonstram que a adoção de IA na administração pública depende de fatores técnicos, organizacionais e institucionais interdependentes, o que justificou tratar a literatura não por compartimentos isolados, mas por nexos entre difusão, governança e capacidade estatal \autocite{madan2022ai}. De forma convergente, Criado, Sandoval-Almazán e Gil-Garcia propõem uma leitura em níveis micro, meso e macro, útil para organizar a revisão em torno de atores, arranjos institucionais e contextos político-regulatórios \autocite{criado2024artificial}.
A composição do corpus privilegiou estudos que oferecessem densidade conceitual e relevância substantiva para o recorte da dissertação. Por essa razão, foram incluídos tanto trabalhos de revisão sistemática e integrativa quanto artigos analíticos voltados a aplicações específicas, desde que contribuíssem para a compreensão do problema central. Revisões como as de Serra \autocite{serra2024a}, Batool et al. \autocite{batool2025ai}, Lyrio \autocite{lyrio2025emerging}, Mariano \autocite{mariano2026inteligencia} e Goulart \autocite{goulart2025processos} são particularmente importantes porque sintetizam tendências, recorrências e lacunas do campo, permitindo identificar padrões de debate. Já estudos como os de Parycek, Schmid e Novak \autocite{parycek2023artificial}, Mökander \autocite{mkander2023auditing}, Minkkinen et al. \autocite{minkkinen2022continuous}, Brkan \autocite{brkan2020legal} e Kürtner \autocite{krtner2021predictive} foram mobilizados porque aprofundam dimensões críticas da implementação, como automação procedimental, auditabilidade, monitoramento contínuo, explicabilidade e riscos associados a decisões preditivas em contextos públicos.
A racionalidade dessa estratégia consistiu em evitar dois desvios metodológicos frequentes. O primeiro seria restringir a revisão a textos estritamente brasileiros, o que produziria uma base empírica mais aderente ao caso nacional, porém insuficiente para captar a evolução conceitual e regulatória internacional que hoje estrutura o debate sobre IA responsável no setor público. O segundo seria diluir o objeto em uma literatura genérica sobre governo digital ou inovação pública, perdendo precisamente as especificidades da inteligência artificial enquanto tecnologia sociotécnica capaz de alterar rotinas decisórias, critérios de elegibilidade, procedimentos de fiscalização e formas de prestação de serviços. Para contornar esses desvios, a revisão operou por pertinência temática forte: foram valorizados estudos em que a IA aparecesse ligada, de maneira substantiva, a governança, ética, compliance, auditoria, gestão de riscos, automação administrativa ou capacidades estatais, e não apenas como marcador genérico de modernização \autocite{alhosani2024opportunities} \autocite{david2023understanding} \autocite{pi2021machine}.
Outro aspecto central da estratégia de pesquisa foi a decisão de incorporar deliberadamente literatura de diferentes níveis territoriais e funcionais da administração pública, inclusive estudos sobre governos locais, gestão de pessoas, saúde e serviços públicos, desde que oferecessem mecanismos analíticos transferíveis ao plano federal. Isso se justifica porque muitos problemas associados à IA no setor público não são exclusivos de um nível de governo, mas reaparecem em formatos variados: baixa qualidade de dados, opacidade decisória, déficit de competências, fragilidade de supervisão humana, dependência de fornecedores, assimetrias de poder informacional e dificuldades de responsabilização \autocite{yiitcanlar2022artificial} \autocite{chilunjika2022artificial} \autocite{cho2023human} \autocite{lemes2020o}. O interesse da dissertação, portanto, não está em transpor automaticamente evidências de contextos distintos para a realidade federal brasileira, mas em utilizá-las comparativamente para identificar categorias, tensões e condições de possibilidade relevantes para esse recorte.
Por fim, a estratégia de pesquisa foi estruturada para sustentar uma leitura crítica da literatura, e não uma simples acumulação de posições autorais. Isso significa que o corpus foi tratado como campo de debate, no qual coexistem promessas de eficiência, argumentos de prudência institucional e propostas de regulação e auditoria. Estudos voltados à governança ética da IA enfatizam a necessidade de princípios, estruturas decisórias e mecanismos de coordenação \autocite{sigfrids2022how} \autocite{batool2025ai}, enquanto a literatura sobre implementação evidencia obstáculos concretos ligados a capacidades organizacionais, resistência burocrática e maturidade institucional \autocite{sharma2021implementing} \autocite{melo2026competencias}. Ao integrar essas vertentes, a revisão busca construir base metodológica capaz de responder ao objetivo do trabalho: compreender não apenas onde a IA vem sendo empregada na administração pública, mas sob quais condições sua adoção pode ser considerada institucionalmente legítima, eticamente defensável, auditável e compatível com exigências de compliance e gestão de riscos no âmbito do governo público federal brasileiro.
\subsection{Critérios de seleção do corpus}
\label{sec:org171c69e}
A definição dos critérios de seleção do corpus foi orientada por um princípio de dupla aderência: aderência temática ao objeto da dissertação e aderência analítica às categorias centrais de investigação, quais sejam governança, ética, compliance e gestão de riscos no uso de inteligência artificial pela administração pública. Isso significou que a seleção não se limitou à identificação de textos que mencionassem IA e setor público de modo genérico, mas exigiu que os estudos efetivamente contribuíssem para compreender como sistemas algorítmicos são introduzidos, governados, supervisionados e problematizados em contextos administrativos. Tal delimitação foi necessária porque a literatura sobre transformação digital no Estado é muito mais ampla do que a literatura especificamente voltada à IA, e a inclusão indiscriminada de trabalhos sobre governo eletrônico, inovação pública ou digitalização administrativa diluiria o foco substantivo da revisão \autocite{madan2022ai} \autocite{david2023understanding}.
Nesse desenho, o corpus foi composto por estudos capazes de oferecer pelo menos uma das seguintes contribuições: enquadramento conceitual sobre adoção e difusão de IA na administração pública; análise de aplicações concretas em funções estatais; discussão de arranjos de governança e capacidades institucionais; problematização de riscos éticos, jurídicos e organizacionais; ou proposição de mecanismos de auditoria, supervisão e responsabilização. Esse critério permitiu integrar revisões amplas e estudos de aprofundamento sem perder coerência interna. Revisões sistemáticas e integrativas foram particularmente valorizadas porque ajudam a reconstruir o estado da arte, identificar recorrências e distinguir agendas emergentes de consensos já relativamente estabilizados. É o caso de trabalhos que sintetizam o campo da IA na administração pública, da governança de IA e dos processos de institucionalização em órgãos públicos \autocite{serra2024a} \autocite{madan2022ai} \autocite{batool2025ai} \autocite{goulart2025processos}.
Ao mesmo tempo, o corpus não foi restrito a estudos de revisão. Foram incluídos textos analíticos sobre procedimentos administrativos automatizados, auditabilidade de sistemas de IA, monitoramento contínuo, explicabilidade e uso preditivo em serviços públicos, porque tais trabalhos fornecem granularidade analítica indispensável para interpretar o caso brasileiro. A literatura de revisão mostra tendências gerais, mas nem sempre esclarece como se materializam os riscos de opacidade, discricionariedade automatizada, dependência de dados ou fragilidade de controles em situações administrativas concretas. Por isso, estudos como os de Parycek, Schmid e Novak sobre automação em procedimentos administrativos, os de Mökander sobre auditoria de IA, os de Minkkinen et al. sobre auditoria contínua e os de Brkan sobre explicação de decisões algorítmicas foram selecionados como peças de aprofundamento teórico-metodológico \autocite{parycek2023artificial} \autocite{mkander2023auditing} \autocite{minkkinen2022continuous} \autocite{brkan2020legal}.
Outro critério central foi a pertinência ao setor público, ainda que nem todos os estudos incidissem diretamente sobre a administração pública federal brasileira. Considerou-se que, em um campo ainda em consolidação, a compreensão do objeto exige uma estratégia de corpus escalonado. No primeiro nível, foram priorizados trabalhos diretamente voltados à administração pública e ao setor público. No segundo, foram admitidos estudos sobre governos subnacionais, comparações internacionais e contextos administrativos análogos, desde que produzissem inferências úteis sobre capacidades organizacionais, governança institucional e efeitos da IA sobre funções públicas. Essa opção é justificada pela própria literatura, que aponta forte heterogeneidade entre níveis de governo, maturidades digitais e capacidades regulatórias, sem que isso elimine a utilidade comparativa dos achados \autocite{criado2024artificial} \autocite{yiitcanlar2022artificial} \autocite{criado2024inteligencia}. Assim, pesquisas sobre governos locais, gestão de pessoas no setor público ou estratégias de adoção tecnológica foram incluídas quando contribuíam para esclarecer desafios que também se colocam no plano federal, como competências, resistência organizacional, infraestrutura de dados e desenho de salvaguardas \autocite{chilunjika2022artificial} \autocite{cho2023human} \autocite{melo2026competencias}.
Em contrapartida, foram excluídos textos cuja relação com o problema de pesquisa fosse apenas periférica. Não integraram o núcleo do corpus estudos sobre digitalização estatal sem referência clara a inteligência artificial; trabalhos centrados exclusivamente em desempenho técnico de modelos computacionais, sem discussão administrativa, institucional ou normativa; e textos excessivamente prospectivos, voltados à celebração da inovação, mas sem densidade analítica sobre governança e risco. A exclusão desses materiais não decorreu de juízo de irrelevância absoluta, mas da necessidade de preservar a inteligibilidade do corpus diante do objetivo da dissertação. A literatura recente adverte, inclusive, que o entusiasmo com benefícios de eficiência, personalização e automação tende a obscurecer assimetrias de poder, falhas de accountability e limitações institucionais quando a seleção documental não distingue inovação digital ampla de IA aplicada ao governo \autocite{alhosani2024opportunities} \autocite{pi2021machine} \autocite{serra2024a}.
Também foi adotado um critério de equilíbrio entre produção internacional e produção em português, especialmente brasileira e ibero-americana. Esse equilíbrio foi importante porque a literatura internacional oferece sofisticação conceitual sobre difusão, governança multinível e mecanismos de supervisão, enquanto a literatura mais próxima do contexto brasileiro e latino-americano ilumina problemas de transplantação institucional, assimetrias de maturidade e lacunas de implementação em ambientes administrativos marcados por desigualdade de capacidades \autocite{criado2024artificial} \autocite{criado2024inteligencia} \autocite{mariano2026inteligencia}. Do lado brasileiro, estudos sobre inserção da IA na administração pública, governança em órgãos públicos e usos setoriais, como na saúde, foram considerados estratégicos por aproximarem a discussão das condições concretas da burocracia nacional e da administração federal \autocite{serra2024a} \autocite{goulart2025processos} \autocite{lemes2020o}.
Por fim, a seleção do corpus observou um critério de utilidade analítica para a síntese final da dissertação. Não bastava que o texto fosse pertinente ao tema; era necessário que permitisse comparação, extração de categorias e articulação argumentativa. Em termos práticos, isso significou privilegiar documentos que explicitassem variáveis, fatores condicionantes, dilemas normativos ou mecanismos institucionais, em vez de meras descrições episódicas. Foi esse critério que justificou a incorporação de estudos sobre implementação, adoção e barreiras organizacionais, mesmo quando produzidos fora do contexto brasileiro estrito, pois eles ajudam a explicar por que a IA no setor público depende simultaneamente de capacidade técnica, desenho regulatório, governança de dados, controles de conformidade e supervisão humana \autocite{sharma2021implementing} \autocite{sigfrids2022how} \autocite{krtner2021predictive}. Desse modo, o corpus foi intencionalmente delimitado para sustentar uma revisão que não apenas descreva aplicações de IA no governo, mas examine criticamente as condições sob as quais tais aplicações podem ser consideradas legítimas, auditáveis e compatíveis com o interesse público.
\subsection{Bases, período e idiomas}
\label{sec:org621d95e}
A definição das bases de busca, do recorte temporal e dos idiomas da revisão foi concebida como parte substantiva do desenho metodológico, e não como uma etapa meramente operacional. Em um campo como o da inteligência artificial na administração pública, a escolha dessas dimensões condiciona diretamente o tipo de evidência capturado, o equilíbrio entre literatura conceitual e aplicada, a visibilidade de experiências nacionais e a capacidade de reconstruir tendências recentes sem perder referências formadoras do debate. Por essa razão, a revisão combinou bases de amplo alcance internacional com atenção específica à literatura em português, de modo a contemplar simultaneamente a produção internacional sobre adoção, governança e auditoria de IA e os estudos mais próximos do contexto administrativo brasileiro \autocite{madan2022ai} \autocite{serra2024a} \autocite{criado2024artificial}.
No plano das bases, o desenho privilegiou repositórios e indexadores com forte cobertura em administração pública, políticas públicas, governo digital, estudos organizacionais e discussões interdisciplinares sobre regulação tecnológica. Essa opção se justifica porque a literatura pertinente ao objeto não está concentrada em um único campo disciplinar. Parte relevante dos estudos aparece em periódicos de administração pública e government studies; outra parte, em revistas de ética aplicada, direito digital, sistemas de informação, tecnologia e sociedade ou governança algorítmica. Revisões recentes mostram precisamente essa dispersão temática e institucional do campo, indicando que a análise da IA no setor público exige estratégia de busca capaz de atravessar fronteiras disciplinares \autocite{madan2022ai} \autocite{batool2025ai}. Assim, a seleção das bases buscou reduzir dois riscos metodológicos opostos: de um lado, a subcobertura da literatura internacional de alto impacto; de outro, a invisibilização de estudos em português e de contribuições mais diretamente conectadas ao caso brasileiro.
Essa racionalidade é particularmente importante porque a produção sobre IA no setor público apresenta assimetrias de circulação. Estudos em língua inglesa tendem a concentrar formulações gerais sobre difusão tecnológica, modelos de adoção, capacidades institucionais, governança e accountability, como se observa em revisões amplas e em análises sobre atores, níveis e agendas de política pública \autocite{madan2022ai} \autocite{criado2024artificial}. Já a literatura em português, embora menos volumosa, oferece maior proximidade com problemas concretos da administração pública brasileira, como o uso da IA em saúde, a institucionalização de processos de governança em órgãos públicos, a inserção da tecnologia no aparato estatal e as assimetrias de maturidade no setor público nacional \autocite{lemes2020o} \autocite{serra2024a} \autocite{goulart2025processos} \autocite{mariano2026inteligencia}. Em termos metodológicos, portanto, as bases não foram tratadas como um funil neutro, mas como instrumentos de composição deliberada de um corpus capaz de articular generalização analítica e aderência contextual.
O recorte temporal adotado para a dissertação concentra-se no período de 2019 a 2026, em consonância com o objeto empírico delimitado para a administração pública federal brasileira. Esse intervalo cobre a fase em que a IA deixa de figurar predominantemente como promessa abstrata de inovação e passa a ser discutida de forma mais sistemática como tecnologia incorporada a fluxos decisórios, triagens documentais, serviços digitais, fiscalização, automação administrativa e apoio à gestão. Também é nesse período que se intensificam as preocupações com explicabilidade, viés, qualidade de dados, supervisão humana, auditoria e estruturas de governança, elementos centrais para esta pesquisa \autocite{parycek2023artificial} \autocite{mkander2023auditing} \autocite{minkkinen2022continuous}.
A escolha de 2019 como marco inicial não implica desconsiderar antecedentes teóricos relevantes, mas priorizar a fase em que o debate ganha densidade institucional e passa a dialogar mais diretamente com implementação pública, regulação e risco. Trabalhos imediatamente anteriores ao recorte principal foram mobilizados apenas quando indispensáveis para esclarecer fundamentos do problema, como ocorre com a discussão sobre a viabilidade jurídica e técnica da explicação de decisões algorítmicas, relevante para compreender os limites da transparência em sistemas complexos \autocite{brkan2020legal}. Em outras palavras, o horizonte principal da busca foi contemporâneo e orientado ao uso efetivo da IA no Estado, mas sem incorrer em presentismo metodológico. Essa combinação permitiu captar tanto a aceleração recente da agenda quanto seus pressupostos normativos e técnicos.
O recorte até 2026, por sua vez, responde à necessidade de incorporar a produção mais recente sobre governança, competências institucionais e IA generativa na administração pública, temas que alteraram significativamente a literatura nos últimos anos. Estudos recentes mostram que a agenda deixou de se concentrar apenas na adoção de ferramentas e passou a enfatizar maturidade organizacional, competências do setor público, mecanismos de governança responsiva e desafios de integração entre inovação e controle \autocite{melo2026competencias} \autocite{mariano2026inteligencia} \autocite{lyrio2025emerging}. A inclusão desse horizonte é especialmente relevante para uma dissertação voltada ao governo federal brasileiro, porque o debate sobre uso responsável de IA já não pode ser examinado apenas sob a ótica da automação tradicional; ele exige considerar também a expansão de aplicações generativas, os novos riscos de opacidade e a pressão por capacidades institucionais mais sofisticadas.
Quanto aos idiomas, a revisão foi delimitada ao português e ao inglês. A opção pelo inglês é metodologicamente incontornável, dada sua centralidade na literatura internacional sobre administração pública digital, governança algorítmica e auditoria de IA. É nessa língua que se encontram formulações de maior difusão sobre padrões de adoção no setor público, arranjos de governança ética, estruturas de auditoria e condicionantes organizacionais da implementação \autocite{madan2022ai} \autocite{sigfrids2022how} \autocite{mkander2023auditing}. Já a inclusão do português decorre da necessidade de capturar trabalhos voltados ao contexto brasileiro, inclusive revisões e análises que examinam aplicações setoriais, governança em órgãos públicos, competências e desafios específicos da administração nacional \autocite{serra2024a} \autocite{goulart2025processos} \autocite{barros2025inteligencia}.
A exclusão sistemática de outros idiomas não significa supor irrelevância da produção em espanhol, francês ou alemão, mas reconhecer um limite pragmático e analítico do desenho. Ainda assim, o corpus incorporou indiretamente contribuições de contextos não anglófonos quando elas estavam disponíveis nas bases selecionadas ou dialogavam com debates comparados relevantes, como no caso de análises sobre o setor público latino-americano e sobre procedimentos administrativos automatizados em tradições jurídico-administrativas europeias \autocite{criado2024inteligencia} \autocite{parycek2023artificial}. Com isso, o recorte linguístico preserva viabilidade e consistência sem fechar completamente a dimensão comparativa.
Em síntese, a combinação entre bases amplas, foco temporal em 2019-2026 e bilinguismo português-inglês foi estruturada para produzir um corpus metodologicamente robusto, comparável e aderente ao problema de pesquisa. Esse desenho permitiu reunir literatura de revisão, estudos conceituais e análises aplicadas em torno de uma questão comum: como a adoção de IA no setor público pode ser compreendida não apenas como inovação tecnológica, mas como transformação institucional sujeita a exigências de governança, ética, compliance e gestão de riscos \autocite{batool2025ai} \autocite{criado2024artificial} \autocite{alhosani2024opportunities}.
\subsection{Procedimentos de análise}
\label{sec:orgdbed2a6}
Os procedimentos de análise foram estruturados para que a revisão de literatura não se limitasse à descrição agregada de achados, mas produzisse uma interpretação comparativa e teoricamente orientada sobre o uso da inteligência artificial no governo público federal, especialmente sob os eixos de governança, ética, compliance e gestão de riscos. Em vez de tratar o corpus como um conjunto homogêneo de textos, a análise partiu do reconhecimento de que a literatura sobre IA no setor público é fragmentada entre abordagens mais tecnicistas, trabalhos voltados à inovação administrativa, estudos normativos sobre accountability e transparência, além de revisões sistemáticas e integrativas que mapeiam tendências do campo \autocite{madan2022ai} \autocite{serra2024a} \autocite{criado2024artificial}. Por essa razão, o desenho analítico adotou uma lógica de síntese qualitativa orientada por categorias, com inspiração em revisão sistemática e integrativa, mas com ênfase interpretativa voltada à construção de nexos entre aplicações, capacidades institucionais e salvaguardas necessárias à implementação responsável.
Após a consolidação do corpus, procedeu-se à leitura integral dos textos selecionados, com registro analítico padronizado de informações referentes a objeto, contexto institucional, tipo de aplicação de IA, nível governamental analisado, riscos identificados, mecanismos de governança propostos e limitações reconhecidas pelos próprios autores. Esse procedimento buscou preservar comparabilidade entre trabalhos muito distintos entre si. Revisões amplas indicam que a adoção de IA na administração pública envolve desde automação de rotinas e triagem documental até analytics preditivo, apoio à decisão e gestão de pessoas, o que exige cuidado para não reduzir o fenômeno a uma única modalidade tecnológica \autocite{madan2022ai} \autocite{pi2021machine} \autocite{chilunjika2022artificial}. No mesmo sentido, estudos sobre procedimentos administrativos automatizados mostram que a capacidade de generalização dos resultados depende de distinguir contextos em que a IA atua como ferramenta de apoio, contextos de automação parcial e situações em que algoritmos passam a influenciar mais diretamente o conteúdo da ação administrativa \autocite{parycek2023artificial}. Assim, a leitura analítica foi conduzida com atenção à função desempenhada pela IA em cada estudo, e não apenas à presença genérica da tecnologia.
A codificação do material foi organizada em duas etapas complementares. Na primeira, empregou-se uma codificação aberta, voltada a identificar recorrências empíricas e conceituais emergentes do corpus. Nessa fase, apareceram com frequência temas como eficiência administrativa, melhoria de serviços, qualidade e governança de dados, viés algorítmico, transparência, supervisão humana, responsabilização institucional, auditoria, competências organizacionais e assimetrias de maturidade. Na segunda etapa, esses elementos foram reorganizados em eixos analíticos mais estáveis, compatíveis com o objetivo da dissertação: a) usos e aplicações de IA na administração pública; b) arranjos de governança e coordenação institucional; c) dilemas éticos e salvaguardas normativas; d) mecanismos de compliance, controle e auditabilidade; e e) práticas de gestão de riscos ao longo do ciclo de vida dos sistemas. Essa estratégia dialoga com revisões que apontam a necessidade de integrar perspectivas micro, meso e macro para compreender a IA no setor público, evitando tanto o reducionismo organizacional quanto abstrações excessivamente normativas \autocite{criado2024artificial}.
No eixo das aplicações, a análise comparou como a literatura descreve diferentes campos de uso, tais como saúde pública, recursos humanos, serviços digitais, fiscalização e procedimentos administrativos. Estudos sobre a administração pública brasileira mostram, por exemplo, que a IA já é discutida em áreas concretas como a saúde, nas quais o debate sobre inovação convive com preocupações relativas a segurança, responsabilidade e proteção de direitos \autocite{lemes2020o}. Em trabalhos internacionais, observa-se padrão semelhante: benefícios esperados em termos de celeridade, capacidade analítica e personalização de serviços são reiteradamente acompanhados de advertências sobre qualidade dos dados, exclusão indevida e dependência de capacidades institucionais ainda insuficientes \autocite{alhosani2024opportunities} \autocite{yiitcanlar2022artificial}. Essa comparação foi central para distinguir promessas de desempenho de condições efetivas de implementação.
No eixo da governança, os textos foram analisados segundo o tipo de arranjo institucional que pressupõem ou recomendam. A literatura recente tem insistido que governança de IA não se reduz à existência de princípios abstratos; ela requer definição de papéis, processos decisórios, instâncias de supervisão, critérios para aquisição ou desenvolvimento de sistemas e mecanismos de revisão contínua \autocite{batool2025ai} \autocite{sigfrids2022how}. Em estudos focados no setor público e no contexto latino-americano ou brasileiro, essa questão aparece ligada à necessidade de traduzir princípios em rotinas administrativas, padrões de controle e coordenação interorganizacional \autocite{criado2024inteligencia} \autocite{goulart2025processos} \autocite{serra2024a}. Por isso, a análise procurou identificar, em cada texto, se a governança era tratada como agenda estratégica, como arquitetura de conformidade, como arranjo sociotécnico ou como mecanismo de mitigação de riscos, comparando essas formulações e suas implicações práticas.
No eixo ético-jurídico, o corpus foi examinado a partir de problemas recorrentes de explicabilidade, discriminação, opacidade e contestabilidade. Em vez de assumir que a transparência resolve por si só tais problemas, a análise considerou trabalhos que demonstram os limites concretos da exigência de explicação em sistemas complexos, especialmente quando modelos sofisticados dificultam a inteligibilidade das decisões algorítmicas \autocite{brkan2020legal}. Também foram comparadas contribuições que defendem supervisão humana significativa com aquelas que alertam para o risco de uma supervisão meramente formal, incapaz de corrigir vieses estruturais ou de enfrentar assimetrias entre fornecedores, burocracias e cidadãos \autocite{parycek2023artificial} \autocite{sigfrids2022how}. Esse tratamento permitiu compreender a ética não como apêndice discursivo, mas como dimensão operacional da governança pública de IA.
Quanto ao compliance e à gestão de riscos, a análise privilegiou a identificação de mecanismos concretos de monitoramento, auditoria e controle. A literatura especializada em auditoria de IA mostra que avaliações ex ante são insuficientes quando desacompanhadas de monitoramento contínuo, documentação robusta, testes de desempenho, revisão de impactos e responsabilização distribuída ao longo do uso do sistema \autocite{mkander2023auditing} \autocite{minkkinen2022continuous}. Tais elementos foram cotejados com estudos sobre adoção da IA no setor público que destacam barreiras organizacionais, escassez de competências e fragilidade de processos internos como fatores que ampliam riscos de implementação \autocite{sharma2021implementing} \autocite{melo2026competencias}. Com isso, o procedimento analítico buscou demonstrar que risco, nesta dissertação, não é entendido apenas como falha técnica, mas como combinação entre desenho algorítmico, qualidade institucional, capacidade de controle e consequências administrativas sobre direitos, acesso e legitimidade.
Por fim, os achados foram sintetizados por convergência e contraste. Em vez de somar evidências de forma indiferenciada, procurou-se reconstruir padrões: onde há consenso robusto, como na centralidade da governança de dados e da supervisão institucional; onde predominam controvérsias, como nos alcances reais da explicabilidade; e onde se observam lacunas, especialmente na literatura aplicada ao governo federal brasileiro. Esse movimento de síntese foi orientado pela compreensão de que revisões de literatura em campos emergentes devem não apenas mapear temas, mas também ordenar o debate e explicitar assimetrias entre formulações normativas e evidências empíricas \autocite{madan2022ai} \autocite{lyrio2025emerging} \autocite{mariano2026inteligencia}. Desse modo, os procedimentos de análise permitiram construir uma leitura substantiva do corpus, apta a sustentar, nas seções seguintes, a discussão sobre usos, modelos de governança, dilemas éticos, exigências de compliance e práticas de gestão de riscos associadas à IA na administração pública federal brasileira.
\subsection{Limitações metodológicas}
\label{sec:org2d79a4b}
Apesar do esforço de sistematização adotado nesta revisão, o desenho metodológico apresenta limitações que precisam ser explicitadas para delimitar o alcance dos achados e evitar inferências excessivas. A primeira delas decorre da própria natureza do estudo: trata-se de uma revisão de literatura, e não de uma investigação empírica direta sobre órgãos, sistemas ou processos específicos da administração pública federal brasileira. Isso significa que as conclusões produzidas dependem da qualidade, da abrangência e do grau de maturidade analítica da produção disponível, o que impõe mediações importantes entre o fenômeno observado na prática administrativa e sua representação nos textos selecionados. Em um campo ainda marcado por aceleração tecnológica, heterogeneidade conceitual e assimetria entre inovação aplicada e reflexão acadêmica, a literatura nem sempre acompanha com a mesma velocidade a difusão efetiva dos usos de IA no setor público \autocite{madan2022ai} \autocite{criado2024artificial}.
Há, além disso, uma limitação substantiva associada ao recorte temático e institucional adotado. O foco desta dissertação recai sobre o governo público federal brasileiro entre 2019 e 2026, com ênfase em governança, ética, compliance e gestão de riscos. Entretanto, parte relevante do corpus é composta por revisões e estudos produzidos em contextos internacionais, comparados ou subnacionais, incluindo governos locais, administrações europeias, experiências latino-americanas e aplicações setoriais específicas \autocite{yiitcanlar2022artificial} \autocite{david2023understanding} \autocite{criado2024inteligencia}. A utilização desse material foi metodologicamente necessária porque a literatura estritamente centrada na administração pública federal brasileira ainda é relativamente dispersa e, em vários temas, insuficiente para sustentar análise exclusiva. O custo analítico dessa opção é que parte das categorias interpretativas deriva de contextos jurídico-institucionais distintos, com capacidades estatais, arranjos regulatórios, culturas organizacionais e graus de digitalização diferentes. Assim, as aproximações realizadas ao longo da revisão têm valor heurístico e comparativo, mas não equivalem automaticamente a generalizações válidas para todos os órgãos federais brasileiros.
Outra limitação importante reside na composição do corpus. Embora o procedimento de seleção tenha buscado equilíbrio entre aderência temática e diversidade analítica, a literatura sobre IA no setor público apresenta forte concentração em textos de caráter exploratório, normativo ou prospectivo. Revisões recentes mostram que o campo ainda é relativamente dominado por mapeamentos, agendas de pesquisa e discussões sobre potenciais, desafios e princípios, com menor presença de avaliações longitudinais, estudos de implementação em profundidade e evidências robustas sobre resultados organizacionais concretos \autocite{madan2022ai} \autocite{serra2024a} \autocite{alhosani2024opportunities}. Isso produz uma assimetria metodológica no próprio material analisado: há abundância de argumentos sobre o que a governança de IA deveria ser, mas menor densidade de evidências sobre como ela efetivamente opera em burocracias públicas concretas. Em consequência, esta subseção reconhece que a revisão tende a identificar com mais clareza repertórios normativos, riscos antecipados e modelos desejáveis do que comprovar, em igual medida, sua implementação institucional efetiva.
Também deve ser reconhecido que a delimitação por idiomas, ainda que adequada ao objetivo do estudo, pode ter excluído contribuições relevantes publicadas em outras línguas. O predomínio de trabalhos em inglês favorece o acesso a debates internacionais consolidados, enquanto o uso do português permite captar melhor a produção voltada ao contexto brasileiro. Ainda assim, discussões importantes sobre IA na administração pública em outros sistemas jurídicos e administrativos podem não ter sido incorporadas. Isso é particularmente relevante em um campo em que experiências regulatórias e institucionais europeias, latino-americanas e asiáticas evoluem de forma desigual e nem sempre circulam plenamente nos mesmos canais de indexação \autocite{batool2025ai} \autocite{sigfrids2022how}.
Uma limitação adicional decorre das bases e estratégias de busca empregadas em revisões sistemáticas e integrativas. Mesmo quando orientadas por protocolos explícitos, como recomenda a literatura metodológica \autocite{prisma2020statement}, tais revisões dependem de descritores, operadores, indexação e metadados nem sempre homogêneos. No caso da IA, esse problema é agravado pela instabilidade terminológica do campo. Muitos estudos tratam conjuntamente IA, automação, machine learning, analytics, sistemas algorítmicos, tecnologias emergentes e transformação digital, nem sempre distinguindo com precisão suas fronteiras conceituais \autocite{pi2021machine} \autocite{lyrio2025emerging}. Em sentido oposto, trabalhos empiricamente relevantes podem descrever usos típicos de IA sem empregar a terminologia que orientou a busca. Assim, mesmo com triagem criteriosa, permanece a possibilidade de subinclusão de estudos aderentes ao fenômeno e, simultaneamente, de inclusão de textos cuja relação com IA seja mais indireta ou ampla do que o núcleo analítico central da dissertação.
No plano conceitual, a própria polissemia da noção de governança de IA introduz limites à comparabilidade. A literatura revisada utiliza o termo para designar, em diferentes casos, princípios éticos, estruturas organizacionais, instrumentos regulatórios, mecanismos de accountability, arranjos de coordenação interinstitucional, práticas de auditoria e procedimentos de gestão de riscos \autocite{batool2025ai} \autocite{sigfrids2022how} \autocite{goulart2025processos}. Esse alargamento semântico dificulta a comparação direta entre estudos e pode produzir aparente convergência onde, na verdade, há objetos analíticos distintos. Problema semelhante ocorre com categorias como transparência, explicabilidade, supervisão humana e responsabilização. A discussão jurídica e técnica demonstra que tais noções possuem diferentes níveis de exigência e viabilidade prática, especialmente em sistemas complexos ou de baixa interpretabilidade \autocite{brkan2020legal} \autocite{parycek2023artificial}. Portanto, os resultados desta revisão devem ser lidos como síntese interpretativa de um campo conceitualmente em disputa, e não como consolidação terminológica definitiva.
A análise também encontra limites quando se volta para compliance, auditoria e gestão contínua de riscos. Embora o corpus contenha contribuições relevantes sobre auditoria de IA e monitoramento ao longo do ciclo de vida dos sistemas, trata-se de um domínio ainda em construção, com baixa padronização e forte distância entre formulações normativas e rotinas organizacionais efetivas \autocite{mkander2023auditing} \autocite{minkkinen2022continuous}. Em especial no setor público, os mecanismos concretos de auditoria algorítmica, documentação técnica, rastreabilidade decisória e revisão independente ainda aparecem mais como exigência emergente do que como prática amplamente consolidada. Isso significa que a revisão consegue mapear repertórios de controle e salvaguarda com razoável consistência, mas enfrenta limitações para afirmar sua difusão real no governo federal brasileiro sem apoio de evidência empírica complementar.
Por fim, há um limite temporal inerente ao objeto. O período de 2019 a 2026 é particularmente dinâmico, marcado pela rápida expansão de sistemas baseados em IA, pela popularização recente da IA generativa e pela evolução contínua dos debates sobre capacidades estatais, maturidade organizacional e marcos de governança \autocite{mariano2026inteligencia} \autocite{melo2026competencias} \autocite{barros2025inteligencia}. Em contextos assim, uma revisão de literatura oferece fotografia analítica robusta, mas necessariamente provisória. Novas regulações, instrumentos institucionais e aplicações concretas podem alterar de forma relativamente rápida o diagnóstico aqui formulado. Por isso, os achados desta dissertação devem ser compreendidos como base analítica para interpretação crítica do estado da arte e para futuras pesquisas empíricas sobre o uso de IA na administração pública federal, e não como palavra final sobre um campo ainda em acelerada transformação.
\section{RESULTADOS E DISCUSSÃO}
\label{sec:org18503ae}
\subsection{Panorama da literatura selecionada}
\label{sec:orga17e527}
A literatura selecionada revela, em primeiro lugar, que o uso de inteligência artificial no setor público deixou de ser tratado como mera extensão da agenda de governo digital e passou a constituir um campo analítico próprio, no qual se entrelaçam capacidades tecnológicas, desenho institucional, regimes de responsabilização e valores públicos. Esse deslocamento é particularmente visível nas revisões mais abrangentes do período recente, que mostram a transição de estudos predominantemente exploratórios sobre automação e dados para abordagens centradas em adoção, difusão, governança e efeitos organizacionais da IA na administração pública \autocite{madan2022ai} \autocite{serra2024a} \autocite{lyrio2025emerging}. No corpus examinado, a IA aparece menos como tecnologia isolada e mais como conjunto heterogêneo de sistemas aplicados a classificação, previsão, recomendação, análise documental, atendimento automatizado e apoio à decisão, sempre condicionado por arranjos institucionais específicos. Essa constatação é central para o recorte desta dissertação, porque indica que a discussão relevante para a administração pública federal brasileira não se resume à incorporação instrumental de soluções algorítmicas, mas exige compreender como tais soluções são enquadradas por estruturas de governança, controles internos, capacidades burocráticas e exigências de legitimidade administrativa.
As revisões sistemáticas e integrativas incluídas no corpus convergem ao apontar que o campo ainda está em consolidação, mas já apresenta regularidades importantes. De um lado, há uma recorrência de narrativas orientadas por promessa: aumento de eficiência, maior capacidade analítica, aceleração de rotinas, personalização de serviços e reforço do monitoramento estatal \autocite{alhosani2024opportunities} \autocite{pi2021machine}. De outro, há crescente reconhecimento de que essas promessas são inseparáveis de problemas persistentes, como opacidade decisória, fragilidade de dados, discriminação algorítmica, dependência de fornecedores, assimetrias de capacidade e dificuldades de auditoria \autocite{batool2025ai} \autocite{sigfrids2022how} \autocite{mkander2023auditing}. O panorama, portanto, não é linear nem celebratório. Ao contrário, a literatura mais madura tende a rejeitar tanto o determinismo tecnológico quanto a recusa abstrata da inovação. Ela descreve a adoção de IA como processo político-organizacional complexo, marcado por escolhas sobre finalidade pública, distribuição de autoridade, graus de automação admissíveis e mecanismos de supervisão humana \autocite{criado2024artificial} \autocite{parycek2023artificial}.
No plano temático, o corpus pode ser lido como organizado em quatro grandes conjuntos analíticos, ainda que fortemente interligados. O primeiro reúne estudos de mapeamento e difusão, interessados em explicar como a IA entra nas organizações públicas, em quais setores é aplicada e quais fatores favorecem ou dificultam sua institucionalização \autocite{madan2022ai} \autocite{david2023understanding}. O segundo concentra-se em governança e regulação, enfatizando princípios, estruturas decisórias, accountability e coordenação interorganizacional \autocite{batool2025ai} \autocite{sigfrids2022how} \autocite{goulart2025processos}. O terceiro examina domínios de aplicação específicos, como procedimentos administrativos, saúde, recursos humanos e serviços de emprego, evidenciando que os riscos e benefícios variam conforme o tipo de decisão, a qualidade dos dados e o grau de impacto sobre direitos \autocite{parycek2023artificial} \autocite{lemes2020o} \autocite{cho2023human} \autocite{krtner2021predictive}. O quarto grupo trata de auditoria, monitoramento contínuo e competências institucionais, chamando atenção para o fato de que governar IA exige capacidades permanentes de avaliação, e não apenas diretrizes ex ante \autocite{minkkinen2022continuous} \autocite{mkander2023auditing} \autocite{melo2026competencias}. Essa organização do corpus é relevante porque mostra que a produção já não se limita a perguntar se a administração pública deve usar IA, mas como, sob quais condições e com quais salvaguardas.
Um aspecto importante do panorama é a coexistência de diferentes escalas de análise. Criado, Sandoval-Almazán e Gil-Garcia propõem compreender IA e administração pública a partir de perspectivas micro, meso e macro, articulando comportamento e capacidades dos agentes públicos, desenho organizacional e contextos regulatórios e políticos mais amplos \autocite{criado2024artificial}. Essa chave interpretativa ajuda a ordenar a diversidade do corpus. No nível micro, ganham destaque questões de discricionariedade burocrática, supervisão humana, confiança nos sistemas e competências dos servidores. No nível meso, emergem os temas de coordenação, aquisição tecnológica, governança de dados e integração entre áreas finalísticas, jurídicas e de controle. No nível macro, entram em cena marcos legais, estratégias nacionais de IA, princípios éticos, regimes de proteção de dados e padrões de responsabilização pública. A principal contribuição dessa abordagem é mostrar que falhas de implementação raramente decorrem apenas de limitações técnicas; elas refletem desajustes entre essas camadas analíticas. Para a administração pública federal, isso significa que a simples contratação ou desenvolvimento de soluções algorítmicas não resolve o problema da adoção responsável se não houver alinhamento entre objetivos administrativos, capacidades organizacionais e regras institucionais.
As revisões de Madan e Ashok e de Serra são particularmente úteis para reconstruir a evolução da agenda de pesquisa \autocite{madan2022ai} \autocite{serra2024a}. Ambas indicam que a literatura avançou de descrições genéricas de potencial para discussões mais estruturadas sobre fatores de adoção e efeitos institucionais. No entanto, também ressaltam a permanência de lacunas empíricas. Há muitos estudos propondo agendas normativas e poucos acompanhando longitudinalmente a implementação de sistemas específicos em organizações públicas. Em consequência, o campo é rico em diagnósticos sobre desafios prováveis, mas ainda relativamente escasso em evidências comparáveis sobre desempenho, legitimidade e resultados distributivos. Isso é especialmente problemático quando se pretende extrair implicações para a APF brasileira, pois a literatura internacional frequentemente trabalha com ambientes institucionais mais digitalizados, bases de dados mais integradas ou regimes regulatórios distintos. Ainda assim, o valor comparativo desses estudos é alto: eles permitem identificar padrões recorrentes de sucesso e fracasso, além de oferecer vocabulário analítico para examinar experiências federais nacionais.
Outro traço marcante do corpus é a ampliação da noção de governança de IA. Em trabalhos mais recentes, governança não significa apenas existência de normas ou comitês, mas um ecossistema de práticas, instrumentos e responsabilidades distribuídas ao longo do ciclo de vida dos sistemas \autocite{batool2025ai} \autocite{sigfrids2022how}. A revisão de Batool et al. mostra a consolidação da ideia de AI governance como campo interdisciplinar, articulando dimensões técnicas, jurídicas, éticas e organizacionais \autocite{batool2025ai}. Já Sigfrids et al. enfatizam que administrações públicas devem fomentar o desenvolvimento ético da IA por meio de estruturas que combinem princípios de alto nível com mecanismos operacionais, como avaliação de impacto, diretrizes de contratação, documentação, supervisão e espaços de deliberação \autocite{sigfrids2022how}. Em vez de tratar governança como camada externa aplicada após a implementação, essa literatura a concebe como condição constitutiva do uso legítimo da IA no setor público. Para a APF, a implicação é clara: programas de IA não podem ficar restritos a laboratórios de inovação, áreas de tecnologia ou iniciativas isoladas; precisam ser integrados às rotinas de governança institucional, controle, integridade e gestão de riscos.
No âmbito das aplicações, o corpus sugere que a administração pública adota IA principalmente em funções de triagem, classificação, previsão e apoio procedimental, e menos frequentemente em substituição plena de decisão humana. O estudo de Parycek, Schmid e Novak é ilustrativo ao analisar automação e IA em procedimentos administrativos e mostrar que o potencial transformador dessas tecnologias depende do tipo de ato administrativo, do grau de estruturação da tarefa e da disponibilidade de dados de qualidade \autocite{parycek2023artificial}. O argumento central é que ambientes decisórios intensamente juridificados e com alta incidência de direitos subjetivos exigem cautela ampliada, especialmente quando a lógica algorítmica afeta motivação, impugnação e revisibilidade dos atos. Essa discussão dialoga diretamente com o contexto federal brasileiro, no qual inúmeras atividades potencialmente automatizáveis envolvem prestações sociais, fiscalização, análise documental, reconhecimento de direitos e sanções administrativas. A literatura, nesse sentido, não sustenta um uso indiferenciado da IA; ela distingue contextos em que automação pode elevar eficiência sem comprometer garantias fundamentais e contextos em que os riscos jurídicos e éticos são substancialmente maiores.
A análise setorial reforça essa conclusão. No campo da saúde, a revisão sobre o uso de IA pela administração pública brasileira aponta oportunidades relevantes em diagnóstico, priorização e apoio à gestão, mas também ressalta desafios ligados à sensibilidade dos dados, à responsabilização por erros e à necessidade de preservar o juízo profissional em situações complexas \autocite{lemes2020o}. Em recursos humanos, tanto a literatura internacional quanto a produção mais recente em português advertem que ferramentas algorítmicas podem ampliar eficiência em recrutamento, análise de perfis e gestão de desempenho, mas carregam forte risco de reproduzir discriminações históricas e naturalizar métricas pouco transparentes \autocite{chilunjika2022artificial} \autocite{cho2023human} \autocite{barros2025inteligencia}. Nos serviços de emprego, estudos sobre algoritmos preditivos mostram que a promessa de alocação mais eficiente de políticas pode entrar em tensão com critérios de justiça, inteligibilidade e contestabilidade \autocite{krtner2021predictive}. O elemento comum entre esses domínios é que o valor da IA depende menos do sofisticação técnica em abstrato e mais da capacidade institucional de enquadrar seu uso segundo finalidades públicas claras e controles proporcionais ao impacto decisório.
Também chama atenção no corpus a crescente importância conferida às competências estatais. A literatura mais recente sobre administração pública e IA insiste que a adoção bem-sucedida não é apenas problema de infraestrutura, mas de formação organizacional. Melo destaca que competências em IA na gestão pública abrangem não só habilidades técnicas, mas também compreensão normativa, capacidade analítica, leitura crítica de modelos, gestão de fornecedores e tradução entre áreas técnicas e finalísticas \autocite{melo2026competencias}. Mariano, ao mapear assimetrias de maturidade no uso de IA e IA generativa na administração pública, reforça que a adoção se distribui de modo desigual entre órgãos e níveis de capacidade, produzindo ilhas de inovação cercadas por fragilidades institucionais \autocite{mariano2026inteligencia}. Esse achado é decisivo para a APF, marcada por elevada heterogeneidade entre ministérios, autarquias, agências, empresas estatais e órgãos de controle. O panorama da literatura sugere que políticas federais para IA precisarão lidar com assimetrias de capacidade administrativa, sob pena de gerar adoção fragmentada, dependência excessiva de soluções terceirizadas e controles inconsistentes.
Na literatura comparada e subnacional, observam-se padrões que ajudam a interpretar desafios da esfera federal. Estudos sobre governos locais mostram que dirigentes públicos percebem a IA simultaneamente como oportunidade de inovação e como fonte de riscos institucionais relacionados a custos, confiança, escassez de talentos e incerteza regulatória \autocite{yiitcanlar2022artificial} \autocite{david2023understanding}. Embora o ambiente municipal difira da APF, essas pesquisas evidenciam um ponto transferível: a decisão de adotar IA raramente decorre apenas de demonstração de eficiência; ela depende de liderança, clareza estratégica, arranjos de coordenação e percepção de legitimidade. Em contextos emergentes, como mostra Sharma, dificuldades de implementação incluem resistência organizacional, baixa qualidade de dados, carência de infraestrutura e insuficiência de competências \autocite{sharma2021implementing}. Tais obstáculos reaparecem, sob diferentes formas, nos estudos sobre setor público em geral, sugerindo que a literatura converge na identificação de barreiras institucionais persistentes. Para o caso brasileiro, isso reforça a necessidade de interpretar a IA menos como solução pronta e mais como objeto de mudança organizacional que exige planejamento, experimentação controlada e governança adaptativa.
Outro eixo decisivo do panorama é a relação entre IA, explicabilidade e auditabilidade. A literatura jurídica e técnico-institucional insiste que, no setor público, não basta que o sistema funcione; é necessário que sua lógica, seus limites e seus efeitos sejam passíveis de escrutínio institucional. Brkan problematiza a própria viabilidade da “explicação” em decisões algorítmicas complexas e mostra que a demanda regulatória por inteligibilidade enfrenta limites técnicos reais, sobretudo em modelos opacos \autocite{brkan2020legal}. Esse debate é aprofundado por Mökander, que revisa abordagens legais, éticas e técnicas de auditoria de IA e argumenta que auditoria deve ser compreendida como mecanismo central de governança, embora ainda cercado por indefinições metodológicas e conflitos de interesse \autocite{mkander2023auditing}. Minkkinen, por sua vez, avança na discussão sobre auditoria contínua, sugerindo que sistemas de IA exigem monitoramento permanente porque seu desempenho e seus riscos podem variar ao longo do tempo em função de drift, mudanças contextuais e usos imprevistos \autocite{minkkinen2022continuous}. Para a administração pública federal, esse conjunto de estudos é especialmente relevante: como muitas aplicações envolvem decisões recorrentes e de alto volume, a auditabilidade não pode ser episódica, devendo integrar arquitetura de governança, documentação, trilhas de decisão e rotinas de revisão.
No plano regional, a literatura latino-americana e ibero-americana cumpre papel corretivo importante diante do predomínio de referenciais do Norte global. O estudo comparado sobre inteligência artificial no setor público latino-americano, orientado pela Carta Ibero-Americana de IA na Administração Pública, evidencia que a região compartilha ambições de modernização e prestação de serviços mais inteligentes, mas enfrenta restrições mais severas de capacidade estatal, coordenação intergovernamental e institucionalização de salvaguardas \autocite{criado2024inteligencia}. Textos brasileiros recentes sobre governança em órgãos públicos confirmam essa percepção ao identificar processos ainda incipientes, frequentemente marcados por descontinuidade, baixa padronização e ausência de estruturas robustas de supervisão \autocite{goulart2025processos} \autocite{serra2024a}. Em consequência, o corpus sugere que o debate aplicável ao governo federal brasileiro não deve simplesmente importar modelos internacionais de responsible AI, mas adaptá-los a um ambiente administrativo no qual convivem forte juridificação, capacidade digital desigual e pressão por eficiência em larga escala.
Em síntese, o panorama da literatura selecionada mostra um campo em rápida expansão, porém ainda tensionado entre entusiasmo inovador e prudência institucional. Há consenso razoável de que a IA pode ampliar capacidades estatais em análise, automação e prestação de serviços, mas também de que seus riscos são particularmente sensíveis quando incorporados a funções públicas que afetam direitos, distribuem recursos, classificam cidadãos ou orientam fiscalização \autocite{pi2021machine} \autocite{parycek2023artificial} \autocite{batool2025ai}. O corpus converge ao indicar que o elemento decisivo não é a adoção em si, mas a forma de institucionalização: qualidade dos dados, clareza de finalidade, supervisão humana, governança de ciclo de vida, capacidade de auditoria, competências organizacionais e alinhamento com valores públicos \autocite{criado2024artificial} \autocite{sigfrids2022how} \autocite{mkander2023auditing}. Para a APF brasileira, esse panorama oferece duas conclusões preliminares. A primeira é que a literatura já fornece massa crítica suficiente para abandonar abordagens genéricas sobre “inovação com IA” e passar a discutir arranjos concretos de governança, ética, compliance e risco. A segunda é que a experiência federal precisa ser lida à luz de suas especificidades institucionais: volume decisório elevado, sensibilidade jurídica das ações administrativas, heterogeneidade de capacidades entre órgãos e crescente pressão por transparência e responsabilização. É a partir desse quadro que os demais subtítulos desta seção aprofundarão os eixos de governança identificados, os dilemas éticos e de compliance, os riscos institucionais e, por fim, as implicações mais diretamente voltadas ao governo público federal brasileiro.
\subsection{Eixos de governança identificados}
\label{sec:org72b98c4}
A leitura integrada do corpus permite sustentar que a governança de IA na administração pública não se organiza em torno de um único modelo, mas de um conjunto de eixos recorrentes que estruturam a adoção, o controle e a legitimação dessas tecnologias. Em vez de aparecer apenas como tema normativo abstrato, a governança surge na literatura como problema operacional e institucional: quem decide sobre a adoção, com base em quais critérios, com que distribuição de responsabilidades, sob quais salvaguardas e com que mecanismos de revisão ao longo do ciclo de vida dos sistemas \autocite{batool2025ai} \autocite{sigfrids2022how}. Quando essas perguntas são projetadas sobre o setor público, a governança deixa de equivaler a mera coordenação tecnológica e passa a significar a articulação entre legalidade administrativa, capacidade organizacional, gestão de dados, responsabilização e proteção do interesse público \autocite{criado2024artificial} \autocite{serra2024a}.
No corpus analisado, é possível identificar ao menos cinco eixos substantivos de governança que se repetem com variações de ênfase: governança estratégica e de alinhamento institucional; governança de dados e infraestrutura informacional; governança decisória e supervisão humana; governança de accountability, transparência e auditoria; e governança de capacidades, competências e coordenação interorganizacional. Esses eixos não são compartimentos estanques. Ao contrário, a literatura sugere que falhas em um deles tendem a contaminar os demais. Um sistema tecnicamente robusto, por exemplo, pode permanecer institucionalmente inadequado se não houver clareza sobre finalidade pública, autoridade decisória ou mecanismos de contestação; da mesma forma, diretrizes éticas genéricas perdem efetividade quando não se traduzem em rotinas de monitoramento, documentação e responsabilização \autocite{goulart2025processos} \autocite{mkander2023auditing}.
O primeiro eixo identificado é o da governança estratégica e do alinhamento institucional. Diversos estudos mostram que a adoção de IA no setor público fracassa ou produz resultados limitados quando é tratada como iniciativa isolada de inovação, desconectada da estratégia organizacional, das competências existentes e da missão institucional \autocite{madan2022ai} \autocite{sharma2021implementing}. A contribuição de Madan e Ashok é especialmente relevante ao demonstrar que adoção e difusão de IA na administração pública dependem de fatores organizacionais, ambientais e tecnológicos que interagem entre si, e não de uma simples decisão de investimento \autocite{madan2022ai}. Essa constatação aproxima-se da revisão integrativa de Serra e Machado, que destaca a necessidade de sistematizar usos e resultados da IA no setor público sob um enquadramento explícito de governança, justamente para evitar iniciativas fragmentadas e orientadas apenas por modismos tecnológicos \autocite{serra2024a}.
Nesse plano estratégico, a literatura distingue organizações que usam IA como camada incremental de automação daquelas que a inserem em arquitetura mais ampla de transformação administrativa. Essa diferença é decisiva para a administração pública federal brasileira, onde coexistem órgãos com alta maturidade digital e outros ainda dependentes de fluxos documentais pouco estruturados. Sem alinhamento institucional, a IA tende a ser implementada como solução pontual para gargalos específicos, mas sem integração com prioridades de política pública, gestão de riscos e controles internos. Parycek, Schmid e Novak argumentam que, em procedimentos administrativos, o valor da automação algorítmica depende da adequada compreensão dos limites do sistema, da natureza da decisão e das condições jurídico-organizacionais que moldam sua utilização \autocite{parycek2023artificial}. Assim, o eixo estratégico não se resume a decidir “usar IA”, mas a definir finalidades legítimas, objetos decisórios admissíveis, métricas de desempenho compatíveis com valor público e critérios para interrupção, revisão ou redimensionamento da solução.
A literatura comparada também sugere que a governança estratégica de IA exige tradução institucional de princípios amplos em estruturas concretas de decisão. Revisões sobre governança de IA mostram recorrência de princípios como transparência, fairness, accountability, segurança e supervisão humana, mas destacam que esses princípios só ganham densidade quando vinculados a papéis organizacionais, fóruns decisórios e processos de aprovação, avaliação e escalonamento \autocite{batool2025ai} \autocite{sigfrids2022how}. Em outras palavras, o problema não é a ausência de princípios, mas sua baixa operacionalização. Para a APF, isso implica reconhecer que instrumentos de estratégia digital, governança corporativa, gestão de riscos e integridade precisam dialogar entre si, evitando a duplicação de estruturas normativas sem efetividade prática.
O segundo eixo é a governança de dados e da infraestrutura informacional. A literatura é quase unânime ao afirmar que a qualidade da governança de IA é inseparável da qualidade da governança de dados. Isso ocorre porque sistemas algorítmicos dependem não apenas de capacidade computacional, mas de bases íntegras, representativas, rastreáveis e juridicamente utilizáveis \autocite{parycek2023artificial} \autocite{pi2021machine}. Em administrações públicas, esse problema assume contornos particulares: dados são produzidos em contextos burocráticos heterogêneos, carregam classificações administrativas historicamente sedimentadas e frequentemente refletem desigualdades ou vieses institucionais anteriores à modelagem estatística. Portanto, a governança de dados não é etapa técnica preliminar; ela é núcleo da governança pública da IA.
A revisão de Batool et al. reforça que incidentes e riscos de IA frequentemente decorrem de fragilidades estruturais na cadeia de dados, incluindo opacidade sobre origem, critérios de curadoria, atualização e uso secundário das bases \autocite{batool2025ai}. Parycek et al., ao discutir procedimentos administrativos, enfatizam que a validade das predições ou classificações depende criticamente da disponibilidade e adequação dos dados, e que o uso de dados inadequados compromete não apenas a eficiência do sistema, mas sua aceitabilidade jurídica e administrativa \autocite{parycek2023artificial}. No contexto federal brasileiro, isso sugere que a expansão de IA em áreas como fiscalização, triagem documental, benefícios, saúde ou gestão de pessoas depende menos da mera aquisição de modelos e mais da institucionalização de rotinas de catalogação, interoperabilidade, controle de qualidade e documentação de bases. Sem isso, a administração incorre no risco de automatizar inconsistências já presentes em seus cadastros e fluxos de trabalho.
Esse eixo de governança de dados também se conecta ao debate sobre privacidade, proteção de dados pessoais e explicabilidade. Ainda que a presente subseção não esgote a dimensão de compliance, importa notar que Brkan mostra os limites da ideia de uma explicação plena e simples para decisões algorítmicas complexas, indicando que a inteligibilidade relevante para fins regulatórios e institucionais pode exigir soluções graduadas, contextuais e compatíveis com a natureza técnica do sistema \autocite{brkan2020legal}. Isso tem implicações diretas para a governança de dados na APF: não basta armazenar dados e modelos; é necessário manter trilhas documentais sobre fontes, transformações, variáveis utilizadas, critérios de exclusão, finalidades e limitações do sistema. A governança de IA, aqui, aparece como governança da memória institucional do sistema.
O terceiro eixo identificado é o da governança decisória e da supervisão humana. Esse é talvez o ponto em que a literatura mais claramente se distancia de leituras tecnocêntricas. Ao analisar IA na administração pública a partir de perspectivas micro, meso e macro, Criado, Sandoval-Almazán e Gil-Garcia mostram que os efeitos da IA dependem de como ela reconfigura papéis, rotinas, autonomia e relações entre agentes públicos, organizações e ambientes regulatórios \autocite{criado2024artificial}. Em termos de governança, isso significa que a questão central não é apenas se há “humano no circuito”, mas como a autoridade humana é preservada, redistribuída ou esvaziada no desenho concreto dos processos.
Nos procedimentos administrativos, Parycek et al. demonstram que a admissibilidade e a utilidade da automação variam conforme a estrutura da tarefa: atividades padronizadas, repetitivas e fortemente baseadas em regras tendem a ser mais suscetíveis à automação, ao passo que decisões complexas, valorativas ou intensamente dependentes de contexto exigem maior cautela e intervenção humana substantiva \autocite{parycek2023artificial}. A literatura, contudo, adverte que a mera previsão formal de revisão humana não garante controle efetivo. Há risco de “automation bias”, deferência excessiva à recomendação algorítmica e redução progressiva da capacidade crítica dos servidores, especialmente em contextos de pressão por produtividade ou baixa compreensão técnica do sistema \autocite{pi2021machine} \autocite{sigfrids2022how}.
Por isso, a governança decisória aparece no corpus como eixo que requer delimitação clara entre apoio à decisão, recomendação, priorização, classificação e decisão automatizada propriamente dita. Essa distinção é crucial para o governo federal brasileiro. Em áreas como análise de benefícios, fiscalização, contratação, processamento documental e atendimento digital, soluções de IA podem operar em graus variados de influência sobre a decisão final. A literatura sugere que governança robusta demanda explicitar esse grau de influência, definir quando a revisão humana é obrigatória, registrar divergências entre servidor e sistema e prever canais de contestação por parte dos administrados \autocite{sigfrids2022how} \autocite{criado2024inteligencia}. Não se trata de preservar o humano como ornamento legitimador, mas de desenhar responsabilidade decisória efetiva, verificável e compatível com os princípios do direito administrativo.
O quarto eixo é a governança de accountability, transparência e auditoria. O corpus mostra que a governança de IA se distingue de arranjos clássicos de TI justamente porque exige capacidades de escrutínio contínuo sobre sistemas que aprendem, se adaptam ou operam sobre cadeias de dados de alta complexidade. Mökander argumenta que a auditoria de IA constitui campo em expansão e que suas práticas podem aprender com tradições de auditoria financeira, segurança e ciências sociais, mas precisam responder a desafios próprios, como ambiguidade de critérios, dependência de fornecedores e dificuldade de acesso a artefatos técnicos \autocite{mkander2023auditing}. A auditoria, nesse sentido, não é elemento periférico; é mecanismo estruturante de governança.
A relevância desse eixo é aprofundada por Minkkinen et al., que tratam da auditoria contínua de IA como resposta à natureza dinâmica desses sistemas \autocite{minkkinen2022continuous}. Em vez de avaliações únicas ex ante, a literatura passa a defender monitoramento ao longo do ciclo de vida, com indicadores de desempenho, fairness, drift, conformidade e robustez. Para a administração pública, essa perspectiva é especialmente importante porque decisões inicialmente aceitáveis podem tornar-se problemáticas em razão de mudança no ambiente regulatório, alteração do perfil dos dados, reconfiguração das populações afetadas ou uso ampliado para finalidades distintas das originais. Assim, governança de IA demanda capacidade de revisar premissas, reavaliar riscos e recalibrar modelos periodicamente.
A transparência, por sua vez, aparece na literatura de forma mais sofisticada do que o simples ideal de abertura irrestrita do código. Mökander e Brkan, por caminhos distintos, ajudam a compreender que transparência relevante para governança pode envolver documentação de finalidade, desenho, dados, limites, responsabilidades e processos de revisão, mesmo quando a plena explicação técnica do modelo é difícil ou insuficiente \autocite{mkander2023auditing} \autocite{brkan2020legal}. No setor público, isso se traduz em exigência de transparência multicamada: transparência para gestores, para órgãos de controle, para equipes técnicas e para cidadãos afetados. A inexistência dessa arquitetura informacional compromete accountability, dificulta contestação e enfraquece a legitimidade da ação estatal mediada por IA.
O quinto eixo é a governança de capacidades, competências e coordenação interorganizacional. A literatura é incisiva ao mostrar que governar IA não é apenas adquirir ferramentas, mas desenvolver capacidades burocráticas para contratá-las, compreendê-las, supervisioná-las e integrá-las aos processos institucionais \autocite{melo2026competencias} \autocite{madan2022ai}. Melo identifica, por meio de revisão sistemática, que competências em IA na gestão pública abrangem dimensões técnicas, analíticas, éticas e gerenciais, o que impede reduzir o tema à formação de especialistas em ciência de dados \autocite{melo2026competencias}. Para além de quadros técnicos, a governança requer dirigentes capazes de decidir sobre riscos e finalidades, áreas jurídicas aptas a traduzir obrigações regulatórias, unidades de controle com repertório para auditar sistemas e servidores finalísticos preparados para interagir criticamente com recomendações algorítmicas.
Esse achado é consistente com estudos sobre recursos humanos e analytics no setor público. Chilunjika et al. mostram que aplicações de IA em gestão de pessoas abrem oportunidades de eficiência, mas também expõem a administração a riscos de exclusão, assimetria de acesso e fragilidade na tomada de decisão se não houver competências e salvaguardas adequadas \autocite{chilunjika2022artificial}. Cho, ao discutir analytics para gestão de pessoal, reforça que ganhos analíticos precisam ser balanceados com cautelas relativas à qualidade dos dados, interpretação dos resultados e impactos sobre direitos e confiança institucional \autocite{cho2023human}. Barros, em diálogo com a realidade da administração pública, igualmente associa IA em gestão de pessoas a desafios de ética, inovação e governança, sugerindo que o problema central não é a utilização em si, mas o desenho institucional que condiciona seus efeitos \autocite{barros2025inteligencia}.
No plano organizacional mais amplo, a coordenação interorganizacional emerge como componente decisivo desse eixo. Goulart, ao revisar processos de governança em órgãos públicos brasileiros, evidencia que a institucionalização da IA depende de arranjos que conectem áreas de tecnologia, governança, integridade, negócio e controle, evitando tanto a centralização excessiva quanto a pulverização decisória \autocite{goulart2025processos}. Essa conclusão dialoga com estudos sobre governos locais e adoção tecnológica, nos quais aparecem com frequência obstáculos como falta de liderança, insuficiência de recursos, baixa cultura analítica e ausência de estratégias coordenadas \autocite{yiitcanlar2022artificial} \autocite{david2023understanding}. Embora parte dessa literatura se concentre em níveis subnacionais ou contextos não brasileiros, os mecanismos explicativos são úteis para a APF: organizações complexas, com múltiplas carreiras, sistemas legados e forte segmentação funcional tendem a enfrentar dificuldades semelhantes de articulação e de compartilhamento de responsabilidades.
Além desses cinco eixos centrais, o corpus sugere um eixo transversal de governança orientada ao valor público. Alhosani et al. e Pi destacam benefícios frequentemente associados à IA em serviços governamentais, como eficiência, melhor alocação de recursos, personalização e apoio analítico \autocite{alhosani2024opportunities} \autocite{pi2021machine}. Entretanto, a literatura mais crítica insiste que esses benefícios só se convertem em valor público quando mediados por critérios de legitimidade, equidade, confiabilidade e responsividade institucional \autocite{criado2024artificial} \autocite{sigfrids2022how}. Esse ponto é importante porque evita uma leitura puramente gerencialista dos eixos de governança. No setor público, governar IA não significa apenas maximizar desempenho; significa decidir que tipo de desempenho importa, para quem, com quais salvaguardas e com quais trade-offs explicitados.
Em termos comparativos, observa-se que as abordagens mais maduras do corpus convergem menos em oferecer modelos fechados de governança e mais em sustentar requisitos estruturais mínimos. Entre eles, destacam-se: definição clara de finalidade pública; classificação dos casos de uso por risco e impacto; governança de dados com documentação e rastreabilidade; papéis institucionais definidos para decisão, supervisão e revisão; mecanismos de auditoria e monitoramento contínuo; e investimento em competências e coordenação \autocite{batool2025ai} \autocite{mkander2023auditing} \autocite{minkkinen2022continuous} \autocite{melo2026competencias}. A heterogeneidade dos contextos administrativos explica por que não há um desenho universal, mas o corpus é consistente ao mostrar que a ausência desses elementos aumenta a probabilidade de opacidade, dependência tecnológica, erros decisórios e perda de confiança.
Para a administração pública federal brasileira, os eixos identificados têm implicações diretas. Primeiro, indicam que iniciativas de IA não deveriam ser avaliadas apenas por métricas de inovação ou digitalização, mas por sua inserção em arranjos de governança capazes de sustentar legalidade, accountability e aprendizado institucional. Segundo, revelam que a maturidade em IA não depende exclusivamente do grau de sofisticação técnica, mas da capacidade de combinar estratégia, dados, supervisão, auditoria e competências. Terceiro, sugerem que a principal fronteira de governança talvez não esteja na formulação de novos princípios abstratos, e sim na tradução desses princípios em rotinas administrativas, protocolos de documentação, instâncias de decisão e circuitos de controle permanentemente ativáveis.
Finalmente, o corpus permite concluir que a governança de IA no setor público deve ser entendida como capacidade institucional de enquadrar tecnologias probabilísticas dentro de um regime administrativo orientado por responsabilidade pública. Essa formulação sintetiza a principal contribuição dos estudos analisados: IA não é ingovernável, mas tampouco é governada apenas por boas intenções, códigos éticos ou decisões de aquisição. Ela exige um ecossistema institucional no qual estratégia, dados, decisão, auditoria e competências se reforcem mutuamente \autocite{serra2024a} \autocite{criado2024artificial} \autocite{goulart2025processos}. É a partir dessa constatação que os dilemas éticos, os desafios de compliance e os riscos institucionais, examinados nas subseções seguintes, podem ser compreendidos não como externalidades da inovação, mas como dimensões constitutivas da própria governança da inteligência artificial na administração pública.
\subsection{Dilemas éticos e desafios de compliance}
\label{sec:org56e5b82}
Se, na subseção anterior, a governança apareceu como arquitetura institucional para orientar a adoção de IA, nesta etapa dos resultados ela se revela também como resposta a um problema mais profundo: a incorporação de sistemas algorítmicos na administração pública federal desloca e reconfigura dilemas éticos já conhecidos da ação estatal, ao mesmo tempo em que produz desafios de compliance que não podem ser resolvidos apenas pela transposição de controles tradicionais. A literatura do corpus converge em um ponto central: no setor público, a questão ética da IA não se limita ao risco de mau uso tecnológico, mas envolve a própria legitimidade da decisão administrativa mediada por dados, modelos estatísticos e formas automatizadas de triagem, recomendação, classificação e priorização \autocite{serra2024a} \autocite{criado2024artificial} \autocite{sigfrids2022how}. Isso significa que o debate sobre ética e compliance não é periférico à implementação, mas constitutivo da definição de quando, como e sob quais condições a IA pode ser usada de modo compatível com o interesse público.
A primeira implicação analítica desse achado é que a literatura revisada rejeita, ainda que por caminhos distintos, uma visão instrumental simplificada segundo a qual bastaria avaliar ganhos de eficiência para justificar a adoção de IA. Em revisões amplas sobre difusão de IA na administração pública, observa-se que os incentivos à automação costumam ser formulados em termos de produtividade, redução de custos, aceleração de fluxos decisórios e melhoria do serviço ao cidadão; contudo, os estudos também demonstram que tais benefícios são condicionados por variáveis institucionais, jurídicas e éticas que, se negligenciadas, convertem vantagens esperadas em novas fontes de erro, opacidade e contestação \autocite{madan2022ai} \autocite{alhosani2024opportunities} \autocite{pi2021machine}. No corpus, ética e compliance surgem precisamente como os domínios que impõem limites à lógica de solução tecnológica. Eles lembram que, na administração pública, decidir mais rápido não equivale necessariamente a decidir melhor, e decidir com base em padrões extraídos de bases históricas tampouco significa decidir de maneira justa.
Esse ponto é particularmente importante para a APF porque grande parte das aplicações descritas na literatura envolve contextos em que a administração não apenas presta serviços, mas exerce autoridade, distribui benefícios, define prioridades, fiscaliza condutas ou interfere em trajetórias funcionais e direitos sociais. Em procedimentos administrativos, a adoção de IA altera a gramática da decisão pública ao introduzir camadas técnicas que podem obscurecer o nexo entre fato, norma e decisão. Parycek, Schmid e Novak mostram que o uso de IA em procedimentos administrativos depende fortemente do tipo de tarefa, da estrutura legal do processo, da qualidade e disponibilidade dos dados e da possibilidade de compreender os limites do sistema \autocite{parycek2023artificial}. Assim, o dilema ético central não é apenas o da automação em si, mas o da deslocação da racionalidade administrativa para instrumentos cuja lógica probabilística nem sempre se ajusta às exigências de motivação, individualização e revisibilidade típicas do direito administrativo.
A opacidade algorítmica ocupa, nesse cenário, posição privilegiada entre os dilemas mapeados. A literatura sugere que a chamada “caixa-preta” não é um problema exclusivamente técnico, mas também jurídico e organizacional. Em sentido estrito, a opacidade diz respeito à dificuldade de compreender o funcionamento interno de certos modelos; em sentido mais amplo, ela decorre da combinação entre complexidade técnica, contratação de soluções externas, documentação insuficiente, ausência de trilhas de decisão e baixa capacitação dos agentes públicos para interpretar outputs algorítmicos \autocite{brkan2020legal} \autocite{mkander2023auditing}. Brkan, ao discutir os limites da pretensão de explicação de decisões algorítmicas, demonstra que o problema não se resolve com promessas genéricas de transparência, pois há tensões reais entre interpretabilidade, desempenho do modelo, segredo comercial e inteligibilidade efetiva para os afetados \autocite{brkan2020legal}. Para a APF, isso tem consequências diretas: não basta exigir que o fornecedor ou a área técnica “explique” o sistema; é preciso definir qual tipo de explicação é relevante para fins de legalidade, accountability, contraditório, auditoria e controle interno.
A literatura de governança ética também mostra que a transparência, quando tratada apenas como princípio abstrato, tende a ser superestimada. Sigfrids e colaboradores observam que propostas para fomentar o desenvolvimento ético da IA nas administrações públicas frequentemente concentram-se em princípios como fairness, transparência, accountability e supervisão humana, mas deixam em aberto os mecanismos institucionais concretos de implementação \autocite{sigfrids2022how}. Batool e coautores, em revisão sistemática sobre governança de IA, reforçam que a efetividade de princípios depende de sua tradução em processos, papéis, indicadores e controles verificáveis \autocite{batool2025ai}. No corpus analisado, portanto, o dilema não é a ausência de consenso normativo sobre valores desejáveis, e sim a fragilidade de arranjos capazes de converter esses valores em rotinas de compliance. Em termos práticos, isso implica reconhecer que códigos de ética, cartas de princípios ou diretrizes de uso responsável têm utilidade limitada se não forem acompanhados de avaliação prévia de impacto, documentação contínua, revisão de bases de treinamento, segregação de funções, comitês de decisão e canais de contestação.
Outro eixo decisivo diz respeito ao viés, à discriminação e à reprodução de desigualdades. A literatura é consistente ao apontar que sistemas de IA aprendem a partir de dados socialmente produzidos e, por isso, podem amplificar padrões históricos de exclusão, invisibilização ou tratamento desigual. Esse problema assume particular gravidade no setor público porque a administração lida com populações vulneráveis, critérios distributivos e decisões que podem afetar acesso a emprego, benefícios, saúde, fiscalização ou prioridades de atendimento \autocite{pi2021machine} \autocite{serra2024a}. Em aplicações de emprego público e gestão de pessoas, por exemplo, a promessa de maior objetividade pode ocultar discriminações reembaladas em linguagem analítica. Cho, ao discutir analytics para gestão de pessoal no setor público, mostra que o uso de dados em recursos humanos oferece potencial relevante, mas traz caveats importantes relacionados a privacidade, validade inferencial, enviesamentos e uso indevido de proxies \autocite{cho2023human}. De modo semelhante, estudos sobre IA e gestão de pessoas na administração pública destacam o risco de transformar avaliações probabilísticas em critérios de seleção, promoção ou monitoramento sem suficiente debate sobre justiça organizacional e direitos dos servidores \autocite{chilunjika2022artificial} \autocite{barros2025inteligencia}.
A relevância desse debate para a APF é ampliada pelo fato de que sistemas algorítmicos podem operar tanto em relações externas com cidadãos quanto em relações internas de gestão. Isso significa que o problema ético não se restringe à prestação de serviços digitais, alcançando também recrutamento, alocação de força de trabalho, avaliação de desempenho, triagem de currículos, predição de rotatividade e classificação de perfis funcionais. A literatura indica que, nesses contextos, a discriminação pode ocorrer não apenas por atributos protegidos diretamente observáveis, mas por correlações indiretas, dados históricos enviesados e desenhos de problema que já embutem determinadas concepções normativas sobre mérito, produtividade ou risco \autocite{cho2023human} \autocite{chilunjika2022artificial}. Para fins de compliance, isso exige que a administração não trate viés como defeito residual a ser corrigido a posteriori, mas como risco inerente ao ciclo de vida do sistema, demandando monitoramento contínuo, testes de desempenho por grupos, validação contextual e revisão periódica de objetivos.
Há, ademais, um dilema ético especificamente relacionado à substituição ou esvaziamento da discricionariedade administrativa. Parte da literatura alerta que a IA não elimina a discricionariedade; ela a redistribui. Escolhas sobre quais dados coletar, como rotular observações, que variável otimizar, que limiares definir e que exceções admitir já são, em si, decisões normativas \autocite{criado2024artificial} \autocite{parycek2023artificial}. Quando essas escolhas são naturalizadas como opções técnicas, ocorre um deslocamento problemático da responsabilidade. O agente público passa a se apoiar no sistema como se este apenas refletisse a realidade, quando na verdade incorpora pressupostos sobre causalidade, prioridade e valor. Em administrações marcadas por pressões por produtividade, esse fenômeno pode produzir o chamado automation bias: a tendência de aceitar recomendações automatizadas com escrutínio insuficiente, especialmente quando os operadores não dispõem de tempo, capacitação ou autoridade para divergir do sistema \autocite{pi2021machine} \autocite{madan2022ai}. Para a APF, isso é crítico em áreas de fiscalização, triagem documental, priorização de casos e apoio à decisão, nas quais a recomendação algorítmica pode ser percebida como tecnicamente neutra e, portanto, politicamente menos contestável.
O problema da supervisão humana, contudo, é tratado no corpus de forma menos ingênua do que em grande parte do debate normativo. A literatura não presume que a mera presença de um humano “no circuito” resolva os dilemas éticos. Ao contrário, sugere que a supervisão humana só é efetiva quando o agente possui compreensão suficiente do sistema, autonomia decisória real, acesso a justificativas úteis e condições organizacionais para revisar ou recusar recomendações \autocite{sigfrids2022how} \autocite{mkander2023auditing}. Se a revisão humana se limitar a validar outputs em massa, sem tempo, sem competência técnica e sem responsabilidade claramente definida, ela se converte em formalidade de compliance e não em salvaguarda substantiva. Esse achado é especialmente relevante para o contexto federal brasileiro, em que assimetrias de capacidade entre órgãos podem estimular adoções heterogêneas e controles meramente nominais.
O tema da privacidade e da proteção de dados também ocupa lugar central entre os dilemas éticos e desafios de conformidade. Embora o corpus não se restrinja a análises jurídicas da proteção de dados, vários trabalhos apontam que a expansão da IA no setor público depende de intensa coleta, integração, limpeza e reutilização de bases administrativas, frequentemente produzidas para finalidades distintas daquelas assumidas em novos modelos analíticos \autocite{parycek2023artificial} \autocite{serra2024a}. Essa reconfiguração funcional dos dados cria tensões entre eficiência administrativa, minimização de dados, finalidade específica, proporcionalidade e segurança da informação. No setor público, tais tensões são agravadas pela assimetria entre Estado e administrado: o cidadão raramente pode escolher não participar de certos sistemas ou negar fornecimento de dados em interações essenciais com a administração. Por isso, a literatura sugere que a conformidade em IA deve ser mais exigente do que a mera observância formal de regras de tratamento de dados; ela deve incluir justificativa robusta de necessidade, governança de acesso, rastreabilidade de transformações dos dados e controles específicos para evitar usos secundários incompatíveis com o interesse público \autocite{batool2025ai} \autocite{mkander2023auditing}.
No campo do compliance, um achado recorrente do corpus é a insuficiência de modelos tradicionais baseados apenas em aderência documental e verificação ex post. A IA desafia programas de integridade e conformidade porque seu comportamento pode mudar ao longo do tempo, seus erros podem ser sistêmicos e seu desempenho depende de fatores sociotécnicos distribuídos por diferentes áreas organizacionais. Mökander argumenta que a auditoria de IA deve articular abordagens legais, éticas e técnicas, aprendendo com tradições de auditoria financeira, engenharia de segurança e ciências sociais, mas sem presumir equivalência automática entre esses campos \autocite{mkander2023auditing}. Minkkinen, ao tratar da auditoria contínua de IA, destaca que controles pontuais são insuficientes em sistemas que dependem de dados dinâmicos, atualizações frequentes e monitoramento persistente de desempenho \autocite{minkkinen2022continuous}. A implicação para a APF é clara: compliance em IA precisa ser concebido como processo longitudinal, cobrindo desenho, desenvolvimento, contratação, implantação, operação, reavaliação e descontinuação.
Essa perspectiva longitudinal expõe outra dificuldade importante: a fragmentação de responsabilidades. Em muitos arranjos públicos, a área demandante entende o sistema como ferramenta técnica; a equipe de TI o vê como solução operacional; a assessoria jurídica intervém apenas em fase tardia; o controle interno atua depois da implantação; e a alta gestão se mobiliza sobretudo em torno de metas de inovação. A literatura sugere que essa distribuição difusa favorece lacunas de accountability, porque nenhum ator detém visão completa do risco ético e regulatório \autocite{criado2024artificial} \autocite{goulart2025processos}. Em revisões sobre processos de governança em órgãos públicos brasileiros, nota-se justamente a necessidade de integrar governança digital, governança de dados, controles internos e diretrizes de integridade para que a adoção de IA não ocorra em zonas cinzentas de responsabilidade \autocite{goulart2025processos}. Nesse sentido, compliance deixa de ser função exclusiva de unidades de controle e passa a exigir desenho institucional transversal, com responsabilidades explícitas entre gestores de negócio, cientistas de dados, equipes jurídicas, áreas de segurança da informação, corregedorias e auditoria interna.
As dificuldades se intensificam quando a administração depende de fornecedores privados, plataformas proprietárias ou serviços de IA generativa. Embora a literatura revisada trate em graus variados dessas tecnologias, o corpus mais recente já identifica uma assimetria crescente entre capacidade estatal de contratar e fiscalizar e a velocidade de evolução das soluções ofertadas \autocite{mariano2026inteligencia} \autocite{lyrio2025emerging}. Isso gera dilemas éticos e de compliance adicionais: como assegurar explicabilidade mínima em sistemas fechados? Como verificar qualidade dos dados de treinamento quando os modelos são desenvolvidos externamente? Como evitar lock-in tecnológico que reduza a possibilidade de auditoria independente ou descontinuação segura? Como compatibilizar cláusulas contratuais, segredo comercial e exigências de accountability pública? Tais perguntas são particularmente sensíveis na APF porque o regime jurídico-administrativo exige motivação, controle e impessoalidade, ao passo que muitos produtos de mercado são desenhados para ambientes de experimentação privada, com menor exigência de publicidade e revisão contestatória.
A literatura comparada também chama atenção para o risco de “ethics washing” ou “compliance washing”, isto é, a adoção de linguagem ética sem transformação real das práticas institucionais \autocite{batool2025ai} \autocite{sigfrids2022how}. No setor público, esse risco se manifesta quando órgãos publicam princípios de IA responsável, criam comitês consultivos ou anunciam supervisão humana, mas mantêm baixa documentação dos sistemas, ausência de avaliação de impacto, indicadores insuficientes de desempenho distributivo e canais frágeis de recurso para os afetados. Em outras palavras, a retórica da inovação responsável pode operar como mecanismo de legitimação simbólica, sem alterar materialmente as condições de governabilidade dos sistemas. Esse é um ponto particularmente relevante para a administração federal brasileira, que convive com forte incentivo à transformação digital e com pressões por rápida entrega de resultados. Sem critérios substantivos de verificação, a ética pode ser absorvida como discurso reputacional, enquanto o compliance se reduz a checklist procedimental.
Os estudos sobre adoção tecnológica em governos locais e organizações públicas reforçam ainda que percepções gerenciais de oportunidade frequentemente coexistem com limitada compreensão dos constrangimentos éticos e regulatórios \autocite{yiitcanlar2022artificial} \autocite{david2023understanding}. Isso ajuda a explicar por que projetos de IA são, por vezes, concebidos sob a lógica de modernização inevitável, minimizando controvérsias sobre exclusão digital, erro classificatório, desvio de finalidade e contestabilidade. A revisão de Madan e Ashok indica que fatores contextuais e organizacionais moldam a difusão da IA tanto quanto seus atributos técnicos \autocite{madan2022ai}. Assim, em ambientes institucionais nos quais inovação é premiada e falhas distributivas são pouco mensuradas, os dilemas éticos tendem a ser percebidos como barreiras externas, e não como critérios internos de desenho. O desafio de compliance, então, não é apenas normativo, mas cultural: criar capacidades para que a organização reconheça que limitações éticas e jurídicas são parte do problema de implementação, e não obstáculos externos à tecnologia.
No caso brasileiro, a literatura sobre saúde pública e administração pública sugere que usos setoriais da IA evidenciam uma tensão recorrente entre potencial de ganho operacional e deveres reforçados de cuidado institucional \autocite{lemes2020o} \autocite{serra2024a}. Em áreas sensíveis, decisões baseadas em predição ou classificação podem afetar diretamente pessoas em situação de vulnerabilidade, o que eleva o ônus de justificativa, precisão, segurança e revisão. O mesmo vale para algoritmos preditivos em serviços de emprego e alocação de políticas, cuja utilização pode racionalizar recursos, mas também cristalizar expectativas normativas sobre quem merece prioridade, risco ou acompanhamento intensivo \autocite{krtner2021predictive}. Nessas situações, compliance não deve ser confundido com mera conformidade processual; ele exige avaliação de adequação material entre o uso da IA e os princípios da administração pública, incluindo impessoalidade, moralidade, publicidade, eficiência em sentido qualificado e proteção de direitos.
Em síntese, os resultados do corpus permitem afirmar que os dilemas éticos e desafios de compliance da IA no setor público federal brasileiro não são acessórios à agenda de transformação digital, mas sua prova mais exigente. A literatura converge em apontar que os principais riscos se concentram na opacidade decisória, no viés e na discriminação, na fragilidade da supervisão humana, na expansão pouco controlada do uso de dados, na fragmentação de responsabilidades e na insuficiência de mecanismos tradicionais de auditoria e conformidade \autocite{mkander2023auditing} \autocite{minkkinen2022continuous} \autocite{criado2024artificial}. Ao mesmo tempo, o corpus sugere que esses problemas não devem levar a uma rejeição abstrata da IA, mas à formulação de condições institucionais mais rigorosas para seu uso legítimo. Para a APF, isso significa deslocar o foco da pergunta “onde a IA pode gerar eficiência?” para uma indagação mais exigente: em quais contextos a IA pode ser usada sem degradar a legalidade, a equidade, a explicabilidade, a contestabilidade e a responsabilidade pública? Essa inflexão é decisiva porque redefine a inovação não como exceção ao controle, mas como objeto ampliado de governança, integridade e accountability \autocite{sigfrids2022how} \autocite{batool2025ai} \autocite{goulart2025processos}.
\subsection{Riscos institucionais e mecanismos de mitigação}
\label{sec:orgdddb57c}
Se a subseção anterior evidenciou que os dilemas éticos e de compliance decorrem da incorporação de sistemas algorítmicos a processos decisórios estatais, o passo seguinte consiste em demonstrar que tais dilemas se materializam, no plano organizacional, como riscos institucionais. A literatura do corpus sugere que a adoção de IA no setor público não deve ser lida apenas como um problema de desempenho tecnológico, mas como uma transformação na ecologia decisória da administração, com efeitos sobre legalidade, coordenação, capacidade burocrática, confiança pública, distribuição de responsabilidades e sustentabilidade da ação estatal \autocite{criado2024artificial} \autocite{madan2022ai} \autocite{batool2025ai}. Em outras palavras, o risco institucional emerge quando falhas, vieses, opacidades ou dependências ligadas à IA deixam de ser incidentes pontuais e passam a afetar rotinas, competências, controles e expectativas normativas que sustentam a legitimidade da administração pública.
Esse enquadramento é especialmente relevante para o governo público federal brasileiro porque a IA tende a ser empregada em domínios de elevada sensibilidade administrativa: automação de triagens, análise documental, apoio à fiscalização, detecção de anomalias, gestão de pessoas, classificação de demandas, serviços digitais e, progressivamente, uso de ferramentas generativas para apoio a atividades internas e comunicação institucional. Em todos esses casos, o risco não reside apenas em o sistema errar, mas em a organização passar a estruturar fluxos de trabalho, alocação de pessoal e critérios de decisão a partir de infraestruturas técnicas cuja governança é frequentemente desigual entre órgãos, áreas e fornecedores \autocite{mariano2026inteligencia} \autocite{lyrio2025emerging} \autocite{goulart2025processos}. Assim, a discussão sobre risco institucional exige deslocar a análise da aplicação isolada para o arranjo organizacional mais amplo em que ela opera.
Uma primeira família de riscos identificada na literatura é a do desalinhamento entre capacidades tecnológicas e capacidades institucionais. Revisões sobre adoção e difusão de IA mostram que o entusiasmo com ganhos de eficiência frequentemente antecede a consolidação de estruturas internas de governança, competências técnicas e modelos de supervisão compatíveis com a complexidade da tecnologia \autocite{madan2022ai} \autocite{alhosani2024opportunities}. Esse descompasso produz um padrão recorrente: órgãos passam a experimentar soluções de IA sem maturidade suficiente para definir escopo, aferir qualidade dos dados, documentar decisões de desenho, monitorar desempenho ao longo do tempo e revisar efeitos adversos. A consequência é um risco organizacional de implementação aparente, em que a tecnologia é formalmente incorporada, mas a instituição não desenvolve meios efetivos para controlá-la. Sharma, ao tratar dos desafios de implementação da IA, indica que barreiras como infraestrutura inadequada, resistência organizacional, deficiência de dados e lacunas de competências não constituem obstáculos periféricos, mas fatores centrais para o fracasso ou uso distorcido da tecnologia \autocite{sharma2021implementing}. No contexto da APF, isso implica reconhecer que programas de IA sem investimento correlato em arquitetura de governança podem aumentar, e não reduzir, a exposição institucional ao erro e à contestação.
Uma segunda família de riscos diz respeito à qualidade, integridade e adequação dos dados. Parycek, Schmid e Novak enfatizam que a validade de sistemas de IA em procedimentos administrativos depende criticamente da disponibilidade, consistência e pertinência dos dados utilizados no treinamento e na operação dos modelos \autocite{parycek2023artificial}. Esse ponto é decisivo no setor público, onde bases históricas refletem assimetrias territoriais, lacunas cadastrais, padrões administrativos contingentes e, muitas vezes, critérios decisórios heterogêneos sedimentados ao longo do tempo. Quando tais bases são tratadas como representação neutra da realidade, os sistemas tendem a reproduzir e reforçar distorções pré-existentes. O risco institucional, nesse caso, não se limita à presença de viés estatístico, mas alcança a cristalização burocrática de desigualdades por meio de instrumentos revestidos de autoridade técnica. Em aplicações voltadas à fiscalização, ao reconhecimento de padrões ou à priorização de atendimento, a má qualidade dos dados pode deslocar recursos públicos para alvos indevidos, produzir falsos positivos ou invisibilizar grupos subatendidos, gerando custos administrativos e reputacionais cumulativos \autocite{pi2021machine} \autocite{serra2024a}.
A literatura também mostra que o risco de viés adquire particular gravidade quando a IA é aplicada à gestão de pessoas e a contextos laborais no setor público. Estudos sobre recursos humanos alertam que sistemas analíticos e algoritmos de apoio à seleção, avaliação ou alocação de servidores podem importar vieses históricos, penalizar trajetórias não lineares e reduzir a complexidade do desempenho funcional a indicadores de fácil mensuração \autocite{cho2023human} \autocite{chilunjika2022artificial} \autocite{barros2025inteligencia}. Em termos institucionais, isso pode afetar não apenas indivíduos, mas a própria capacidade estatal, ao induzir padrões de gestão que privilegiam conformidade quantitativa em detrimento de julgamento profissional, diversidade de perfis e aprendizado organizacional. Para a APF, em que carreiras, concursos, avaliações e mobilidade funcional possuem forte densidade normativa, o uso de IA em gestão de pessoas demanda cautela redobrada, sob pena de converter instrumentos de modernização em fontes de litigiosidade, desmotivação e erosão da confiança interna na imparcialidade administrativa.
Uma terceira família de riscos está associada à opacidade organizacional e à difusão da irresponsabilidade. A literatura sobre governança de IA destaca que, em instituições complexas, a introdução de sistemas algorítmicos tende a fragmentar a cadeia de responsabilização: gestores contratam ou aprovam soluções sem domínio técnico; equipes técnicas modelam sistemas sem controle sobre os usos concretos; operadores utilizam outputs sem compreender seus limites; e fornecedores preservam partes relevantes do funcionamento sob proteção contratual ou segredo comercial \autocite{batool2025ai} \autocite{mkander2023auditing} \autocite{brkan2020legal}. O resultado é um ambiente em que todos participam do sistema, mas poucos conseguem responder integralmente por seus efeitos. Mökander argumenta que a auditoria de IA vem justamente suprir essa lacuna, ao buscar estruturar evidências, critérios e procedimentos de verificação que permitam avaliar conformidade legal, ética e técnica \autocite{mkander2023auditing}. Todavia, a literatura também sugere que auditar não é um ato isolado, e sim uma capacidade institucional contínua. Quando essa capacidade não existe, a administração fica exposta ao risco de não conseguir explicar, contestar ou corrigir decisões mediadas por IA em tempo hábil.
Essa observação se conecta a outra dimensão importante: o risco de automação indevida ou de sobredelegação decisória. Em estudos sobre procedimentos administrativos, a automação aparece como promissora em tarefas repetitivas, classificatórias e padronizáveis, mas torna-se problemática quando transborda para contextos em que a apreciação individualizada, a ponderação normativa e a interpretação contextual são essenciais \autocite{parycek2023artificial}. A literatura não afirma que a IA deva ser excluída desses processos, mas insiste que a distribuição entre apoio e decisão precisa ser institucionalmente desenhada. O risco surge quando recomendações algorítmicas passam a operar como decisões de fato, seja porque os agentes públicos confiam excessivamente no sistema, seja porque pressões por produtividade reduzem a revisão humana a uma etapa meramente formal. Esse fenômeno, por vezes descrito como automation bias, tem consequências particularmente severas na administração pública, na medida em que compromete a exigência de motivação, contraditório e revisibilidade dos atos administrativos \autocite{pi2021machine} \autocite{serra2024a}.
No corpus, os riscos institucionais também se manifestam como dependência tecnológica e fragilidade estratégica. A adoção de IA no setor público frequentemente ocorre por meio de contratação externa, parcerias ou aquisição de soluções prontas, o que pode acelerar a implementação, mas também ampliar dependências de fornecedores quanto a manutenção, customização, atualização de modelos, documentação e acesso a parâmetros críticos \autocite{yiitcanlar2022artificial} \autocite{david2023understanding}. Quando a instituição pública não possui massa crítica interna para formular requisitos, validar resultados e negociar cláusulas de transparência e auditoria, a assimetria informacional se aprofunda. Esse risco é institucional porque afeta a autonomia decisória do Estado e sua capacidade de preservar continuidade administrativa. Na APF, marcada por diversidade de capacidades entre órgãos e por processos de contratação nem sempre preparados para lidar com especificidades algorítmicas, o risco de aprisionamento tecnológico não deve ser tratado como problema apenas de TI, mas como questão de governança pública e de soberania organizacional.
Outra contribuição importante da literatura é mostrar que o risco institucional não se distribui de forma homogênea. Criado, Sandoval-Almazán e Gil-Garcia propõem compreender a IA na administração pública a partir de perspectivas micro, meso e macro, o que ajuda a distinguir riscos relacionados ao comportamento individual de agentes, à organização e aos arranjos político-institucionais mais amplos \autocite{criado2024artificial}. No plano micro, destacam-se déficit de competências, uso acrítico de recomendações algorítmicas e incapacidade de interpretar evidências produzidas por sistemas complexos. No plano meso, surgem problemas de coordenação entre áreas jurídicas, tecnológicas, de negócio e controle. No plano macro, entram em cena diretrizes regulatórias, padrões nacionais, accountability democrática e confiança social. Essa leitura multiescalar é particularmente útil para a APF porque impede respostas simplistas. Não basta treinar servidores se o órgão não possui governança clara; tampouco basta criar comitês se faltam padrões de documentação e monitoramento; e nada disso é suficiente se o ambiente regulatório e institucional não oferecer critérios comuns para classificação de risco, prestação de contas e supervisão.
No caso latino-americano, a literatura comparada reforça que os riscos de adoção desigual de IA se combinam com assimetrias de capacidade estatal, heterogeneidade normativa e diferentes estágios de digitalização administrativa \autocite{criado2024inteligencia}. Esse ponto merece destaque porque o debate internacional, muitas vezes orientado por administrações com alta maturidade digital, pode subestimar problemas típicos de contextos em que coexistem excelência técnica localizada e fragilidade institucional difusa. Para a administração pública federal brasileira, isso sugere que mecanismos de mitigação não podem ser copiados de modo acrítico de modelos estrangeiros. Eles precisam ser ajustados a um ambiente em que convivem órgãos com bases de dados relativamente sofisticadas, estruturas robustas de controle e expertise analítica, ao lado de unidades com baixa capacidade de modelagem, pouca integração de dados e limitada familiaridade com auditoria algorítmica. O risco institucional, portanto, é também risco de fragmentação federativa interna à própria União, entendida aqui como conjunto heterogêneo de órgãos e entidades com maturidades distintas.
Diante desse quadro, a literatura converge para a necessidade de mecanismos de mitigação que combinem instrumentos ex ante, monitoramento contínuo e capacidade corretiva. Um primeiro mecanismo é a classificação prévia de casos de uso por nível de risco, vinculando exigências de governança à criticidade da aplicação. Revisões sobre governança de IA mostram que estruturas mais maduras evitam tratar todos os sistemas como equivalentes e distinguem aplicações de baixo impacto operacional de usos que afetam direitos, acesso a benefícios, fiscalização, pessoal ou priorização de políticas \autocite{batool2025ai} \autocite{sigfrids2022how}. Para a APF, isso significa abandonar uma lógica genérica de inovação e adotar um regime escalonado de salvaguardas: quanto maior o potencial de impacto jurídico, distributivo ou reputacional, maior deve ser o grau de documentação, validação, supervisão humana, testagem e controle independente.
Um segundo mecanismo consiste na institucionalização de avaliação de impacto e de documentação ao longo do ciclo de vida do sistema. Embora a literatura apresente diferentes nomenclaturas e formatos, há relativa convergência em torno da ideia de que sistemas de IA precisam ser acompanhados por registros sobre finalidade, base de dados, critérios de modelagem, métricas de desempenho, limitações conhecidas, riscos esperados, público afetado, responsáveis institucionais e procedimentos de revisão \autocite{mkander2023auditing} \autocite{minkkinen2022continuous}. A documentação não é mero formalismo: ela cria memória organizacional, facilita auditoria, permite aprendizado intersetorial e reduz o risco de decisões opacas dependentes de indivíduos específicos. Em ambientes públicos marcados por rotatividade de equipes e mudanças de gestão, esse aspecto é crucial. Sem documentação robusta, a continuidade administrativa fica refém de conhecimento tácito, o que amplia vulnerabilidades e dificulta correções.
A auditoria, por sua vez, aparece na literatura como mecanismo central de mitigação, mas com uma formulação mais ampla do que a simples inspeção técnica do código. Mökander destaca que auditorias de IA podem assumir naturezas legais, éticas e técnicas, e sua utilidade depende de clareza sobre objeto, escopo, independência e critérios de avaliação \autocite{mkander2023auditing}. Minkkinen, ao discutir auditoria contínua, argumenta que sistemas de IA não devem ser avaliados apenas antes da implantação, pois seu comportamento pode se alterar em razão de deriva de dados, mudanças contextuais e novos padrões de uso \autocite{minkkinen2022continuous}. Para a APF, isso tem implicações diretas: mecanismos de controle precisam acompanhar a operação efetiva do sistema, e não apenas sua fase de projeto ou contratação. Um modelo responsivo exigiria métricas de monitoramento em produção, gatilhos para reavaliação, revisão periódica de desempenho e canais para incorporação de reclamações, erros detectados e contestação pelos afetados.
Outro mecanismo recorrente é o fortalecimento da supervisão humana qualificada. A literatura é cautelosa em relação ao uso retórico do “humano no circuito”, mostrando que a simples presença formal de um servidor não garante controle substantivo sobre o sistema \autocite{sigfrids2022how} \autocite{parycek2023artificial}. A supervisão só mitiga riscos quando o agente dispõe de tempo, autoridade, capacitação e informação suficiente para discordar da recomendação algorítmica, registrar justificativas e acionar rotinas de revisão. Isso conecta a governança de IA ao tema das competências. A revisão de Melo aponta que a gestão pública necessita desenvolver competências técnicas, analíticas, éticas e gerenciais para lidar com IA de forma responsável \autocite{melo2026competencias}. Sem esse investimento, a supervisão humana tende a ser ilusória: o servidor torna-se mero homologador de uma saída técnica que não consegue examinar criticamente.
A formação de capacidades institucionais, portanto, não é medida acessória, mas eixo estruturante de mitigação. O corpus mostra que a adoção sustentável de IA requer equipes interdisciplinares, articulação entre áreas finalísticas e de tecnologia, letramento algorítmico para gestores, e integração entre governança digital, compliance, gestão de riscos e controle interno \autocite{goulart2025processos} \autocite{madan2022ai}. Esse ponto é particularmente sensível na APF, onde a existência de competências avançadas em alguns nichos não elimina déficits generalizados de compreensão sobre limites, riscos e requisitos de accountability da IA. A resposta institucional mais adequada parece ser a criação de capacidades distribuídas: núcleos especializados para suporte metodológico e de auditoria, combinados com capacidades mínimas em cada órgão para formulação de demanda, acompanhamento e avaliação de uso.
A literatura também sugere mecanismos de mitigação associados ao desenho organizacional e à coordenação. Em vez de tratar a IA como tema exclusivo de unidades de inovação ou tecnologia, os estudos mais consistentes apontam para arranjos colegiados, políticas internas, fluxos de aprovação e integração com estruturas já existentes de integridade, proteção de dados, segurança da informação e gestão de riscos \autocite{serra2024a} \autocite{goulart2025processos} \autocite{batool2025ai}. Esse aspecto importa porque muitos riscos institucionais decorrem justamente da compartimentalização: a área técnica avalia desempenho, a área jurídica olha conformidade formal, a área finalística busca eficiência operacional e nenhuma instância assume a visão sistêmica do caso de uso. Mecanismos colegiados, quando bem estruturados, podem reduzir essa fragmentação ao exigir deliberação transversal antes da implantação e revisão periódica durante a operação.
Por fim, o corpus indica que a mitigação de riscos institucionais depende de reconhecer a reversibilidade como princípio de prudência administrativa. Em contextos de incerteza elevada, a boa governança não se mede apenas pela capacidade de implantar, mas também pela capacidade de suspender, corrigir, redesenhar ou descontinuar sistemas que produzam efeitos inadequados \autocite{batool2025ai} \autocite{pi2021machine}. Essa orientação é especialmente importante para a APF, onde a pressão por inovação pode induzir irreversibilidades prematuras, seja por contratos, seja por reorganização de fluxos internos. A literatura revisada sugere, ao contrário, uma abordagem incremental e reflexiva: pilotos com critérios claros, avaliação progressiva, escopo delimitado, documentação robusta e possibilidade real de rollback institucional.
Em síntese, os estudos analisados permitem concluir que os principais riscos institucionais da IA na administração pública não são externos à organização, mas produzidos na interseção entre tecnologia, dados, rotinas, competências, estruturas de responsabilidade e ambiente regulatório. Vieses, opacidade, automação indevida, dependência de fornecedores, fragmentação de accountability e insuficiência de capacidades aparecem como riscos centrais porque afetam a legitimidade, a confiabilidade e a continuidade da ação estatal \autocite{criado2024artificial} \autocite{parycek2023artificial} \autocite{mkander2023auditing}. Em resposta, a literatura aponta um repertório de mitigação baseado em classificação de risco, documentação, auditoria contínua, supervisão humana qualificada, capacitação, coordenação intersetorial e reversibilidade. Para a administração pública federal brasileira, a principal implicação é que a adoção de IA só se torna institucionalmente defensável quando acompanhada de mecanismos capazes de transformar inovação tecnológica em capacidade pública controlável, auditável e orientada pelo interesse público, e não apenas em modernização aparente \autocite{minkkinen2022continuous} \autocite{melo2026competencias} \autocite{criado2024inteligencia}.
\subsection{Implicações para o governo público federal brasileiro}
\label{sec:org60cbea2}
Os achados consolidados nas subseções anteriores permitem sustentar que, para a administração pública federal brasileira, a inteligência artificial não deve ser tratada como mero recurso incremental de modernização administrativa, mas como tecnologia organizacional e institucional capaz de reconfigurar processos, competências, controles, responsabilidades e formas de produção da decisão pública. A literatura revisada converge ao indicar que os efeitos da IA no setor público decorrem menos de suas funcionalidades abstratas e mais do modo como ela é incorporada a arranjos concretos de governança, coordenação e responsabilização \autocite{madan2022ai} \autocite{criado2024artificial} \autocite{batool2025ai}. Essa constatação tem implicações diretas para o governo federal brasileiro, cujo desenho administrativo combina forte heterogeneidade entre órgãos, grande diversidade funcional, coexistência de estruturas altamente digitalizadas com outras de baixa maturidade analítica e presença crescente de pressões por eficiência, personalização de serviços, automação de rotinas e intensificação do controle sobre resultados. Nesse contexto, a questão central não é saber se a APF utilizará IA, mas sob quais condições institucionais ela conseguirá fazê-lo sem comprometer legalidade, isonomia, devido processo, proteção de dados, rastreabilidade decisória e confiança pública.
Uma primeira implicação diz respeito à necessidade de substituir uma lógica de adoção oportunística por uma lógica de governança pública da IA. As revisões sobre difusão da tecnologia na administração pública mostram que a adoção tende a ocorrer de forma fragmentada, impulsionada por demandas setoriais, disponibilidade de fornecedores, iniciativas de inovação isoladas ou pressão por ganhos imediatos de produtividade \autocite{madan2022ai} \autocite{alhosani2024opportunities}. Em termos brasileiros, isso significa que diferentes ministérios, autarquias, fundações, agências e empresas estatais podem avançar em ritmos distintos, com critérios também distintos para aquisição, desenvolvimento, documentação e supervisão de sistemas. A consequência provável, se não houver balizas comuns, é a formação de um ecossistema assimétrico em que órgãos mais capacitados internalizam práticas robustas de governança, enquanto outros passam a utilizar soluções opacas, pouco auditáveis ou excessivamente dependentes de fornecedores. A literatura brasileira já aponta que processos de governança em IA em órgãos públicos ainda se encontram em estágio desigual, frequentemente marcados por iniciativas embrionárias, dispersão decisória e institucionalização incompleta \autocite{goulart2025processos} \autocite{serra2024a}. Para a APF, portanto, a implicação não é apenas criar diretrizes genéricas, mas estruturar capacidades centrais de coordenação capazes de induzir padrões mínimos de qualidade institucional sem sufocar adaptações setoriais necessárias.
Isso conduz a uma segunda implicação: o governo federal precisa desenvolver uma arquitetura multinível de governança da IA. A literatura recente sugere que a compreensão da IA em administração pública exige olhar simultaneamente para dimensões micro, meso e macro: práticas dos agentes, arranjos organizacionais e marcos políticos-institucionais mais amplos \autocite{criado2024artificial}. Transposta para a APF, essa abordagem implica reconhecer que uma estratégia responsável de IA não pode se limitar a normativos centrais nem ser deixada integralmente à autonomia das unidades executoras. No nível macro, são necessárias diretrizes transversais sobre princípios, critérios de risco, accountability, registro de sistemas, exigências de documentação, hipóteses de revisão humana e deveres de transparência. No nível meso, cada órgão deve traduzir tais diretrizes em comitês, fluxos de aprovação, papéis funcionais, protocolos de avaliação de impacto e rotinas de monitoramento compatíveis com seus processos substantivos. No nível micro, gestores, analistas, procuradores, equipes de TI e áreas finalísticas precisam de competências para interpretar limites da tecnologia, questionar resultados, identificar degradação de desempenho e interromper usos indevidos. A implicação prática é clara: sem articulação entre esses níveis, a governança da IA permanecerá formal no topo e frágil na operação cotidiana.
Uma terceira implicação refere-se à natureza dos casos de uso prioritários no governo federal. O corpus sugere que a IA produz benefícios mais sustentáveis quando aplicada a tarefas com objetivos delimitados, critérios verificáveis, boa qualidade de dados e possibilidade de supervisão institucional efetiva \autocite{parycek2023artificial} \autocite{pi2021machine}. Em contrapartida, aplicações em ambientes de alta ambiguidade normativa, forte discricionariedade, impacto severo sobre direitos ou baixa explicabilidade tendem a exigir controles muito mais rigorosos. Para a APF, isso significa que a expansão da IA deveria seguir uma lógica de criticidade e não apenas de oportunidade tecnológica. Ferramentas de apoio à triagem documental, classificação de demandas, organização de acervos, extração de informações e detecção preliminar de anomalias podem representar espaços mais promissores para consolidação inicial de capacidades, desde que acompanhadas de rastreabilidade e revisão adequada. Já o uso em domínios como priorização de fiscalização, concessão ou restrição de benefícios, gestão de pessoas, avaliação de desempenho, monitoramento de condutas e apoio a decisões sancionatórias exige grau muito superior de governança, porque o risco de conversão de correlações estatísticas em decisões materialmente injustas é mais alto \autocite{krtner2021predictive} \autocite{cho2023human} \autocite{barros2025inteligencia}. A principal implicação, portanto, é que a APF precisa instituir uma matriz de risco dos usos de IA, diferenciando aplicações assistivas, analíticas, preditivas e decisórias, em vez de tratá-las como um conjunto homogêneo.
Essa diferenciação é especialmente importante diante das assimetrias de maturidade observadas no corpus. Estudos recentes sobre o setor público indicam que a incorporação de IA e, em particular, de IA generativa, ocorre em contextos organizacionais muito desiguais, tanto em termos de infraestrutura quanto de capacidades de governança \autocite{mariano2026inteligencia} \autocite{lyrio2025emerging}. No governo federal brasileiro, a existência de alguns polos avançados de transformação digital pode criar a falsa percepção de que a administração como um todo está preparada para absorver tecnologias mais complexas. Contudo, a literatura sobre implementação mostra que lacunas de infraestrutura, integração de dados, interoperabilidade, cultura organizacional e gestão da mudança continuam sendo determinantes para o sucesso ou fracasso da adoção \autocite{sharma2021implementing} \autocite{david2023understanding}. A implicação prática é que políticas federais de IA não podem pressupor homogeneidade institucional. Será necessário combinar regulação interna, apoio técnico compartilhado, mecanismos de cooperação entre órgãos e programas de elevação gradual de maturidade, sob pena de converter a desigualdade de capacidades em desigualdade de proteção de direitos e de qualidade do serviço público.
Outra implicação decisiva para a APF envolve a centralidade da qualidade dos dados e da infraestrutura informacional. A literatura destaca que o desempenho de sistemas de IA em procedimentos administrativos depende da integridade, representatividade, atualidade e pertinência dos dados utilizados \autocite{parycek2023artificial}. Em ambientes públicos, porém, dados são produzidos em ecossistemas marcados por legados burocráticos, cadastros incompletos, critérios administrativos mutáveis e registros feitos para fins diversos daqueles exigidos por modelos analíticos. Assim, para o governo federal, governança de IA e governança de dados tornam-se dimensões inseparáveis. Não basta estimular projetos de IA; é preciso instituir rotinas de curadoria, padronização semântica, documentação de proveniência, gestão do ciclo de vida dos dados e validação contextual de bases históricas. Sem isso, algoritmos serão alimentados por registros que não apenas contêm ruído e lacunas, mas também cristalizam práticas pretéritas potencialmente desiguais. A implicação substantiva é que a APF deve tratar a preparação e a qualificação dos dados como investimento institucional prévio, e não como etapa acessória da implementação tecnológica \autocite{serra2024a} \autocite{pi2021machine}.
No caso brasileiro, essa preocupação assume relevância adicional porque muitas funções federais lidam com públicos vulneráveis, políticas distributivas complexas e decisões em larga escala. Em tais contextos, erros sistêmicos não permanecem localizados: podem afetar benefícios, acesso a serviços, priorização territorial, tratamento de denúncias, alocação de fiscalização e definição de perfis de risco. A literatura sobre IA ética e responsável mostra que fairness, accountability, transparency e inclusivity não são slogans abstratos, mas dimensões práticas de desenho institucional \autocite{batool2025ai} \autocite{sigfrids2022how}. Para a APF, isso implica incorporar avaliações de impacto que examinem não apenas desempenho técnico, mas também efeitos distributivos, grupos potencialmente prejudicados, possibilidade de contestação, proporcionalidade do uso e compatibilidade com finalidades públicas legítimas. O ponto central é que a racionalidade administrativa baseada em IA não pode ser aceita simplesmente porque promete eficiência; ela precisa demonstrar que sua eficiência não é obtida à custa da opacidade, da discriminação indireta ou da erosão de garantias procedimentais.
Daí decorre uma implicação particularmente sensível: o fortalecimento do dever de explicabilidade institucional. A literatura jurídica e técnica ressalta que a ideia de “explicação” em sistemas algorítmicos é mais complexa do que uma exigência simplificada de abertura da caixa-preta \autocite{brkan2020legal}. Em muitos casos, mesmo quando o funcionamento interno do modelo é tecnicamente descrito, isso não se traduz automaticamente em inteligibilidade útil para administrados, gestores, auditores ou órgãos de controle. Para a APF, a consequência é que a explicabilidade relevante deve ser concebida em camadas: explicação sobre finalidade do sistema, dados utilizados, critérios gerais de operação, grau de autonomia, limites de uso, impactos possíveis e responsáveis pela supervisão. Em procedimentos administrativos, essa exigência se conecta diretamente ao dever de motivação e à possibilidade de revisão. Parycek, Schmid e Novak mostram que a automação em procedimentos administrativos requer atenção específica às condições jurídicas e organizacionais que permitem manter a legitimidade decisória \autocite{parycek2023artificial}. Logo, a implicação para o governo federal é que sistemas de IA não podem substituir a fundamentação administrativa por escores, classificações ou recomendações cujo sentido normativo permaneça inacessível para o processo decisório e para o controle subsequente.
Associada à explicabilidade está a necessidade de institucionalizar auditoria e monitoramento contínuos. O corpus é enfático ao mostrar que a governança da IA não se esgota na fase de desenho ou contratação; ela exige supervisão permanente, capaz de detectar deriva de modelo, mudança de contexto, efeitos adversos emergentes e desconexão entre objetivos declarados e resultados efetivos \autocite{mkander2023auditing} \autocite{minkkinen2022continuous}. Para o governo federal brasileiro, essa implicação é estratégica porque grande parte dos riscos institucionais da IA é dinâmica. Um sistema inicialmente adequado pode tornar-se problemático com alteração do perfil dos dados, mudança normativa, reconfiguração da política pública ou simples ampliação indevida de seu escopo de uso. A APF, portanto, precisa avançar para além de controles ex ante e estruturar mecanismos de auditoria contínua que integrem perspectivas jurídicas, éticas, técnicas e operacionais. Isso supõe registro formal de versões de modelos, métricas periódicas, testes de robustez, análise de falsos positivos e falsos negativos, trilhas de decisão e canais de reporte para incidentes. Também sugere que instâncias de controle interno e externo precisarão ampliar repertório metodológico para auditar sistemas algorítmicos não apenas como ativos tecnológicos, mas como dispositivos com efeitos normativos e organizacionais \autocite{mkander2023auditing}.
Uma implicação adicional incide sobre a relação entre IA, trabalho público e gestão de pessoas. A literatura mostra que a adoção de IA no setor público altera perfis ocupacionais, redistribui tarefas cognitivas, redefine fronteiras entre análise humana e processamento automatizado e pode introduzir novas tensões em seleção, avaliação e alocação de pessoal \autocite{chilunjika2022artificial} \autocite{cho2023human} \autocite{barros2025inteligencia}. Para a APF, isso significa que a política de IA não pode ser formulada apenas como agenda tecnológica; ela deve ser também uma agenda de desenvolvimento de capacidades estatais. Melo destaca, em revisão sistemática, a centralidade de competências em IA para a gestão pública, o que reforça a ideia de que capacidade institucional não se resume à existência de sistemas, mas inclui alfabetização algorítmica, competência analítica, discernimento ético e aptidão para supervisão \autocite{melo2026competencias}. A implicação prática é dupla. De um lado, a APF precisará formar quadros capazes de contratar, acompanhar e avaliar soluções de IA com independência crítica. De outro, terá de preparar servidores de áreas finalísticas para trabalhar com ferramentas que sugerem, priorizam e sintetizam informações sem transferir automaticamente a elas a autoridade decisória. Em vez de imaginar substituição linear do trabalho humano, a literatura recomenda desenhar complementaridades institucionais entre automação e julgamento administrativo \autocite{criado2024artificial} \autocite{madan2022ai}.
No plano federativo e comparado, o corpus também sugere uma lição relevante para o governo federal: a liderança nacional em IA pública não decorre apenas de adotar tecnologia antes dos demais, mas de construir referenciais replicáveis de uso responsável. Estudos sobre governos locais e experiências internacionais mostram que percepções de oportunidade frequentemente convivem com forte incerteza sobre custos, governança e impactos \autocite{yiitcanlar2022artificial} \autocite{david2023understanding}. No espaço latino-americano, a análise comparada associada à Carta Ibero-americana de Inteligência Artificial na Administração Pública indica que o avanço do tema tem sido acompanhado por diferentes graus de institucionalização normativa e capacidade estatal \autocite{criado2024inteligencia}. Para a APF brasileira, isso implica papel duplo: organizar internamente seus próprios usos de IA e, ao mesmo tempo, funcionar como centro difusor de padrões para outras esferas e entidades públicas. Dada a densidade administrativa e regulatória do governo federal, suas escolhas em matéria de documentação, contratação, avaliação de risco e prestação de contas tendem a irradiar efeitos sobre tribunais, órgãos subnacionais, empresas públicas e ecossistemas de fornecedores. A implicação, portanto, é que falhas federais podem multiplicar problemas sistêmicos, enquanto acertos federais podem induzir convergência virtuosa.
Outra consequência importante diz respeito ao desenho de compras públicas e à relação com fornecedores. Embora parte da literatura destaque inovação e flexibilidade como condições favoráveis à adoção, as revisões sobre governança alertam para o risco de dependência excessiva de soluções proprietárias e de assimetrias informacionais entre contratantes públicos e provedores tecnológicos \autocite{batool2025ai} \autocite{madan2022ai}. No caso da APF, isso significa que a contratação de IA não deve ser tratada como simples aquisição de software. É necessário assegurar requisitos de documentação técnica, acesso a informações relevantes para auditoria, definição de responsabilidades por desempenho e incidente, critérios de atualização, governança sobre dados e previsões sobre descontinuidade ou substituição do sistema. Sem tais cautelas, o Estado pode perder capacidade de compreender e controlar tecnologias que orientam suas próprias rotinas administrativas. A implicação é especialmente forte em sistemas de maior criticidade, nos quais a opacidade contratual pode se converter em opacidade institucional.
Por fim, a literatura do corpus leva a uma implicação normativa de fundo para o governo público federal brasileiro: a adoção de IA deve ser subordinada a uma concepção de valor público, e não apenas de eficiência administrativa. Revisões sobre governança ética insistem que confiança, legitimidade, inclusão e responsabilização constituem condições de sustentabilidade do uso de IA no setor público \autocite{sigfrids2022how} \autocite{batool2025ai}. No contexto da APF, isso significa que a pergunta relevante não é apenas se a tecnologia reduz custos, acelera fluxos ou amplia capacidade analítica, mas se o faz de maneira compatível com os fins constitucionais da administração pública. Em uma burocracia orientada pelo interesse público, a IA precisa ser desenhada para ampliar consistência, tempestividade e capacidade estatal sem reduzir espaços de contestação, sem obscurecer a motivação do ato administrativo e sem reproduzir assimetrias sociais sob aparência de neutralidade técnica. A implicação conclusiva, portanto, é que o governo federal brasileiro necessita de uma política de IA fundada em prudência institucional, diferenciação por risco, capacidades permanentes de auditoria, governança de dados, formação de servidores e mecanismos robustos de accountability. Somente assim a IA poderá deixar de ser um vetor de vulnerabilidade administrativa e passar a constituir instrumento legítimo de governança, ética, compliance e gestão de riscos na administração pública federal \autocite{serra2024a} \autocite{goulart2025processos} \autocite{mkander2023auditing} \autocite{minkkinen2022continuous}.
\section{CONSIDERAÇÕES FINAIS}
\label{sec:org6df1468}
\subsection{Síntese final}
\label{sec:orge67c198}
A revisão desenvolvida nesta dissertação permite concluir que o uso de inteligência artificial na administração pública federal brasileira deve ser compreendido menos como simples incorporação de ferramentas digitais e mais como processo de transformação institucional que afeta simultaneamente capacidades estatais, arranjos de governança, critérios de responsabilização e condições de legitimidade da ação administrativa. A literatura examinada converge ao demonstrar que a IA no setor público produz valor quando articulada a finalidades públicas claras, bases de dados adequadas, desenho organizacional consistente e mecanismos de supervisão capazes de prevenir danos, corrigir desvios e justificar decisões perante cidadãos e órgãos de controle \autocite{madan2022ai} \autocite{criado2024artificial} \autocite{serra2024a}. Em sentido oposto, quando adotada sob lógica meramente tecnocrática, orientada por promessas de eficiência desancoradas de salvaguardas institucionais, a IA tende a ampliar opacidade, dependência tecnológica, assimetrias de capacidade entre órgãos e riscos de violação de direitos.
O corpus analisado indica que a difusão da IA na administração pública ocorre em ambientes marcados por heterogeneidade organizacional, múltiplos atores e pressões simultâneas por inovação, austeridade, personalização de serviços e controle de resultados. Por isso, a literatura mais robusta afasta leituras lineares segundo as quais o avanço tecnológico produziria automaticamente melhor desempenho estatal. Ao contrário, a adoção e a difusão da IA dependem de fatores como liderança, prontidão institucional, qualidade da infraestrutura de dados, competências dos servidores, disponibilidade de financiamento, clareza regulatória e adequação dos casos de uso às rotinas decisórias de cada órgão \autocite{madan2022ai} \autocite{sharma2021implementing} \autocite{david2023understanding}. No recorte federal brasileiro, essa constatação é particularmente relevante, pois a administração reúne órgãos com diferentes graus de maturidade digital e distintas exposições a riscos jurídicos, reputacionais e operacionais. Assim, a questão decisiva não é a presença ou ausência da IA, mas a capacidade estatal de governá-la de modo proporcional ao risco e compatível com os princípios que estruturam a administração pública.
Outro resultado sintético importante é que governança, ética, compliance e gestão de riscos não aparecem na literatura como dimensões periféricas ou sucessivas à implantação, mas como condições constitutivas de uma adoção responsável. Revisões recentes sobre governança da IA mostram que princípios como transparência, accountability, justiça, robustez e supervisão humana somente ganham efetividade quando traduzidos em processos, papéis, métricas, documentação e auditoria \autocite{batool2025ai} \autocite{sigfrids2022how}. Isso significa que, no setor público, não basta afirmar compromisso abstrato com “IA ética”. É necessário definir quem autoriza o uso, quais critérios classificam riscos, quando a revisão humana é obrigatória, como são tratados vieses e erros, quais informações devem ser registradas e em que condições um sistema deve ser interrompido. A literatura sobre procedimentos administrativos automatizados reforça essa conclusão ao mostrar que a viabilidade da IA depende de condições de enquadramento normativo, disponibilidade e qualidade dos dados, distinguibilidade entre tarefas estruturadas e decisões complexas, além de margens claras para contestação e revisão \autocite{parycek2023artificial}.
Nesse sentido, a revisão também evidencia que o problema da explicabilidade não pode ser tratado de forma simplista. A demanda por explicação de decisões algorítmicas é juridicamente e tecnicamente complexa, especialmente em sistemas baseados em modelos menos interpretáveis, o que exige distinguir explicação do modelo, justificativa da decisão, documentação do processo e possibilidade efetiva de contestação \autocite{brkan2020legal}. Para a administração pública federal, essa distinção é crucial: em muitos casos, a exigência institucional mais relevante não será tornar integralmente transparente toda a arquitetura matemática do sistema, mas assegurar rastreabilidade suficiente para demonstrar finalidade, critérios operacionais, bases de dados utilizadas, limites de desempenho, papel do agente humano e meios de revisão. A literatura de auditoria de IA reforça que governança responsável depende menos de soluções retóricas e mais de mecanismos contínuos de verificação, combinando enfoques legais, éticos e técnicos \autocite{mkander2023auditing} \autocite{minkkinen2022continuous}.
As aplicações setoriais analisadas ao longo do corpus tornam esse argumento mais concreto. Em áreas como saúde, gestão documental, serviços públicos digitais, fiscalização e recursos humanos, a IA revela potencial para ampliar triagem, acelerar fluxos, melhorar alocação de atenção administrativa e apoiar análises em larga escala \autocite{lemes2020o} \autocite{alhosani2024opportunities} \autocite{lyrio2025emerging}. Contudo, esse potencial varia de acordo com a natureza da tarefa. Em domínios mais estruturados e com critérios verificáveis, a automação tende a ser mais defensável; em contextos de forte discricionariedade, impacto sobre direitos ou dados sensíveis, os riscos de reprodução de desigualdades, erro sistêmico e esvaziamento da responsabilidade pública tornam-se mais elevados \autocite{pi2021machine} \autocite{krtner2021predictive}. Estudos sobre gestão de pessoas e analytics no setor público mostram, por exemplo, que ganhos de eficiência em seleção, avaliação ou planejamento de força de trabalho convivem com riscos de viés, vigilância excessiva e naturalização de inferências estatísticas sobre trajetórias profissionais \autocite{cho2023human} \autocite{chilunjika2022artificial} \autocite{barros2025inteligencia}. Isso confirma que a adequação da IA depende de juízo institucional situado, e não de entusiasmo generalista com automação.
A literatura brasileira e latino-americana reforça, ainda, que a importação acrítica de modelos internacionais de governança é insuficiente. Embora princípios globais sejam relevantes, sua operacionalização no setor público depende de tradições administrativas, capacidades burocráticas, desenho federativo, regime jurídico e trajetórias de transformação digital específicas \autocite{criado2024inteligencia} \autocite{serra2024a} \autocite{goulart2025processos}. No caso brasileiro, a administração pública federal opera em ambiente no qual convivem exigências de legalidade estrita, controle externo intenso, pluralidade de sistemas informacionais e diferentes níveis de capacidade analítica. Por isso, a maturidade em IA tende a ser assimétrica, inclusive com diferenças entre usos mais convencionais de automação analítica e a incorporação recente de IA generativa \autocite{mariano2026inteligencia}. Tal assimetria recomenda cautela contra políticas uniformizantes excessivamente abstratas, mas também contra a fragmentação completa, que deixaria cada órgão definir isoladamente seus padrões de governança.
Em termos substantivos, a síntese final desta dissertação sustenta que a contribuição mais promissora da IA para o governo público federal brasileiro não reside em substituir o juízo administrativo, mas em reorganizar o apoio à decisão, a priorização de fluxos, a extração de padrões em grandes bases de dados e a ampliação da capacidade operacional do Estado sob condições de controle público. Isso recoloca o papel do servidor e da organização: em vez de uma burocracia dispensável, a literatura aponta para a centralidade de burocracias mais qualificadas, capazes de compreender limitações dos modelos, supervisionar fornecedores, interpretar resultados e assumir responsabilidade pelas consequências administrativas e sociais do uso da tecnologia \autocite{criado2024artificial} \autocite{melo2026competencias}. A IA responsável, portanto, exige investimento continuado em capacidades humanas, e não apenas em infraestrutura tecnológica.
Em conclusão, o balanço crítico do corpus permite afirmar que a adoção de IA na administração pública federal brasileira será institucionalmente virtuosa apenas se for tratada como agenda de governança pública e não como agenda restrita de inovação tecnológica. O valor público da IA depende da combinação entre utilidade administrativa, conformidade normativa, prudência ética, auditabilidade e gestão contínua de riscos \autocite{batool2025ai} \autocite{mkander2023auditing} \autocite{parycek2023artificial}. Essa conclusão sintetiza o argumento central da dissertação: no setor público federal, a pergunta correta não é quanto de inteligência artificial incorporar, mas quais usos podem ser legitimamente sustentados por dados adequados, controles verificáveis, competências institucionais e compromisso inequívoco com o interesse público.
\subsection{Limitações do estudo}
\label{sec:orgb0c5a78}
Esta dissertação apresenta limitações que decorrem, em primeiro lugar, de sua própria natureza metodológica. Por se tratar de uma revisão de literatura, o trabalho não realizou investigação empírica direta junto a órgãos e entidades da administração pública federal, nem observação in loco de processos decisórios, fluxos de contratação, rotinas de monitoramento ou práticas concretas de auditoria algorítmica. Em consequência, as conclusões aqui formuladas dependem da qualidade, da profundidade analítica e do grau de comparabilidade do corpus selecionado. Essa característica não invalida os achados, mas impõe cautela quanto ao seu alcance: o estudo permite mapear tendências, categorias analíticas, lacunas e convergências da literatura, porém não autoriza inferências fortes sobre a efetividade real de arranjos específicos de governança de IA no governo federal brasileiro sem validação empírica complementar \autocite{madan2022ai} \autocite{serra2024a}.
Uma segunda limitação reside na própria configuração do campo estudado. A literatura sobre inteligência artificial no setor público ainda é marcada por rápida expansão, heterogeneidade conceitual e forte assimetria entre temas mais visíveis e temas menos documentados. Parte dos estudos usa “IA” em sentido amplo, agregando automação, análise preditiva, aprendizado de máquina, sistemas especialistas e, mais recentemente, IA generativa, sem sempre distinguir com precisão capacidades técnicas, graus de autonomia e perfis de risco. Essa elasticidade conceitual dificulta comparações mais finas entre aplicações e pode produzir generalizações excessivas sobre benefícios, limitações ou exigências regulatórias \autocite{criado2024artificial} \autocite{lyrio2025emerging} \autocite{mariano2026inteligencia}. No caso desta revisão, embora tenha havido esforço de delimitação analítica, o corpus inevitavelmente reflete esse problema de origem. Assim, alguns resultados precisam ser lidos como tendências do debate acadêmico e técnico-institucional, e não como descrição homogênea de uma tecnologia única ou de um padrão estabilizado de adoção estatal.
Também há limitações associadas ao recorte federativo e temporal adotado. O foco na administração pública federal brasileira entre 2019 e 2026 foi adequado ao objetivo de examinar governança, ética, compliance e gestão de riscos em contexto institucional específico, mas esse recorte restringe a generalização para estados, municípios, empresas estatais ou outros países. Mesmo quando o trabalho dialoga com literatura internacional e latino-americana, permanece o problema de transferibilidade institucional: instrumentos considerados viáveis em determinados contextos dependem de capacidades burocráticas, maturidade digital, densidade regulatória e estruturas de controle muito distintas \autocite{criado2024inteligencia} \autocite{yiitcanlar2022artificial} \autocite{david2023understanding}. Além disso, o período analisado coincide com forte aceleração do debate sobre IA generativa, o que significa que parte da literatura mais recente ainda está em formação e pode não captar com estabilidade seus efeitos organizacionais, jurídicos e éticos de médio prazo \autocite{mariano2026inteligencia}.
Outra limitação importante diz respeito à distribuição temática do corpus. Embora a dissertação tenha buscado integrar dimensões de governança, ética, compliance, auditoria, capacidades estatais e aplicações setoriais, a literatura disponível não é equilibrada entre esses eixos. Há maior densidade de trabalhos de caráter programático ou de revisão sobre princípios, oportunidades e desafios gerais da IA, enquanto há menor disponibilidade de estudos que descrevam, com o mesmo rigor, mecanismos concretos de implementação, indicadores de desempenho, protocolos de classificação de risco, padrões de documentação, rotinas de revisão humana e resultados de auditorias em organizações públicas \autocite{batool2025ai} \autocite{sigfrids2022how} \autocite{alhosani2024opportunities}. Em outras palavras, o campo ainda oferece mais prescrições do que evidências comparáveis sobre o funcionamento real das salvaguardas. Isso afeta diretamente esta dissertação, pois parte da análise de governança responsável precisa ser construída a partir de recomendações normativas e conceituais, e não apenas de evidência consolidada sobre arranjos testados e avaliados.
Há, ademais, um limite específico relacionado à opacidade empírica dos próprios sistemas estudados. Em matéria de IA no setor público, a literatura registra obstáculos recorrentes de acesso a informações sobre bases de dados, parâmetros de modelos, contratos com fornecedores, métricas de desempenho, critérios de intervenção humana e registros de incidentes. Tais barreiras reduzem a auditabilidade externa e dificultam a produção de conhecimento independente sobre riscos e impactos \autocite{mkander2023auditing} \autocite{minkkinen2022continuous}. Esse problema é agravado quando se trata de sistemas com baixa explicabilidade técnica ou quando a organização pública depende de soluções proprietárias. Como argumenta a literatura sobre explicação e decisões algorítmicas, nem sempre é juridicamente ou tecnicamente simples converter o funcionamento interno de modelos complexos em justificativas compreensíveis e institucionalmente úteis \autocite{brkan2020legal}. Consequentemente, esta revisão consegue discutir de forma robusta o problema da transparência e da rastreabilidade, mas encontra limites objetivos para avaliar comparativamente, em profundidade empírica, o grau efetivo de explicabilidade praticado em diferentes experiências governamentais.
O estudo também enfrenta a limitação de trabalhar com uma base bibliográfica majoritariamente composta por revisões, artigos conceituais e análises de tendências, com menor presença de avaliações longitudinais de implementação. Isso significa que a literatura frequentemente enfatiza potenciais e riscos em termos prospectivos, mas oferece menos evidência sobre efeitos institucionais acumulados, aprendizagem organizacional ao longo do tempo e sustentabilidade de mecanismos de governança após a fase inicial de adoção \autocite{parycek2023artificial} \autocite{pi2021machine}. Em especial, permanecem insuficientemente documentados os efeitos de segunda ordem da IA sobre discricionariedade administrativa, redistribuição interna de poder técnico, dependência de fornecedores, reconfiguração de competências e mudança nas relações entre áreas finalísticas, jurídicas, de tecnologia e controle. Embora a dissertação tenha explorado essas dimensões de modo analítico, elas ainda aparecem mais como fronteiras interpretativas do que como objetos já plenamente consolidados na literatura.
Outra limitação relevante refere-se ao tratamento das capacidades humanas e organizacionais. O corpus confirma que competências técnicas, letramento algorítmico, governança interdisciplinar e preparo da burocracia são variáveis centrais para o uso responsável da IA \autocite{melo2026competencias} \autocite{chilunjika2022artificial} \autocite{cho2023human} \autocite{barros2025inteligencia}. Contudo, o próprio material disponível ainda é desigual na articulação entre competências individuais, desenho organizacional e estruturas formais de accountability. Em muitos casos, a literatura identifica a necessidade de qualificação, mas não detalha suficientemente como capacidades são construídas, medidas, institucionalizadas e preservadas em ambientes públicos marcados por rotatividade, restrições orçamentárias e assimetrias salariais frente ao setor privado. Assim, esta dissertação consegue demonstrar a centralidade do fator humano para governança e compliance, mas não esgota as formas pelas quais tais capacidades podem ser operacionalizadas de maneira sustentável na administração pública federal.
Por fim, convém reconhecer uma limitação interpretativa mais ampla: ao organizar um corpus multidisciplinar e internacional para responder a uma questão situada no contexto federal brasileiro, o estudo inevitavelmente opera por aproximações analíticas. Essa estratégia é útil para construir quadro conceitual abrangente e identificar padrões recorrentes, mas pode atenuar especificidades jurídico-administrativas, setoriais e organizacionais de casos concretos. Em síntese, a principal força desta dissertação — integrar literatura dispersa para oferecer visão sistemática sobre IA, governança, ética, compliance e riscos — também define seu principal limite: ela esclarece o estado do debate e os contornos normativos e institucionais do problema, mas não substitui investigação empírica aprofundada sobre sistemas, órgãos, contratos, fluxos de decisão e mecanismos efetivos de supervisão no governo federal brasileiro \autocite{goulart2025processos} \autocite{serra2024a} \autocite{madan2022ai}.
\subsection{Agenda futura de pesquisa}
\label{sec:org0227407}
A agenda futura de pesquisa sobre inteligência artificial no governo público federal brasileiro precisa avançar para além do diagnóstico geral de oportunidades, riscos e princípios normativos, concentrando-se na produção de evidência comparável sobre como arranjos institucionais concretos condicionam resultados, falhas e capacidades de controle. A literatura revisada mostra que a adoção de IA no setor público tem sido discutida, em larga medida, por meio de revisões sistemáticas, modelos conceituais e proposições de governança, o que foi decisivo para consolidar o campo, mas ainda insuficiente para responder com precisão quais combinações de liderança, capacidades técnicas, desenho regulatório, infraestrutura de dados e mecanismos de responsabilização produzem implementação responsável em contextos administrativos específicos \autocite{madan2022ai} \autocite{serra2024a} \autocite{criado2024artificial}. Desse modo, um primeiro eixo promissor consiste em pesquisas empíricas de maior profundidade, com estudos de caso comparados entre órgãos federais, capazes de examinar diferenças de maturidade institucional entre aplicações de apoio à decisão, automação de procedimentos, fiscalização, triagem documental e atendimento digital.
Esse deslocamento analítico é importante porque a literatura já sugere que a IA pública não opera como tecnologia neutra ou homogênea, mas como conjunto de sistemas dependentes de contexto organizacional, qualidade de dados, enquadramento jurídico e aceitação burocrática \autocite{parycek2023artificial} \autocite{pi2021machine}. Assim, futuras pesquisas devem mapear não apenas “se” a administração pública utiliza IA, mas “como”, “com quais salvaguardas”, “em que etapas do processo decisório” e “sob quais critérios de revisão humana”. Há espaço, por exemplo, para investigações que distingam aplicações meramente instrumentais, voltadas à eficiência administrativa, de aplicações com maior potencial de afetar direitos, priorização distributiva, fiscalização e tratamento diferenciado de cidadãos. Essa diferenciação permitiria construir tipologias de risco mais aderentes ao setor público federal brasileiro, evitando o problema recorrente de tratar toda IA estatal como se apresentasse o mesmo perfil de impacto \autocite{batool2025ai} \autocite{sigfrids2022how}.
Um segundo eixo prioritário diz respeito à governança multinível e à ecologia institucional da IA pública. A literatura recente destaca a necessidade de articular perspectivas micro, meso e macro para compreender atores, incentivos, capacidades e estruturas normativas envolvidas na adoção de IA \autocite{criado2024artificial}. No caso brasileiro, isso sugere pesquisas voltadas à interação entre órgãos setoriais, instâncias centrais de governo digital, unidades jurídicas, controles internos, auditoria, corregedoria, proteção de dados e fornecedores privados. Em vez de presumir uma governança unitária, futuros estudos devem investigar como responsabilidades são distribuídas, sobrepostas ou difusas ao longo do ciclo de vida dos sistemas. Tal agenda é particularmente relevante porque revisões sobre governança de IA têm mostrado que princípios abstratos de transparência, justiça e accountability só se tornam efetivos quando convertidos em processos, papéis organizacionais, documentação e mecanismos verificáveis de supervisão \autocite{batool2025ai} \autocite{goulart2025processos}.
Nesse ponto, uma lacuna central refere-se à auditoria algorítmica. A literatura tem avançado no reconhecimento da auditoria como mecanismo promissor de governança, mas ainda carece de evidência sistemática sobre sua viabilidade operacional, seus custos, seus limites metodológicos e sua capacidade real de prevenir danos no setor público \autocite{mkander2023auditing} \autocite{minkkinen2022continuous}. No contexto federal brasileiro, pesquisas futuras podem examinar como estruturar trilhas de auditoria para sistemas baseados em aprendizado de máquina e, mais recentemente, em IA generativa, incluindo registro de versões de modelos, documentação de bases de dados, parâmetros de performance, tratamento de incidentes, monitoramento de deriva e revisão periódica dos critérios de uso. Também seria valioso investigar quais modelos de auditoria são mais adequados em diferentes situações: auditorias internas contínuas, avaliações independentes ex ante, revisões ex post orientadas por impacto ou combinações dessas modalidades. Esse esforço permitiria deslocar o debate de uma confiança retórica em “IA responsável” para práticas verificáveis de controle institucional.
Um quarto eixo envolve a tensão entre explicabilidade, legalidade e efetividade administrativa. A literatura indica que a exigência de explicação de decisões algorítmicas não pode ser tratada de forma simplista, sobretudo em sistemas complexos, nos quais inteligibilidade técnica, compreensão organizacional e justificativa juridicamente adequada nem sempre coincidem \autocite{brkan2020legal}. Em vez de buscar respostas universalistas, pesquisas futuras devem examinar empiricamente quais níveis de explicação são necessários e possíveis em diferentes tipos de uso estatal da IA. Em aplicações de triagem, apoio analítico ou priorização de casos, por exemplo, pode ser mais relevante a reconstrução procedimental da decisão administrativa do que a abertura irrestrita do modelo em termos puramente técnicos. Já em contextos com maior potencial de restrição de direitos, a necessidade de explicabilidade substantiva e contestabilidade tende a ser mais exigente \autocite{parycek2023artificial}. Há, portanto, espaço para estudos jurídico-organizacionais que aproximem teoria regulatória, administração pública e ciência de dados.
A agenda futura também deve incorporar com mais centralidade a dimensão do trabalho público e das capacidades burocráticas. Parte importante da literatura mostra que a adoção de IA depende menos de aquisição tecnológica isolada e mais da formação de competências, da reorganização de rotinas e da construção de discernimento institucional para uso, supervisão e contestação de sistemas automatizados \autocite{melo2026competencias} \autocite{chilunjika2022artificial}. No governo federal, isso abre uma frente de pesquisa sobre competências híbridas: não apenas habilidades técnicas em dados e modelos, mas capacidades de tradução entre linguagem jurídica, lógica administrativa, gestão de riscos, proteção de dados, auditoria e desenho de políticas. Estudos futuros podem investigar quais perfis profissionais são necessários, como se distribuem entre carreiras e como programas de capacitação afetam a qualidade da governança da IA. Essa agenda deve incluir ainda o campo de gestão de pessoas, no qual o uso de analytics e automação pode tanto qualificar decisões quanto reproduzir vieses e opacidades em recrutamento, avaliação e alocação de servidores \autocite{cho2023human} \autocite{barros2025inteligencia}.
Outro eixo incontornável refere-se às assimetrias de maturidade entre setores, órgãos e tipos de tecnologia. O avanço recente da IA generativa tornou ainda mais visível que a administração pública convive simultaneamente com aplicações tradicionais de automação, modelos preditivos relativamente delimitados e ferramentas generativas de uso difuso, muitas vezes incorporadas de forma experimental e descentralizada \autocite{mariano2026inteligencia} \autocite{lyrio2025emerging}. Pesquisas futuras devem distinguir essas camadas tecnológicas e comparar seus efeitos sobre governança, risco e compliance. Em particular, será importante examinar se a IA generativa introduz novos padrões de informalidade organizacional, uso não autorizado, vazamento de dados, dependência de plataformas privadas e fragilização de registros decisórios. Também merecem atenção os efeitos de acoplamento entre ferramentas generativas e sistemas administrativos existentes, sobretudo quando saídas produzidas por modelos de linguagem passam a influenciar análise documental, elaboração de pareceres, comunicação com usuários ou triagem de demandas.
Por fim, a agenda de pesquisa precisa fortalecer perspectivas comparativas, sem perder aderência ao caso brasileiro. Estudos internacionais e latino-americanos sugerem que fatores como capacidade estatal, orientação estratégica, pressão por inovação, percepção gerencial e desenho federativo afetam intensamente a adoção de tecnologias digitais e de IA \autocite{yiitcanlar2022artificial} \autocite{david2023understanding} \autocite{criado2024inteligencia}. No entanto, ainda há pouca acumulação comparativa capaz de mostrar o que é transferível e o que depende de contextos normativos e organizacionais específicos. Para o governo federal brasileiro, isso aponta para pesquisas comparadas entre setores como saúde, trabalho, controle, previdência, compras públicas e serviços digitais, bem como entre experiências nacionais e referenciais ibero-americanos. A contribuição mais relevante dessa agenda não será apenas ampliar o inventário de casos, mas identificar mecanismos causais: que tipo de arranjo institucional reduz riscos, melhora auditabilidade, preserva discricionariedade legítima e reforça o interesse público. Em síntese, o próximo ciclo de estudos precisa deslocar o campo da descrição geral para a avaliação crítica de capacidades estatais, instrumentos de controle e efeitos concretos da IA sobre a ação administrativa \autocite{madan2022ai} \autocite{serra2024a} \autocite{criado2024artificial}.
:UNNUMBERED: notoc
\section*{GLOSSÁRIO}
\label{sec:orgea7f853}
\begin{itemize}
\item Accountability algorítmica — capacidade de atribuir responsabilidades e justificar decisões mediadas por sistemas algorítmicos.
\item Auditabilidade — possibilidade de examinar, verificar e rastrear o funcionamento, os dados, os registros e os resultados de um sistema de IA.
\item Compliance público — conjunto de mecanismos de conformidade normativa, integridade institucional e aderência a controles aplicáveis à atuação administrativa.
\item Explicabilidade — grau em que o funcionamento ou o resultado de um sistema de IA pode ser compreendido e comunicado de modo adequado ao contexto decisório.
\item Governança da IA — arranjo de princípios, estruturas, processos e controles voltados à direção, supervisão e uso responsável de sistemas de inteligência artificial.
\item Supervisão humana significativa — intervenção humana qualificada e efetiva sobre o desenho, a operação e a revisão de sistemas automatizados com impacto administrativo.
\item Viés algorítmico — distorção sistemática em dados, modelos ou resultados que pode produzir tratamento inadequado, desigual ou discriminatório.
\item Avaliação de impacto algorítmico — procedimento de identificação, análise e mitigação prévia e continuada de efeitos jurídicos, éticos, sociais e institucionais associados ao uso de IA.
\end{itemize}
\section*{ÍNDICE}
\label{sec:org575e249}
\printbibliography
\end{document}
